% concatenated from bare_jrnl.tex
%% bare_jrnl.tex
%% V1.4a
%% 2014/09/17
%% by Michael Shell
%% see http://www.michaelshell.org/
%% for current contact information.
%%
%% This is a skeleton file demonstrating the use of IEEEtran.cls
%% (requires IEEEtran.cls version 1.8a or later) with an IEEE
%% journal paper.
%%
%% Support sites:
%% http://www.michaelshell.org/tex/ieeetran/
%% http://www.ctan.org/tex-archive/macros/latex/contrib/IEEEtran/
%% and
%% http://www.ieee.org/

%%*************************************************************************
%% Legal Notice:
%% This code is offered as-is without any warranty either expressed or
%% implied; without even the implied warranty of MERCHANTABILITY or
%% FITNESS FOR A PARTICULAR PURPOSE! 
%% User assumes all risk.
%% In no event shall IEEE or any contributor to this code be liable for
%% any damages or losses, including, but not limited to, incidental,
%% consequential, or any other damages, resulting from the use or misuse
%% of any information contained here.
%%
%% All comments are the opinions of their respective authors and are not
%% necessarily endorsed by the IEEE.
%%
%% This work is distributed under the LaTeX Project Public License (LPPL)
%% ( http://www.latex-project.org/ ) version 1.3, and may be freely used,
%% distributed and modified. A copy of the LPPL, version 1.3, is included
%% in the base LaTeX documentation of all distributions of LaTeX released
%% 2003/12/01 or later.
%% Retain all contribution notices and credits.
%% ** Modified files should be clearly indicated as such, including  **
%% ** renaming them and changing author support contact information. **
%%
%% File list of work: IEEEtran.cls, IEEEtran_HOWTO.pdf, bare_adv.tex,
%%                    bare_conf.tex, bare_jrnl.tex, bare_conf_compsoc.tex,
%%                    bare_jrnl_compsoc.tex, bare_jrnl_transmag.tex
%%*************************************************************************


% *** Authors should verify (and, if needed, correct) their LaTeX system  ***
% *** with the testflow diagnostic prior to trusting their LaTeX platform ***
% *** with production work. IEEE's font choices and paper sizes can       ***
% *** trigger bugs that do not appear when using other class files.       ***                          ***
% The testflow support page is at:
% http://www.michaelshell.org/tex/testflow/



\documentclass[conference]{IEEEtran}
%
% If IEEEtran.cls has not been installed into the LaTeX system files,
% manually specify the path to it like:
% \documentclass[journal]{../sty/IEEEtran}





% Some very useful LaTeX packages include:
% (uncomment the ones you want to load)


% *** MISC UTILITY PACKAGES ***
%
%\usepackage{ifpdf}
\usepackage{amsmath,amssymb,amsfonts}
\usepackage{algorithmic}
\usepackage{graphicx}
\usepackage{textcomp}
\usepackage{xcolor}
\usepackage{times}
\usepackage{enumerate}
\usepackage{amsmath,amssymb,amsfonts}
\usepackage{mathtools} % for \mathclap
\usepackage{amsthm} % for theorem/lemma environments
\newcommand{\bball}{\mathbb{B}^{\mathclap{d}}(0,1)}

\usepackage{cleveref}
\usepackage{mathrsfs}
\usepackage{comment}
\usepackage{cite}
\newcommand\myworries[1]{\textcolor{red}{#1}}
\newcommand\numberthis{\addtocounter{equation}{1}\tag{\theequation}}
\newenvironment{hproof}{%
  \renewcommand{\proofname}{Sketch of proof}\proof}{\endproof}
\def\BibTeX{{\rm B\kern-.05em{\sc i\kern-.025em b}\kern-.08em
    T\kern-.1667em\lower.7ex\hbox{E}\kern-.125emX}}
% Heiko Oberdiek's ifpdf.sty is very useful if you need conditional
% compilation based on whether the output is pdf or dvi.
% usage:
% \ifpdf
%   % pdf code
% \else
%   % dvi code
% \fi
% The latest version of ifpdf.sty can be obtained from:
% http://www.ctan.org/tex-archive/macros/latex/contrib/oberdiek/
% Also, note that IEEEtran.cls V1.7 and later provides a builtin
% \ifCLASSINFOpdf conditional that works the same way.
% When switching from latex to pdflatex and vice-versa, the compiler may
% have to be run twice to clear warning/error messages.


\usepackage{placeins}
\usepackage{amsthm}
\newtheorem{definition}{Definition}
\newtheorem{proposition}{Proposition}
\newtheorem{theorem}{Theorem}
\newtheorem{corollary}{Corollary}
\newtheorem{remark}{Remark}
\newtheorem{lemma}{Lemma}
\addtolength{\topmargin}{0.05in}
\usepackage{geometry}

% default for the rest of the document
\geometry{
  top=72pt,
  left=54pt,
  right=54pt,
  bottom=54pt
}

% *** CITATION PACKAGES ***
%
%\usepackage{cite}
% cite.sty was written by Donald Arseneau
% V1.6 and later of IEEEtran pre-defines the format of the cite.sty package
% \cite{} output to follow that of IEEE. Loading the cite package will
% result in citation numbers being automatically sorted and properly
% "compressed/ranged". e.g., [1], [9], [2], [7], [5], [6] without using
% cite.sty will become [1], [2], [5]--[7], [9] using cite.sty. cite.sty's
% \cite will automatically add leading space, if needed. Use cite.sty's
% noadjust option (cite.sty V3.8 and later) if you want to turn this off
% such as if a citation ever needs to be enclosed in parenthesis.
% cite.sty is already installed on most LaTeX systems. Be sure and use
% version 5.0 (2009-03-20) and later if using hyperref.sty.
% The latest version can be obtained at:
% http://www.ctan.org/tex-archive/macros/latex/contrib/cite/
% The documentation is contained in the cite.sty file itself.






% *** GRAPHICS RELATED PACKAGES ***
%
\ifCLASSINFOpdf
  % \usepackage[pdftex]{graphicx}
  % declare the path(s) where your graphic files are
  % \graphicspath{{../pdf/}{../jpeg/}}
  % and their extensions so you won't have to specify these with
  % every instance of \includegraphics
  % \DeclareGraphicsExtensions{.pdf,.jpeg,.png}
\else
  % or other class option (dvipsone, dvipdf, if not using dvips). graphicx
  % will default to the driver specified in the system graphics.cfg if no
  % driver is specified.
  % \usepackage[dvips]{graphicx}
  % declare the path(s) where your graphic files are
  % \graphicspath{{../eps/}}
  % and their extensions so you won't have to specify these with
  % every instance of \includegraphics
  % \DeclareGraphicsExtensions{.eps}
\fi
% graphicx was written by David Carlisle and Sebastian Rahtz. It is
% required if you want graphics, photos, etc. graphicx.sty is already
% installed on most LaTeX systems. The latest version and documentation
% can be obtained at: 
% http://www.ctan.org/tex-archive/macros/latex/required/graphics/
% Another good source of documentation is "Using Imported Graphics in
% LaTeX2e" by Keith Reckdahl which can be found at:
% http://www.ctan.org/tex-archive/info/epslatex/
%
% latex, and pdflatex in dvi mode, support graphics in encapsulated
% postscript (.eps) format. pdflatex in pdf mode supports graphics
% in .pdf, .jpeg, .png and .mps (metapost) formats. Users should ensure
% that all non-photo figures use a vector format (.eps, .pdf, .mps) and
% not a bitmapped formats (.jpeg, .png). IEEE frowns on bitmapped formats
% which can result in "jaggedy"/blurry rendering of lines and letters as
% well as large increases in file sizes.
%
% You can find documentation about the pdfTeX application at:
% http://www.tug.org/applications/pdftex





% *** MATH PACKAGES ***
%
%\usepackage[cmex10]{amsmath}
% A popular package from the American Mathematical Society that provides
% many useful and powerful commands for dealing with mathematics. If using
% it, be sure to load this package with the cmex10 option to ensure that
% only type 1 fonts will utilized at all point sizes. Without this option,
% it is possible that some math symbols, particularly those within
% footnotes, will be rendered in bitmap form which will result in a
% document that can not be IEEE Xplore compliant!
%
% Also, note that the amsmath package sets \interdisplaylinepenalty to 10000
% thus preventing page breaks from occurring within multiline equations. Use:
%\interdisplaylinepenalty=2500
% after loading amsmath to restore such page breaks as IEEEtran.cls normally
% does. amsmath.sty is already installed on most LaTeX systems. The latest
% version and documentation can be obtained at:
% http://www.ctan.org/tex-archive/macros/latex/required/amslatex/math/





% *** SPECIALIZED LIST PACKAGES ***
%
%\usepackage{algorithmic}
% algorithmic.sty was written by Peter Williams and Rogerio Brito.
% This package provides an algorithmic environment fo describing algorithms.
% You can use the algorithmic environment in-text or within a figure
% environment to provide for a floating algorithm. Do NOT use the algorithm
% floating environment provided by algorithm.sty (by the same authors) or
% algorithm2e.sty (by Christophe Fiorio) as IEEE does not use dedicated
% algorithm float types and packages that provide these will not provide
% correct IEEE style captions. The latest version and documentation of
% algorithmic.sty can be obtained at:
% http://www.ctan.org/tex-archive/macros/latex/contrib/algorithms/
% There is also a support site at:
% http://algorithms.berlios.de/index.html
% Also of interest may be the (relatively newer and more customizable)
% algorithmicx.sty package by Szasz Janos:
% http://www.ctan.org/tex-archive/macros/latex/contrib/algorithmicx/




% *** ALIGNMENT PACKAGES ***
%
%\usepackage{array}
% Frank Mittelbach's and David Carlisle's array.sty patches and improves
% the standard LaTeX2e array and tabular environments to provide better
% appearance and additional user controls. As the default LaTeX2e table
% generation code is lacking to the point of almost being broken with
% respect to the quality of the end results, all users are strongly
% advised to use an enhanced (at the very least that provided by array.sty)
% set of table tools. array.sty is already installed on most systems. The
% latest version and documentation can be obtained at:
% http://www.ctan.org/tex-archive/macros/latex/required/tools/


% IEEEtran contains the IEEEeqnarray family of commands that can be used to
% generate multiline equations as well as matrices, tables, etc., of high
% quality.




% *** SUBFIGURE PACKAGES ***
%\ifCLASSOPTIONcompsoc
%  \usepackage[caption=false,font=normalsize,labelfont=sf,textfont=sf]{subfig}
%\else
%  \usepackage[caption=false,font=footnotesize]{subfig}
%\fi
% subfig.sty, written by Steven Douglas Cochran, is the modern replacement
% for subfigure.sty, the latter of which is no longer maintained and is
% incompatible with some LaTeX packages including fixltx2e. However,
% subfig.sty requires and automatically loads Axel Sommerfeldt's caption.sty
% which will override IEEEtran.cls' handling of captions and this will result
% in non-IEEE style figure/table captions. To prevent this problem, be sure
% and invoke subfig.sty's "caption=false" package option (available since
% subfig.sty version 1.3, 2005/06/28) as this is will preserve IEEEtran.cls
% handling of captions.
% Note that the Computer Society format requires a larger sans serif font
% than the serif footnote size font used in traditional IEEE formatting
% and thus the need to invoke different subfig.sty package options depending
% on whether compsoc mode has been enabled.
%
% The latest version and documentation of subfig.sty can be obtained at:
% http://www.ctan.org/tex-archive/macros/latex/contrib/subfig/




% *** FLOAT PACKAGES ***
%
%\usepackage{fixltx2e}
% fixltx2e, the successor to the earlier fix2col.sty, was written by
% Frank Mittelbach and David Carlisle. This package corrects a few problems
% in the LaTeX2e kernel, the most notable of which is that in current
% LaTeX2e releases, the ordering of single and double column floats is not
% guaranteed to be preserved. Thus, an unpatched LaTeX2e can allow a
% single column figure to be placed prior to an earlier double column
% figure. The latest version and documentation can be found at:
% http://www.ctan.org/tex-archive/macros/latex/base/


%\usepackage{stfloats}
% stfloats.sty was written by Sigitas Tolusis. This package gives LaTeX2e
% the ability to do double column floats at the bottom of the page as well
% as the top. (e.g., "\begin{figure*}[!b]" is not normally possible in
% LaTeX2e). It also provides a command:
%\fnbelowfloat
% to enable the placement of footnotes below bottom floats (the standard
% LaTeX2e kernel puts them above bottom floats). This is an invasive package
% which rewrites many portions of the LaTeX2e float routines. It may not work
% with other packages that modify the LaTeX2e float routines. The latest
% version and documentation can be obtained at:
% http://www.ctan.org/tex-archive/macros/latex/contrib/sttools/
% Do not use the stfloats baselinefloat ability as IEEE does not allow
% \baselineskip to stretch. Authors submitting work to the IEEE should note
% that IEEE rarely uses double column equations and that authors should try
% to avoid such use. Do not be tempted to use the cuted.sty or midfloat.sty
% packages (also by Sigitas Tolusis) as IEEE does not format its papers in
% such ways.
% Do not attempt to use stfloats with fixltx2e as they are incompatible.
% Instead, use Morten Hogholm'a dblfloatfix which combines the features
% of both fixltx2e and stfloats:
%
% \usepackage{dblfloatfix}
% The latest version can be found at:
% http://www.ctan.org/tex-archive/macros/latex/contrib/dblfloatfix/




%\ifCLASSOPTIONcaptionsoff
%  \usepackage[nomarkers]{endfloat}
% \let\MYoriglatexcaption\caption
% \renewcommand{\caption}[2][\relax]{\MYoriglatexcaption[#2]{#2}}
%\fi
% endfloat.sty was written by James Darrell McCauley, Jeff Goldberg and 
% Axel Sommerfeldt. This package may be useful when used in conjunction with 
% IEEEtran.cls'  captionsoff option. Some IEEE journals/societies require that
% submissions have lists of figures/tables at the end of the paper and that
% figures/tables without any captions are placed on a page by themselves at
% the end of the document. If needed, the draftcls IEEEtran class option or
% \CLASSINPUTbaselinestretch interface can be used to increase the line
% spacing as well. Be sure and use the nomarkers option of endfloat to
% prevent endfloat from "marking" where the figures would have been placed
% in the text. The two hack lines of code above are a slight modification of
% that suggested by in the endfloat docs (section 8.4.1) to ensure that
% the full captions always appear in the list of figures/tables - even if
% the user used the short optional argument of \caption[]{}.
% IEEE papers do not typically make use of \caption[]'s optional argument,
% so this should not be an issue. A similar trick can be used to disable
% captions of packages such as subfig.sty that lack options to turn off
% the subcaptions:
% For subfig.sty:
% \let\MYorigsubfloat\subfloat
% \renewcommand{\subfloat}[2][\relax]{\MYorigsubfloat[]{#2}}
% However, the above trick will not work if both optional arguments of
% the \subfloat command are used. Furthermore, there needs to be a
% description of each subfigure *somewhere* and endfloat does not add
% subfigure captions to its list of figures. Thus, the best approach is to
% avoid the use of subfigure captions (many IEEE journals avoid them anyway)
% and instead reference/explain all the subfigures within the main caption.
% The latest version of endfloat.sty and its documentation can obtained at:
% http://www.ctan.org/tex-archive/macros/latex/contrib/endfloat/
%
% The IEEEtran \ifCLASSOPTIONcaptionsoff conditional can also be used
% later in the document, say, to conditionally put the References on a 
% page by themselves.




% *** PDF, URL AND HYPERLINK PACKAGES ***
%
%\usepackage{url}
% url.sty was written by Donald Arseneau. It provides better support for
% handling and breaking URLs. url.sty is already installed on most LaTeX
% systems. The latest version and documentation can be obtained at:
% http://www.ctan.org/tex-archive/macros/latex/contrib/url/
% Basically, \url{my_url_here}.




% *** Do not adjust lengths that control margins, column widths, etc. ***
% *** Do not use packages that alter fonts (such as pslatex).         ***
% There should be no need to do such things with IEEEtran.cls V1.6 and later.
% (Unless specifically asked to do so by the journal or conference you plan
% to submit to, of course. )


% correct bad hyphenation here
\hyphenation{op-tical net-works semi-conduc-tor}

\IEEEoverridecommandlockouts

\begin{document}
%
% paper title
% Titles are generally capitalized except for words such as a, an, and, as,
% at, but, by, for, in, nor, of, on, or, the, to and up, which are usually
% not capitalized unless they are the first or last word of the title.
% Linebreaks \\ can be used within to get better formatting as desired.
% Do not put math or special symbols in the title.
\title{Instantiating Bayesian CVaR lower bounds in Interactive Decision Making Problems}
%
%
% author names and IEEE memberships
% note positions of commas and nonbreaking spaces ( ~ ) LaTeX will not break
% a structure at a ~ so this keeps an author's name from being broken across
% two lines.
% use \thanks{} to gain access to the first footnote area
% a separate \thanks must be used for each paragraph as LaTeX2e's \thanks
% was not built to handle multiple paragraphs
%

\author{
\IEEEauthorblockN{Raghav Bongole, Tobias J. Oechtering, and Mikael Skoglund}\\
\IEEEauthorblockA{Department of Information Science and Engineering (ISE)\\
KTH Royal Institute of Technology\\
%\texttt{\{bongole, oech, skoglund\}@kth.se}
}% <-this % stops a space
\thanks{This work is supported by the Knut and Wallenberg Foundation.}
}

% note the % following the last \IEEEmembership and also \thanks - 
% these prevent an unwanted space from occurring between the last author name
% and the end of the author line. i.e., if you had this:
% 
% \author{....lastname \thanks{...} \thanks{...} }
%                     ^------------^------------^----Do not want these spaces!
%
% a space would be appended to the last name and could cause every name on that
% line to be shifted left slightly. This is one of those "LaTeX things". For
% instance, "\textbf{A} \textbf{B}" will typeset as "A B" not "AB". To get
% "AB" then you have to do: "\textbf{A}\textbf{B}"
% \thanks is no different in this regard, so shield the last } of each \thanks
% that ends a line with a % and do not let a space in before the next \thanks.
% Spaces after \IEEEmembership other than the last one are OK (and needed) as
% you are supposed to have spaces between the names. For what it is worth,
% this is a minor point as most people would not even notice if the said evil
% space somehow managed to creep in.



% The paper headers
\markboth{Journal of \LaTeX\ Class Files,~Vol.~13, No.~9, September~2014}%
{Shell \MakeLowercase{\textit{et al.}}: Bare Demo of IEEEtran.cls for Journals}
% The only time the second header will appear is for the odd numbered pages
% after the title page when using the twoside option.
% 
% *** Note that you probably will NOT want to include the author's ***
% *** name in the headers of peer review papers.                   ***
% You can use \ifCLASSOPTIONpeerreview for conditional compilation here if
% you desire.




% If you want to put a publisher's ID mark on the page you can do it like
% this:
%\IEEEpubid{0000--0000/00\$00.00~\copyright~2014 IEEE}
% Remember, if you use this you must call \IEEEpubidadjcol in the second
% column for its text to clear the IEEEpubid mark.



% use for special paper notices
%\IEEEspecialpapernotice{(Invited Paper)}




% make the title area
\maketitle

% As a general rule, do not put math, special symbols or citations
% in the abstract or keywords.

{\begin{abstract}
Recent work established a generalized-Fano framework for lower bounding prior-predictive (Bayesian) CVaR in interactive statistical decision making. In this paper, we show how to instantiate that framework in concrete interactive problems and derive explicit Bayesian CVaR lower bounds from its abstract corollaries. Our approach compares a hard model with a reference model using squared Hellinger distance, and combines a lower bound on a reference hinge term with a bound on the distinguishability of the two models. We apply this approach to canonical examples, including Gaussian bandits, and obtain explicit bounds that make the dependence on key problem parameters transparent. These results show how the generalized-Fano Bayesian CVaR framework can be used as a practical lower-bound tool for interactive learning and risk-sensitive decision making.
\end{abstract}
}


% Note that keywords are not normally used for peerreview papers.
\begin{IEEEkeywords}
information theory and control, statistical learning
\end{IEEEkeywords}

\section{Introduction}
Lower bounds play a central role in statistical decision making by quantifying the smallest loss or cost that any problem must incur. In a Bayesian formulation of statistical decision making, one assumes a prior \(\mu\) on the model class \(\mathcal{M}\) and evaluates performance under the prior-predictive probability. The resulting Bayes risk provides a natural benchmark for what can be achieved on average under the prior, and lower bounds on this quantity reveal information-theoretic limits of learning and decision making. Classical tools such as Le Cam's method, Assouad's lemma, and Fano's inequality are fundamental techniques for deriving such lower bounds in passive statistical estimation \cite{tsybakov2009introduction}.

Many modern learning and control problems are interactive: the learner's actions affect what data are observed. Chen et al.~\cite{chen2024assouad} formalize this through the interactive statistical decision making (ISDM) framework. In ISDM, a model \(M\in\mathcal{M}\) specifies the environment, an algorithm \(\textup{Alg}\in\mathcal{D}\) chooses actions or decisions, and the resulting observation or transcript \(X\) is distributed according to an algorithm-dependent law
\(
X\sim P^{M,\textup{Alg}}.
\)
This algorithm dependence is the key distinction from passive estimation, where the data law is fixed and independent of the estimator. The ISDM viewpoint encompasses passive estimation, multi-armed bandits, and reinforcement learning within a common framework \cite{chen2024assouad,lattimore2020bandit,sutton2018reinforcement}. It is especially relevant in control and sequential decision making, where information acquisition and performance are inherently coupled: one must act in order to learn, and those actions shape both future information and eventual cost. Most lower-bound results in these settings focus on expected loss or expected regret. While expected risk is a natural measure of performance, it can conceal rare but costly failures. This motivates risk-sensitive criteria that emphasize the tail of the loss distribution. A standard example is Conditional Value-at-Risk,
\(
\textup{CVaR}_\alpha(L),
 \alpha\in[0,1),
\)
which measures the expected loss in the worst \((1-\alpha)\)-tail. For losses, \(\textup{CVaR}\) is more sensitive to adverse events than the mean and is widely used in finance, optimization, and risk-sensitive control \cite{rockafellar2000optimization,rockafellar2002general,acerbi2002expected}. Related tail criteria, such as quantiles and Value-at-Risk, have also been studied in statistical learning; see, for example, \cite{ma2024highprobability}. In contrast to Bayes expected-risk lower bounds, Bayesian \(\mathrm{CVaR}\) lower bounds quantify how small the tail-sensitive
performance criterion can be under the prior-predictive probability law, and hence what no algorithm can improve on. A recent step in this direction is the work of Bongole et al.~\cite{bongole2026generalizing}, which extends the interactive Fano method of Chen et al.~\cite{chen2024assouad}. They consider bounded transforms of the loss and derive a generalized interactive Fano inequality for expected bounded transforms under the prior-predictive law. As a consequence, they obtain a general lower-bound template for Bayesian \(\textup{CVaR}\) of bounded losses in ISDM. Their results show that tail-sensitive Bayesian performance criteria can be treated within the same information-theoretic framework as expected-risk and quantile-type criteria.

The main challenge left open by \cite{bongole2026generalizing} is instantiation. The corollaries identify the quantities: a benchmark term under a reference law and a divergence budget, but applying them in a concrete problem still requires a problem-specific calculation. In this work, the objective is to make this framework usable in concrete settings. To this end, a reusable two-point Bayesian \(\textup{CVaR}\) template is extracted from the generalized-Fano corollaries and then instantiated in both a passive Gaussian mean estimation problem and an interactive two-armed Gaussian bandit problem. These examples show how explicit worst-case Bayesian \(\textup{CVaR}\) lower bounds can be obtained and how the resulting bounds compare with the corresponding Bayesian expected-loss benchmarks.

\newgeometry{top=0.75in,left=0.75in,right=0.75in,bottom=0.75in}
Recent work on \(\mathrm{CVaR}\) has largely focused on risk-sensitive criteria that are different from the prior-predictive quantity studied here. In stochastic bandits, Baudry et al.~\cite{baudry2021optimal} evaluate each arm by the \(\mathrm{CVaR}\) of its reward distribution and analyze the corresponding cumulative regret relative to the arm with best \(\mathrm{CVaR}\). Tan and Weng~\cite{tan2023cvar} instead study upper bounds on the \(\mathrm{CVaR}\) of the regret random variable itself. Wang et al.~\cite{wang2023near} consider a cumulative \(\mathrm{CVaR}\)-regret criterion in bandits and reinforcement learning and analyze its expectation, so the \(\mathrm{CVaR}\)-based performance gap is defined within each environment and then aggregated. Du et al.~\cite{du2022provably} study an iterated-\(\mathrm{CVaR}\) objective in episodic reinforcement learning, which yields a dynamic nested risk measure. 

In passive settings, Thomas and Learned-Miller~\cite{thomas2019concentration} study estimation of \(\mathrm{CVaR}\) from i.i.d.\ samples, Soma and Yoshida~\cite{soma2020statistical} study learning with a population objective of the form \(\mathrm{CVaR}\) of the loss under a fixed data-generating distribution, and Mhammedi et al.~\cite{mhammedi2020pacbayes} derive PAC-Bayesian generalization bounds for empirical \(\mathrm{CVaR}\); in each case, \(\mathrm{CVaR}\) is evaluated under a single fixed model dependent law. By contrast, our setting is Bayesian: one first draws a model \(M\) from a prior \(\mu\), then generates the transcript under \(P^{M,\mathrm{Alg}}\), and finally evaluates the tail risk of the resulting loss under this prior-predictive mixture law. 

In this work, we study Bayesian \(\textup{CVaR}\) under the prior-predictive law, which is a natural risk-sensitive analogue of Bayes risk: it captures not only average performance under model uncertainty, but also the severity of unfavorable outcomes induced jointly by the prior and the algorithm. Our goal is to derive explicit information-theoretic lower bounds for this prior-predictive \(\textup{CVaR}\) criterion by instantiating the generalized-Fano framework of Bongole et al.~\cite{bongole2026generalizing}.



\subsection{Contributions}

The main contribution is to make the abstract Bayesian \(\textup{CVaR}\) lower-bound framework of \cite{bongole2026generalizing} concrete and usable in canonical problems. In particular:
\begin{itemize}
    \item We extract a reusable two-point Hellinger--\(\textup{CVaR}\) lower-bound template from the generalized interactive Fano corollaries. 

    \item We instantiate this template in both a passive and an interactive setting, namely two-point Gaussian mean estimation and the two-armed Gaussian bandit.

     \item For both examples, we derive explicit fixed-prior Bayesian \(\textup{CVaR}\) lower bounds that recover the same order as the corresponding Bayes-risk benchmarks, while also retaining dependence on the risk level \(\alpha\). In particular, these immediately imply corresponding worst-case (over priors) Bayesian \(\textup{CVaR}\) lower bounds.

    \item We clarify the relation between Bayesian \(\textup{CVaR}\) and Bayesian expected loss through the inequality \(\textup{CVaR}_\alpha[L]\ge \mathbb{E}[L]\), which provides a natural benchmark for interpreting the derived bounds and highlights what is gained by a tail-sensitive lower-bound analysis.


\end{itemize}







\iffalse
\section{Introduction}
%Talk about minimax risk and how people in RL use it.
The concept of minimax risk has become a staple in learning theory and statistics \cite{ma2024high}\cite{audibert2009minimax} \cite{azar2017minimax}.  In interactive or online learning settings the minimax risk is replaced by its counterpart the minimax regret. Both the minimax risk and regret intuitively capture what the best algorithm can achieve when faced with the worst case environment.

In interactive decision making problems such as bandits and reinforcement learning there have been considerable amount of work with the focus being either designing algorithms achieving near minimax-optimal performance or proving through mainly information theoretic principles, a lower bound on a minimax criterion. This lower bound on the minimax risk is a converse result i.e. it shows that
no algorithm can perform better than the proven lower bound.
The minimax risk widely used in these works mainly factors in the expected risk, where the expectation is taken under all possible trajectories under the interaction process.  Although widely used, with a lot of existing tools for analysis such as the Fano method \cite{cover1999elements,yu1997assouad,ma2024high,chen2016bayes,chen2024assouad, duchi2013distance}, the criterion misses important information that can be crucial in certain safety critical applications- it is agnostic to the tail behavior of the interaction process.
Therefore, a tool is required to study the minimax tail behavior in interactive processes.
One such tool is the minimax quantile studied in \cite{ma2024high}. 
The minimax quantile captures the tail behavior of the loss function under a given risk level. \cite{ma2024high} provide tools to study the minimax quantile lower bound using high probability variants of tools used to study the minimax risk such as the Fano method. Their analysis is however limited to estimation problems which do not involve interaction. 
\textcolor{red}{Maybe give one example demonstrating the tail risk behavior}

\cite{chen2024assouad} introduce ISDM as an easy way to include both passive estimation and interaction protocols. It is not that DMSO cannot include passive but ISDM does it easier without need to invent null actions.


In contrast \cite{chen2024assouad} give general techniques for the minimax risk in interactive problems through a notion of quantile probability bound however the minimax quantile behaviour is not of focus in their work. Their focus lies mainly in unifying techniques that lower bound the minimax risk such as the Fano method and Le Cam method \cite{lecam1973convergence, polyanskiy2014lecture} using the quantile lower bound while also providing a generalization of the technique that holds for interactive processes. Their focus also lies in providing lower bounds for the minimax expected risk. As far as we know, however explicit exploration of minimax quantile behaviour in interactive processes remains underexplored.
\textcolor{red}{Give citations to Le cam and Fano etc.}
\textcolor{red}{Write about high probability minimax statements from Lattimore for example}
Previous works have proven similar results i.e. high probability versions of minimax lower bounds \cite{lattimore2020bandit}. But work is scarse and there is no unifying framework and interactive versions of Fano's inequality that can be directly applied as far as we know that give risk level specific bounds on the minimax quantile.

The aim of this work is explore techniques that lower bound the minimax quantile behaviour of interactive problems. This work will serve as a stepping stone to gain larger understanding of the fundamental limits of tail behaviour in interactive problems. Inspired from \cite{ma2024high} that provide minimax quantile lower bounding techniques for non-interactive problems and \cite{chen2024assouad} that provide a general framework for interaction along unifying lower bounding techniques for minimax expected risk we aim to find counterparts of techniques used in non-interactive settings that hold even in interactive ones. Another unexplored territory is finding user friendly techniques that determine how the risk level can affect that tail behavior for various problems. Specifically we make the following contributions: 

%Talk about lower bounds for minimax risk using Fano etc.

%Talk about ISDM.
\fi

\subsection{Notation}

Random variables are denoted by capital letters, such as \(M\), \(X\), \(Y_t\), and \(H_T\); realizations are denoted by lowercase letters. We write \(\mathbb{E}\) for expectation and \(\mathbf{1}\{\cdot\}\) for indicators. For \(z\in\mathbb{R}\), let
\(
[z]_+ := \max\{z,0\}.
\)

For a model \(m\in\mathcal M\) and an algorithm \(\textup{Alg}\), we write
\(
P^{m,\textup{Alg}}
\)
for the law of the observation or transcript \(X\) induced by running \(\textup{Alg}\) in model \(m\). Equivalently,
\(
X\sim P^{m,\textup{Alg}}
\)
means that \(P^{m,\textup{Alg}}\) is the probability distribution of \(X\) given model \(m\) and algorithm \(\textup{Alg}\). In two-point constructions, we use the shorthand
\(
P_i:=P^{M_i,\textup{Alg}}, \qquad i\in\{1,2\}.
\)

For a prior \(\mu\in\mathcal{P}(\mathcal M)\), where \(\mathcal{P}(\mathcal M)\) is the set of all probability distributions on \(\mathcal M\), we write
\(
\mathbb{E}_{\mu,P^{M,\textup{Alg}}}[\cdot]
\)
for expectation under the joint law obtained by first drawing \(M\sim\mu\) and then \(X\sim P^{M,\textup{Alg}}\). Similarly, for a reference law \(Q\) on the transcript space, we write
\(
\mathbb{E}_{\mu,Q}[\cdot]
\)
when \(M\sim\mu\) and \(X\sim Q\) are independent. 
We write \(\textup{Bern}(p)\) for the Bernoulli distribution with mean \(p\). For probability measures \(P\) and \(Q\) on a common measurable space, \(D_{\textup{KL}}(P\|Q)\) denotes the Kullback--Leibler divergence, and \(D_{H^2}(P\|Q)\) denotes the squared Hellinger divergence, where
\(
D_{\textup{KL}}(P\|Q)
:=
\int \log\!\left(\frac{dP}{dQ}\right)\,dP \text{ and }
D_{H^2}(P\|Q)
:=
1-\int \sqrt{\frac{dP}{d\lambda}\frac{dQ}{d\lambda}}\,d\lambda,
\)
where \(\lambda\) is any common dominating measure. We use the standard inequality
\(
D_{H^2}(P\|Q)\le D_{\textup{KL}}(P\|Q).
\)

For a loss random variable \(L\), \(\textup{CVaR}_\alpha(L)\) denotes the upper-tail Conditional Value-at-Risk at level \(\alpha\in[0,1)\), defined by
\(
\textup{CVaR}_\alpha(L)
:=
\min_{t\in\mathbb{R}}
\left\{
t+\frac{1}{1-\alpha}\,\mathbb{E}[(L-t)_+]
\right\}.
\)


\section{Problem setup}

Interactive statistical decision making (ISDM) provides a common language for passive estimation, sequential decision making, and learning under uncertainty. Following \cite{chen2024assouad}, an ISDM problem is specified by
\begin{itemize}
    \item a model class \(\mathcal{M}\),
    \item an algorithm class \(\mathcal{D}\),
    \item a transcript space \(\mathcal{X}\),
    \item and a nonnegative loss function
    \(
    L:\mathcal{M}\times \mathcal{X}\to \mathbb{R}_+.
    \)
\end{itemize}

For each model \(M\in\mathcal{M}\) and algorithm \(\textup{Alg}\in\mathcal{D}\), the interaction between the algorithm and the model induces a law
\(
X\sim P^{M,\textup{Alg}}
\)
on the transcript \(X\in\mathcal{X}\). Thus:
\begin{itemize}
    \item in a passive estimation problem, if one observes samples \(Y_1,\dots,Y_n\) and then outputs an estimate \(\widehat\theta=\textup{Alg}(Y_1,\dots,Y_n)\), the transcript \(X\) may be taken to include both the batch of observations and the final decision, e.g.
    \(
    X=(Y_1,\dots,Y_n,\widehat\theta);
    \)
    \item in a sequential decision problem, \(X\) may be the full interaction history, including actions, observations, and any terminal decision, for example
    \(
    X=(A_1,Y_1,\dots,A_T,Y_T,\widehat d).
    \)
\end{itemize}
This formulation therefore covers both passive and interactive settings within the same framework.

Under a prior \(\mu\in\mathcal{P}(\mathcal{M})\), the standard Bayesian performance criterion is the Bayes risk
\(
\mathbb{E}_{\mu,P^{M,\textup{Alg}}}[L(M,X)].
\)
Much of the existing lower-bound literature focuses on this expected loss. In this work, the focus is instead on the risk-sensitive criterion
\(
\textup{CVaR}^{\,\mu\otimes P^{M,\mathrm{Alg}}}_\alpha(L) =
\)
\[
\min_{t\in\mathbb{R}}
\left\{
t+\frac{1}{1-\alpha}
\int (L(m,x)-t)_+\,(\mu\otimes P^{M,\mathrm{Alg}})(dm,dx)
\right\}.
\]

 This criterion captures tail-sensitive performance and is especially relevant when rare but costly failures matter.

The main aim in this work is to turn the abstract Bayesian \(\textup{CVaR}\) lower-bound framework of \cite{bongole2026generalizing} into concrete and usable converse results. The emphasis is on showing how its corollaries can be instantiated in canonical problems. In particular, the paper develops a reusable two-point Bayesian \(\textup{CVaR}\) template and applies it in both a passive example and an interactive example, thereby illustrating how explicit \(\textup{CVaR}\) lower bounds can be derived directly from the generalized interactive Fano method.

%---------------------------------------------
\iffalse
%---------------------------------------------
\section{Prior results}

\begin{lemma}[Data processing lemma]
    Let \(P_X,Q_X\) be laws on \((\mathcal X,\mathcal F)\). 
Let \(U\sim\nu\) be independent of \(X\), and let \(\Phi:\mathcal X\times\mathcal U\to\mathcal Y\) be measurable.
Define the (randomized) kernel
\[
T(A\mid x)\;=\;\int \mathbf 1\{\Phi(x,u)\in A\}\,\nu(du),\qquad A\in\mathcal G.
\]
Then for any \(f\)-divergence,
\[
D_f(P_X\|Q_X)\ \ge\ D_f\big(P_X T\ \|\ Q_X T\big)
\]
\[=D_f\!\Big(\mathsf{Law}\big(\Phi(X,U)\,\big|\,X\sim P_X\big)\ \Big\|\ 
\mathsf{Law}\big(\Phi(X,U)\,\big|\,X\sim Q_X\big)\Big).
\]
\end{lemma}

Chen et al. \cite{chen2024assouad} prove the following lemma.
\begin{lemma}[Bernoulli monotonicity]
\label{lem:bern-mono}
For fixed $y\in[0,1]$, the map $x\mapsto D_f(\mathrm{Bern}(x)\,\|\,\mathrm{Bern}(y))$
is nondecreasing on $[y,1]$ and nonincreasing on $[0,y]$.
\end{lemma}

\begin{theorem}[Fano inequality for expected risk]
\label{thm:Fano-expected-risk}
Fix an $f$-divergence $D_f$. Let $\textup{ALG}$ be a given algorithm, let $\mu\in\Delta(\mathcal M)$ be a prior, and let $\Delta>0$ be a risk level. For any reference $Q\in\Delta(\mathcal X)$ define
\[
\rho_{\Delta,Q}\;:=\; \mathbb P_{M\sim\mu,\; X\sim Q}\!\big(L(M,X)<\Delta\big).
\]
Then
\[
\begin{aligned}
&\mathbb P_{M\sim\mu,\; X\sim P^{M,\textup{ALG}}}\!\big(L(M,X)\ge \Delta\big)
\ge
\\
&\sup_{Q\in\Delta(\mathcal X),\; \delta\in[0,1]} \Big\{\,
\delta \;:\; \mathbb{E}_{M\sim\mu}\!\big[D_f\!\big(P^{M,\textup{ALG}}\,\|\,Q\big)\big]
\;<\; d_{f,\delta}\!\big(\rho_{\Delta,Q}\big)
\,\Big\},
\end{aligned}
\]
where
\[
d_{f,\delta}(p)\;:=\;
\begin{cases}
D_f\!\big(\mathrm{Bern}(1-\delta)\,\big\|\,\mathrm{Bern}(p)\big), & \text{if } p\le 1-\delta,\\[4pt]
0, & \text{otherwise.}
\end{cases}
\]
\end{theorem}

\begin{proof}[Proof]
We work with the two transcript laws
$P_0(m,x)=\mu(m)P^{M,\textup{ALG}}(x)$ and $P_1(m,x)=\mu(m)Q(x)$.
The argument has three steps.

\emph{i) (indicator reduction).}
Compress $(M,X)$ to the success bit
$S=\mathbf 1\{L(M,X)<\Delta\}$. Then, 
\[
\mathbb P_{(M,X) \sim P_i}(S=1)=\mathbb P_{(M,X) \sim P_i}(L(M,X)< \Delta)
\]

Hence,
$\textup{Law}(S\mid P_0)=\mathrm{Bern}(\bar\rho_\Delta)$ and
$\textup{Law}(S\mid P_1)=\mathrm{Bern}(\rho_{\Delta,Q})$ with
$\bar\rho_\Delta=\mathbb P_{P_0}(L<\Delta)$ and
$\rho_{\Delta,Q}=\mathbb P_{P_1}(L<\Delta)$.

\emph{ii) (data processing).}
By the data-processing lemma for $f$-divergences applied to the indicator
kernel $(m,x)\mapsto S$,
\begin{equation}
\label{eq:dp-bern}
\begin{aligned}
    &D_f\!\big(\mathrm{Bern}(\bar\rho_\Delta)\,\big\|\,\mathrm{Bern}(\rho_{\Delta,Q})\big)
    \ \le\
    D_f(P_0\|P_1)
    =
    \\
    &\mathbb{E}_{M\sim\mu} D_f\!\big(P^{M,\textup{ALG}}\,\|\,Q\big).
\end{aligned}
\end{equation}


\emph{iii) (Bernoulli reduction).}

Fix $Q$ and $\delta\in[0,1]$ such that
\[
\begin{aligned}
&\mathbb{E}_{M\sim\mu} D_f\!\big(P^{M,\textup{ALG}}\,\|\,Q\big)\ <\ d_{f,\delta}(\rho_{\Delta,Q}) =
\\
&\begin{cases}
D_f\!\big(\mathrm{Bern}(1-\delta)\,\big\|\,\mathrm{Bern}(\rho_{\Delta,Q})\big), & \rho_{\Delta,Q}\le 1-\delta,\\
0, & \rho_{\Delta,Q}>1-\delta.
\end{cases}
\end{aligned}
\]
Since the left-hand side is nonnegative, the strict inequality forces $\rho_{\Delta,Q}<1-\delta$.
Combining with \eqref{eq:dp-bern} gives
\[
D_f\!\big(\mathrm{Bern}(\bar\rho_\Delta)\,\big\|\,\mathrm{Bern}(\rho_{\Delta,Q})\big)
\ <\
D_f\!\big(\mathrm{Bern}(1-\delta)\,\big\|\,\mathrm{Bern}(\rho_{\Delta,Q})\big).
\]

\emph{Case 1: $\bar\rho_\Delta \le \rho_{\Delta,Q}$.}
Then $\bar\rho_\Delta \le \rho_{\Delta,Q} < 1-\delta$, hence directly $\bar\rho_\Delta<1-\delta$.

\emph{Case 2: $\bar\rho_\Delta > \rho_{\Delta,Q}$.}


Using lemma \ref{lem:bern-mono} with $y=\rho_{\Delta,Q}$ and $1-\delta\ge \rho_{\Delta,Q}$, the strict inequality above implies
$\bar\rho_\Delta < 1-\delta$.

In both cases we conclude $\bar\rho_\Delta<1-\delta$, i.e.,
\[
\mathbb P_{M\sim\mu,\,X\sim P^{M,\textup{ALG}}}\!\big(L(M,X)\ge \Delta\big)=1-\bar\rho_\Delta>\delta.
\]
Taking the supremum over all such $(Q,\delta)$ completes the proof.


\begin{remark}[From tail to expectation]
For any nonnegative loss, $\mathbb{E}[L]\ge \Delta\,\mathbb P(L\ge \Delta)$, so Theorem~\ref{thm:Fano-expected-risk}
immediately implies an \emph{expected-risk} lower bound:
\[
\begin{aligned}
&\mathbb{E}_{M\sim\mu,\,X\sim P^{M,\textup{ALG}}}[L(M,X)]
\ \ge\
\\
&\Delta\cdot
\sup_{Q,\delta}\Big\{\delta:\ \mathbb{E}_M D_f(P^{M,\textup{ALG}}\|Q)< d_{f,\delta}(\rho_{\Delta,Q})\Big\}.
\end{aligned}
\]
\end{remark}
\fi
\iffalse
\section{Prior results}
Chen et al.~\cite{chen2024assouad} established an interactive Fano inequality by combining data processing for \(f\)-divergences with a Bernoulli monotonicity argument. In this section, we re-derive their result with an eye toward generalization, isolating the proof steps that can be extended and those that are specific to the original formulation.


\begin{lemma}[Data processing inequality for \(f\)-divergences \cite{polyanskiy2025information}]\label{lem:dpi-f}
Let \((\mathcal X,\mathcal F)\), \((\mathcal U,\mathcal H)\), and \((\mathcal Y,\mathcal G)\) be measurable spaces.
Let \(P_X,Q_X\) be probability measures on \((\mathcal X,\mathcal F)\), and let \(\nu\) be a probability measure on \((\mathcal U,\mathcal H)\).
Let \(\Phi:(\mathcal X\times\mathcal U,\mathcal F\otimes\mathcal H)\to(\mathcal Y,\mathcal G)\) be measurable.

Define a Markov kernel \(T:\mathcal X\times\mathcal G\to[0,1]\) by
\[
T(A\mid x)\;:=\;\int_{\mathcal U}\mathbf 1\{\Phi(x,u)\in A\}\,\nu(du),
\qquad x\in\mathcal X,\; A\in\mathcal G.
\]
Let \(T\circ P_X\) denote the pushforward (output) measure on \((\mathcal Y,\mathcal G)\) induced by \(P_X\) through \(T\),
\[
(T\circ P_X)(A)\;:=\;\int_{\mathcal X}T(A\mid x)\,P_X(dx),\qquad A\in\mathcal G,
\]
and similarly define \(T\circ Q_X\).

Then for any \(f\)-divergence \(D_f\),
\[
D_f(P_X\|Q_X)\ \ge\ D_f\big(T\circ P_X\ \|\ T\circ Q_X\big).
\]
Equivalently, if \((X,U)\sim P_X\otimes \nu\) and \((X',U)\sim Q_X\otimes \nu\), then
\[
T\circ P_X \;=\; \textup{Law}_{P_X\otimes \nu}\big(\Phi(X,U)\big)
\]
\[
T\circ Q_X \;=\; \textup{Law}_{Q_X\otimes \nu}\big(\Phi(X',U)\big),
\]
and hence
\[
D_f(P_X\|Q_X)\ \ge\ \]
\[
D_f\!\Big(\textup{Law}_{P_X\otimes \nu}\big(\Phi(X,U)\big)\ \Big\|\ 
\textup{Law}_{Q_X\otimes \nu}\big(\Phi(X',U)\big)\Big).
\]
\end{lemma}

Chen et al.~\cite{chen2024assouad} prove the following lemma.


\begin{lemma}[Bernoulli monotonicity]
\label{lem:bern-mono}
For fixed \(y\in[0,1]\), the map \(x\mapsto D_f(\mathrm{Bern}(x)\,\|\,\mathrm{Bern}(y))\)
is nondecreasing on \([y,1]\) and nonincreasing on \([0,y]\).
\end{lemma}
Based on the above lemmas, Chen et al.~\cite{chen2024assouad} prove an interactive Fano inequality
which says if, for some reference distribution $Q$, the average $f$-divergence
\(
\mathbb{E}_{M\sim\mu}\!\left[D_f\!\left(P^{M,\mathrm{ALG}}\,\|\,Q\right)\right]
\)
stays below a certain threshold $d_{f,\delta}\!\left(\rho_{\Delta,Q}\right)$, then the algorithm must incur loss at least $\Delta$ with probability at least $\delta$, i.e.,
\[
\mathbb{E}_{M\sim\mu}\!\left[D_f\!\left(P^{M,\mathrm{ALG}}\,\|\,Q\right)\right]
< d_{f,\delta}\!\left(\rho_{\Delta,Q}\right)
\]
\[
\quad\Longrightarrow\quad
\mathbb{P}_{M\sim\mu,\;X\sim P^{M,\mathrm{ALG}}}\!\big(L(M,X)\ge \Delta\big)\ge \delta.
\]
\begin{theorem}[Interactive quantile Fano]
\label{thm:Fano-expected-risk}
Fix an \(f\)-divergence \(D_f\). Let \(\textup{ALG}\) be a given algorithm, let \(\mu\in\Delta(\mathcal M)\) be a prior, and let \(\Delta>0\) be a risk level. For any reference \(Q\in\Delta(\mathcal X)\) define
\(
\rho_{\Delta,Q}\;:=\; \mathbb P_{M\sim\mu,\; X\sim Q}\!\big(L(M,X)<\Delta\big).
\)
Then we have the following quantile lower bound 
\[
\begin{aligned}
&\mathbb P_{M\sim\mu,\; X\sim P^{M,\textup{ALG}}}\!\big(L(M,X)\ge \Delta\big)
\ge
\\
&\sup_{Q\in\Delta(\mathcal X),\; \delta\in[0,1]} \{\,
\delta \;:\; \mathbb{E}_{M\sim\mu}\!\big[D_f\!\big(P^{M,\textup{ALG}}\,\|\,Q\big)\big]
\!<\! d_{f,\delta}\!\big(\rho_{\Delta,Q}\big)
\,\},
\end{aligned}
\]
where
\[
d_{f,\delta}(p)\;:=\;
\begin{cases}
D_f\!\big(\mathrm{Bern}(1-\delta)\,\big\|\,\mathrm{Bern}(p)\big), & \text{if } p\le 1-\delta,\\[4pt]
0, & \text{otherwise.}
\end{cases}
\]
\end{theorem}

\begin{proof}[Proof]
We work with the two joint laws
\[
P_0(m,x)\;=\;\mu(m)P^{M,\textup{ALG}}(x), \qquad 
P_1(m,x)\;=\;\mu(m)Q(x).
\]
The argument proceeds in three steps.

\emph{i) (Indicator reduction).}
Compress \((M,X)\) to the success indicator
\[
S\;:=\;\mathbf 1\{L(M,X)<\Delta\}.
\]
Then, for \(i\in\{0,1\}\),
\[
\mathbb P_{(M,X)\sim P_i}(S=1)
= \mathbb P_{(M,X)\sim P_i}\!\big(L(M,X)<\Delta\big).
\]
Hence, under \(P_0\),
\[
S\ \sim\ \mathrm{Bern}(\bar\rho_\Delta),
\qquad 
\bar\rho_\Delta\;:=\;P_0\!\big(L(M,X)<\Delta\big),
\]
and under \(P_1\),
\[
S\ \sim\ \mathrm{Bern}(\rho_{\Delta,Q}),
\qquad 
\rho_{\Delta,Q}\;:=\;P_1\!\big(L(M,X)<\Delta\big).
\]
Equivalently,
\[
\mathsf{Law}_{P_0}(S)=\mathrm{Bern}(\bar\rho_\Delta),
\qquad
\mathsf{Law}_{P_1}(S)=\mathrm{Bern}(\rho_{\Delta,Q}).
\]

\emph{ii) (Data processing).}
By the data-processing inequality for \(f\)-divergences applied to the measurable map \((m,x)\mapsto S\),
\begin{equation}
\label{eq:dp-bern}
\begin{aligned}
    &D_f\!\big(\mathrm{Bern}(\bar\rho_\Delta)\,\big\|\,\mathrm{Bern}(\rho_{\Delta,Q})\big)
    \ \le\
    D_f(P_0\|P_1)
    \\
    &=\
    \mathbb{E}_{M\sim\mu}\, D_f\!\big(P^{M,\textup{ALG}}\,\|\,Q\big).
\end{aligned}
\end{equation}

\emph{iii) (One-sided Inversion).}

Fix \(Q\) and \(\delta\in[0,1]\) such that
\[
\begin{aligned}
&\mathbb{E}_{M\sim\mu}\, D_f\!\big(P^{M,\textup{ALG}}\,\|\,Q\big)
\ <\ d_{f,\delta}(\rho_{\Delta,Q})
\\
&=
\begin{cases}
D_f\!\big(\mathrm{Bern}(1-\delta)\,\big\|\,\mathrm{Bern}(\rho_{\Delta,Q})\big), & \rho_{\Delta,Q}\le 1-\delta,\\
0, & \rho_{\Delta,Q}>1-\delta.
\end{cases}
\end{aligned}
\]
Since the left-hand side is nonnegative, the strict inequality forces \(\rho_{\Delta,Q}<1-\delta\).
Combining with \eqref{eq:dp-bern} yields
\[
D_f\!\big(\mathrm{Bern}(\bar\rho_\Delta)\big\|\,\mathrm{Bern}(\rho_{\Delta,Q})\big)
\!\!<\!\!\!\
D_f\!\big(\mathrm{Bern}(1-\delta)\,\big\|\,\mathrm{Bern}(\rho_{\Delta,Q})\big).
\]

\emph{Case 1:} \(\bar\rho_\Delta \le \rho_{\Delta,Q}\).
Then \(\bar\rho_\Delta \le \rho_{\Delta,Q} < 1-\delta\), hence directly \(\bar\rho_\Delta<1-\delta\).

\emph{Case 2:} \(\bar\rho_\Delta > \rho_{\Delta,Q}\).
Using Lemma~\ref{lem:bern-mono} with \(y=\rho_{\Delta,Q}\) and \(1-\delta\ge \rho_{\Delta,Q}\), the strict inequality above implies
\(\bar\rho_\Delta < 1-\delta\).

In both cases we conclude \(\bar\rho_\Delta<1-\delta\), i.e.,
\[
\mathbb P_{M\sim\mu,\,X\sim P^{M,\textup{ALG}}}\!\big(L(M,X)\ge \Delta\big)
= 1-\bar\rho_\Delta > \delta.
\]
Taking the supremum over all such pairs \((Q,\delta)\) completes the proof.
\end{proof}
\begin{remark}[Discussion: proof steps in the interactive quantile Fano]
The interactive Fano argument for quantile risk can be organized into three conceptual steps:
(i) an \emph{indicator reduction} that compresses the loss into an indicator,
(ii) a \emph{data-processing} bound,
and (iii) an \emph{inversion} step that translates the bound into a lower bound on the risk.
This decomposition also indicates how to extend the method. For more general loss functionals, one must replace the indicator reduction and the resulting \emph{one-sided} inversion by a construction that supports a \emph{two-sided} inversion.
\end{remark}


The tail risk bound gives rise to an in-expectation bound from Markov's inequality.
\begin{remark}[From tail to expectation]
For any nonnegative loss, \(\mathbb{E}[L]\ge \Delta\,\mathbb P(L\ge \Delta)\), so Theorem~\ref{thm:Fano-expected-risk}
immediately implies an \emph{expected-risk} lower bound:
\[
\begin{aligned}
&\mathbb{E}_{M\sim\mu,\,X\sim P^{M,\textup{ALG}}}[L(M,X)]
\ \ge\
\\
&\Delta\cdot
\sup_{Q,\delta}\Big\{\delta:\ \mathbb{E}_M D_f(P^{M,\textup{ALG}}\|Q)< d_{f,\delta}(\rho_{\Delta,Q})\Big\}.
\end{aligned}
\]
\end{remark}
\fi

\subsection{Previous results}
\label{sec:previous-results}

The analysis below builds on the Bayesian \(\textup{CVaR}\) lower-bound framework of \cite{bongole2026generalizing}. We recall the two corollaries that are used in our instantiations.

\begin{corollary}[Bounded-hinge lower bound {\cite[Cor.~2]{bongole2026generalizing}}]
\label{cor:hinge-lower}
Assume \(0\le L(M,X)\le L_{\max}\), and fix \(t\in\mathbb R\). Define
\[
\phi_t(\ell):=\frac{(\ell-t)_+}{L_{\max}}\in[0,1],
\]
\[
b_t
:=
\mathbb{E}_{\mu,Q}[\phi_t(L)]
=
\frac{1}{L_{\max}}\mathbb{E}_{\mu,Q}[(L-t)_+].
\]
Let
\[
B:=\mathbb{E}_{M\sim\mu}D_f(P^{M,\textup{Alg}}\|Q),
\]
and
\[
a_f^{-}(B;b)
:=
\inf\Bigl\{a\in[0,1]:
D_f\big(\textup{Bern}(a)\,\big\|\,\textup{Bern}(b)\big)\le B
\Bigr\}.
\]
Then
\[
\mathbb{E}_{\mu,P^{M,\textup{Alg}}}[(L-t)_+]
\ge
L_{\max}\,a_f^{-}(B;b_t).
\]
\end{corollary}

\begin{corollary}[Bayesian \(\textup{CVaR}\) lower bound {\cite[Cor.~3]{bongole2026generalizing}}]
\label{cor:cvar-lower-bt}
Under the assumptions of Corollary~\ref{cor:hinge-lower},
\[
\textup{CVaR}_\alpha(L)
\ge
\min_{t\in\mathbb R}
\left\{
t+\frac{L_{\max}}{1-\alpha}\,a_f^{-}(B;b_t)
\right\}.
\]
\end{corollary}

Corollaries~\ref{cor:hinge-lower}--\ref{cor:cvar-lower-bt} reduce the derivation of explicit tail-sensitive lower bounds to a concrete instantiation problem: choose a reference law \(Q\), lower bound the corresponding hinge term \(b_t\), and upper bound the divergence budget \(B\). 

\section{Results}
\label{sec:results}
This section presents the reusable two-point Hellinger--\(\textup{CVaR}\) bound and the baseline comparison with Bayes risk that underlie the later examples. Full proofs are deferred to Appendix~\ref{app:proofs}. 

The first result is a reusable two-point Hellinger--\(\textup{CVaR}\) bound that will be instantiated in the examples below. 

\begin{lemma}[Two-point Hellinger--\(\textup{CVaR}\) lower bound]
\label{lem:two-point-hellinger-cvar}
Let \(\mathcal M=\{M_1,M_2\}\), let
\[
\mu=\frac12\delta_{M_1}+\frac12\delta_{M_2},
\]
and fix an algorithm \(\textup{Alg}\). Let
\[
P_i:=P^{M_i,\textup{Alg}},\qquad i=1,2.
\]
Let \(L:\mathcal M\times\mathcal X\to[0,L_{\max}]\) be a bounded loss. Assume that, for some \(C\in[0,2L_{\max}]\),
\[
L(M_1,x)+L(M_2,x)\ge C
\qquad
\text{for all }x\in\mathcal X,
\tag{H1}
\]
and that
\[
D_{H^2}(P_1\|P_2)\le \Gamma_H.
\tag{H2}
\]
Then, for every \(\alpha\in[0,1)\),
\[
\textup{CVaR}_\alpha^{\mu\otimes P^{M,\textup{Alg}}}(L)
\ge\]
\[
\min_{t\in\mathbb R}
\left\{
t+\frac{L_{\max}}{1-\alpha}
\left(
\sqrt{\left(\frac{C/2-t}{L_{\max}}\right)_+}
-\sqrt{\Gamma_H}
\right)_+^2
\right\}.
\]
\end{lemma}

\textit{Proof sketch:}
Apply Corollary~\ref{cor:cvar-lower-bt} with \(f=H^2\) and reference law \(Q=P_2\). Use \((\mathrm{H1})\) to lower bound the hinge benchmark \(b_t\), use \((\mathrm{H2})\) to bound the Hellinger budget, and then apply the Bernoulli Hellinger inversion
\[
a^-_{H^2}(B;b)\ge (\sqrt b-\sqrt{2B})_+^2.
\]
\qed

A simplified version is convenient for the balanced two-point constructions used below.

\begin{corollary}[Balanced-pair Hellinger simplification]
\label{cor:balanced-two-point-hellinger}
Under the assumptions of Lemma~\ref{lem:two-point-hellinger-cvar}, if in addition
\[
L(M_1,x)+L(M_2,x)\ge L_{\max}
\qquad
\text{for all }x\in\mathcal X,
\tag{H3}
\]
then, for every \(\alpha\in[0,1)\),
\[
\textup{CVaR}_\alpha^{\mu\otimes P^{M,\textup{Alg}}}(L)
\ge
\]
\[
\min_{t\in\mathbb R}
\left\{
t+\frac{L_{\max}}{1-\alpha}
\left(
\sqrt{\left(\frac12-\frac{t}{L_{\max}}\right)_+}
-\sqrt{\Gamma_H}
\right)_+^2
\right\}.
\]
\end{corollary}

\textit{Proof sketch:}
Set \(C=L_{\max}\) in Lemma~\ref{lem:two-point-hellinger-cvar}. \qed

The next proposition gives a baseline comparison with Bayesian expected loss.

\begin{proposition}[Bayesian \(\textup{CVaR}\) dominates Bayesian expected loss]
\label{prop:cvar-dominates-bayes-risk}
Let \(L\) be any integrable loss, and let \(\alpha\in[0,1)\). Then
\(
\textup{CVaR}_\alpha(L)\ge \mathbb{E}[L].
\)
Consequently, for any prior \(\mu\),
\[
\textup{CVaR}_\alpha^{\mu\otimes P^{M,\textup{Alg}}}(L)
\ge
\mathbb{E}_{\mu\otimes P^{M,\textup{Alg}}}[L].
\]
In particular, if for some prior \(\mu\) and some quantity \(\psi>0\) one has
\[
\mathbb{E}_{\mu\otimes P^{M,\textup{Alg}}}[L]\ge c\,\psi
\]
for a constant \(c>0\), then
\[
\textup{CVaR}_\alpha^{\mu\otimes P^{M,\textup{Alg}}}(L)\ge c\,\psi.
\]
Moreover,
\[
\sup_{\mu}\textup{CVaR}_\alpha^{\mu\otimes P^{M,\textup{Alg}}}(L)
\ge
\sup_{\mu}\mathbb{E}_{\mu\otimes P^{M,\textup{Alg}}}[L].
\]
\end{proposition}


\section{Examples}
\label{sec:examples}

We now instantiate the two-point Hellinger template in a passive and an interactive Gaussian problem. In both cases, the lower bound is obtained by restricting the full problem class to a carefully chosen two-model subfamily and applying the fixed-prior two-point result on that hard pair. Since a lower bound for one fixed prior immediately implies a lower bound for the supremum over priors, the resulting theorems also yield worst-case (over priors) Bayesian \(\textup{CVaR}\) lower bounds. For readability, we state the main bounds here and defer the full proofs to Appendix~\ref{app:proofs}.

Both examples lead to the same one-dimensional minimization, so we first introduce the closed-form function that appears in the resulting bounds. For \(\alpha\in[0,1)\), define
\[
\Psi_\alpha(\rho):=
\begin{cases}
\displaystyle
\frac12-\frac{\rho^2}{2\alpha},
& 0\le \rho\le \alpha,\ \alpha>0,\\[2ex]
\displaystyle
\frac{(1-\rho)^2}{2(1-\alpha)},
& \alpha<\rho\le 1,\\[2ex]
0,
& \rho\ge 1,
\end{cases}
\]
and, for \(\alpha=0\),
\[
\Psi_0(\rho):=\frac{(1-\rho)_+^2}{2}.
\]
The function \(\Psi_\alpha\) is continuous at the breakpoints \(\rho=\alpha\) and \(\rho=1\); these branch values are obtained by the scalar minimization carried out in Lemma~\ref{lem:scalar-min} in the appendix. We will also use the optimized constant
\[
c_\alpha:=
\begin{cases}
\displaystyle
\frac{2}{27(1-\alpha)},
& 0\le \alpha\le \frac13,\\[2ex]
\displaystyle
\frac{\sqrt{\alpha}}{3\sqrt{3}},
& \frac13\le \alpha<1.
\end{cases}
\]

\subsection{A passive Gaussian mean estimation example}

We begin with a passive Gaussian example, where the transcript consists of the observations together with the final estimator output.

Consider the Gaussian mean estimation problem in which one observes
\[
Y_1,\dots,Y_n \;\;\textup{i.i.d.}\;\; \mathcal N(\theta,1),
\]
with unknown parameter \(\theta\in\mathbb R\), and outputs an estimator
\[
\widehat\theta=\widehat\theta(Y_1,\dots,Y_n).
\]
In the ISDM formulation, the transcript is
\[
X=(Y_1,\dots,Y_n,\widehat\theta).
\]

To apply the two-point Hellinger template, we restrict attention to the hard two-point subfamily
\[
\theta\in\{-\Delta,+\Delta\},
\]
where \(\Delta>0\) will be chosen later. Let
\[
M_+ := M_{+\Delta},\qquad M_-:=M_{-\Delta},
\qquad
\mu:=\frac12\delta_{M_+}+\frac12\delta_{M_-}.
\]
We equip this subfamily with the bounded estimation loss
\[
L(\theta,X):=\min\{|\widehat\theta-\theta|,\;2\Delta\}.
\]

The following theorem gives a fixed-prior Bayesian lower bound on this hard subfamily and hence also a worst-case Bayesian lower bound for the Gaussian mean estimation problem.

\begin{theorem}[Gaussian mean estimation: \(\alpha\)-dependent Bayesian \(\textup{CVaR}\) lower bound]
\label{prop:gaussian-mean-hellinger-cvar}
For every estimator \(\widehat\theta\) and every \(\alpha\in[0,1)\),
\[
\textup{CVaR}_\alpha^{\mu\otimes P^{M,\textup{Alg}}}(L)
\ge
2\Delta\,\Psi_\alpha(2\sqrt{n}\,\Delta).
\]
Equivalently,
\[
\textup{CVaR}_\alpha^{\mu\otimes P^{M,\textup{Alg}}}(L)
\ge\]
\[
\min_{t\in\mathbb R}
\left\{
t+\frac{2\Delta}{1-\alpha}
\left(
\sqrt{\left(\frac{\Delta-t}{2\Delta}\right)_+}
-\sqrt{2n}\,\Delta
\right)_+^2
\right\}.
\]
In particular, optimizing over \(\Delta\) and \(\nu\) gives the worst-case lower bound
\[
\sup_{\nu \in \mathcal{P}(\mathcal{M})}\textup{CVaR}_\alpha^{\nu\otimes P^{M,\textup{Alg}}}(L)
\ge
\sup_{\Delta>0}\textup{CVaR}_\alpha^{\mu\otimes P^{M,\textup{Alg}}}(L)
\ge
\frac{c_\alpha}{\sqrt n}.
\]
\end{theorem}

\textit{Proof sketch:}
The proof has three steps. First, identify the bounded loss level \(L_{\max}=2\Delta\). Second, verify the balanced-pair condition via the triangle inequality:
\[
L(+\Delta,x)+L(-\Delta,x)\ge 2\Delta.
\]
Third, bound the Hellinger divergence between the transcript laws induced by \(M_+\) and \(M_-\) using the Gaussian KL comparison,
\[
D_{H^2}(P^{M_+,\textup{Alg}}\|P^{M_-,\textup{Alg}})
\le
D_{\textup{KL}}(P^{M_+,\textup{Alg}}\|P^{M_-,\textup{Alg}})
=
2n\Delta^2,
\]
and apply Corollary~\ref{cor:balanced-two-point-hellinger}. The closed form then follows from Lemma~\ref{lem:scalar-min}, and optimizing over \(\Delta\) and subsequently over \(\nu\) yields the stated worst-case \(n^{-1/2}\) lower bound. \qed

\subsection{Two-armed Gaussian bandits}

We then turn to an interactive Gaussian example, where the transcript is the full action--observation history.

Consider the two-armed Gaussian bandit problem with unit-variance rewards and horizon \(T\). At each round \(t\), the algorithm \(\textup{Alg}\) chooses an arm \(A_t\in\{1,2\}\) based on the past history and then observes a reward \(Y_t\). The resulting transcript is
\[
H_T=(A_1,Y_1,\dots,A_T,Y_T).
\]

To apply the two-point Hellinger template, we restrict attention to the two-environment subfamily
\[
M_1:=\left(+\frac{g}{2},-\frac{g}{2}\right),
\qquad
M_2:=\left(-\frac{g}{2},+\frac{g}{2}\right),
\]
where \(g>0\) will be chosen later, and equip this subfamily with the symmetric prior
\[
\mu=\frac12\delta_{M_1}+\frac12\delta_{M_2}.
\]
For \(a\in\{1,2\}\), let \(N_a(H_T)\) denote the number of pulls of arm \(a\) up to time \(T\).

The following theorem gives a fixed-prior Bayesian lower bound on this hard two-environment family and hence also a worst-case Bayesian lower bound for the two-armed Gaussian bandit problem.

\begin{theorem}[Two-armed Gaussian bandits: \(\alpha\)-dependent worst-case Bayesian \(\textup{CVaR}\) lower bound]
\label{thm:two-armed-gaussian-hellinger-cvar}
For every algorithm \(\textup{Alg}\) and every \(\alpha\in[0,1)\),
\[
\textup{CVaR}_\alpha^{\mu\otimes P^{M,\textup{Alg}}}(R_T)
\ge
gT\,\Psi_\alpha(g\sqrt{T}).
\]
Equivalently,
\[
\textup{CVaR}_\alpha^{\mu\otimes P^{M,\textup{Alg}}}(R_T)
\ge\]
\[
\min_{t\in\mathbb R}
\left\{
t+\frac{gT}{1-\alpha}
\left(
\sqrt{\left(\frac{gT/2-t}{gT}\right)_+}
-
\frac{g\sqrt{T}}{\sqrt2}
\right)_+^2
\right\},
\]
and optimizing over \(g\) and \(\nu\) gives the worst-case lower bound
\[
\sup_{\nu \in \mathcal{P}(\mathcal{M})}\textup{CVaR}_\alpha^{\nu\otimes P^{M,\textup{Alg}}}(R_T)
\ge
\sup_{g>0}\textup{CVaR}_\alpha^{\mu\otimes P^{M,\textup{Alg}}}(R_T)
\ge
c_\alpha \sqrt{T}.
\]
\end{theorem}

\textit{Proof sketch:}
The proof has the same three-step structure as in the passive example. First, identify the bounded loss level \(L_{\max}=gT\). Second, verify the balanced-pair condition:
\[
R_T(M_1,h)+R_T(M_2,h)=gT
\]
for every transcript \(h\). Third, bound the Hellinger divergence through the Gaussian KL comparison
\[
D_{H^2}(P^{M_1,\textup{Alg}}\|P^{M_2,\textup{Alg}})
\le
D_{\textup{KL}}(P^{M_1,\textup{Alg}}\|P^{M_2,\textup{Alg}})
=
\frac{g^2T}{2},
\]
and apply Corollary~\ref{cor:balanced-two-point-hellinger}. The closed form follows from Lemma~\ref{lem:scalar-min}, and optimizing over \(g\) and subsequently over \(\nu\) yields the stated worst-case \(\sqrt{T}\) lower bound. \qed

\begin{remark}[Interpretation and relation to the examples]
\label{rem:comparison-with-bayes-risk}
Proposition~\ref{prop:cvar-dominates-bayes-risk} shows that any Bayes-risk lower bound immediately yields a \(\textup{CVaR}_\alpha\) lower bound of the same order. Thus, Bayes expected risk provides a natural baseline for interpreting tail-sensitive converse bounds. The value of the generalized-Fano corollaries is that they target the tail directly under the prior-predictive law, provide an instantiation template for both non-interactive and interactive regimes, and can yield more informative statements than expectation alone through explicit dependence on \(\alpha\).

For the Gaussian estimation and two-armed bandit examples, the bounds obtained from the two-point Hellinger template recover the same scaling as the corresponding Bayes-risk benchmarks, namely \(n^{-1/2}\) and \(\sqrt{T}\), respectively. At the same time, unlike the expected Bayes-risk bound, they retain \(\alpha\)-dependence through \(\Psi_\alpha(\cdot)\) and \(c_\alpha\). In this sense, the examples show that the generalized-Fano corollaries are both calibrated to the known Bayesian scaling in canonical two-point problems and capable of capturing tail-sensitive behavior.
\end{remark}
\iffalse
\begin{remark}[Bayesian versus minimax \(\textup{CVaR}\)]
Unlike expected loss, \(\textup{CVaR}_\alpha\) is not linear in the prior, so
\[
\textup{CVaR}_\alpha^{\mu\otimes P^{M,\textup{Alg}}}(L)
\neq
\mathbb{E}_{M\sim\mu}\!\left[\textup{CVaR}_\alpha^{P^{M,\textup{Alg}}}(L)\right]
\]
in general. Therefore, the usual Bayes-risk-to-minimax argument for expected loss does not apply directly to \(\textup{CVaR}\), and the minimax \(\textup{CVaR}\) analysis is left for future work.
\end{remark}
\fi
\section{Conclusion and Future Work}

We instantiate the generalized-Fano Bayesian \(\textup{CVaR}\) framework through a reusable two-point Hellinger template. We derive worst-case Bayesian CVaR lower bounds for canonical problems such as Gaussian mean estimation and two-armed Gaussian bandits showing the applicability of the Hellinger template. The resulting bounds recover the known Bayesian expected-risk scaling while also revealing tail risk level \(\alpha\)-dependence, which is invisible to expected-risk lower bounds. Natural next steps are to understand whether multi-point versions of the present Hellinger template can recover sharp Bayesian \(\textup{CVaR}\) scalings in \(K\)-armed bandits and episodic Markov decision processes while preserving explicit dependence on the risk level \(\alpha\), and to compare those lower bounds with matching upper bounds for risk-sensitive algorithms.

\IEEEpeerreviewmaketitle

\bibliographystyle{IEEEtran}
\bibliography{references}
\newpage
\appendix
\section{Proofs}
\label{app:proofs}

\begin{lemma}[Scalar minimization]
\label{lem:scalar-min}
For \(\alpha\in[0,1)\) and \(\rho\ge 0\), define
\[
F_{\alpha,\rho}(x):=
\frac12-x+\frac{(\sqrt{x}-\rho/\sqrt2)_+^2}{1-\alpha},
\qquad x\in[0,1/2].
\]
Then
\[
\inf_{x\in[0,1/2]}F_{\alpha,\rho}(x)=\Psi_\alpha(\rho).
\]
Moreover,
\[
\sup_{\rho\ge 0}\rho\,\Psi_\alpha(\rho)=c_\alpha.
\]
\end{lemma}

\begin{proof}
Write \(x=s^2\), with \(s\in[0,1/\sqrt2]\). Then
\[
F_{\alpha,\rho}(x)
=
\frac12-s^2+\frac{(s-\rho/\sqrt2)_+^2}{1-\alpha}.
\]

We first identify \(\inf_{x\in[0,1/2]}F_{\alpha,\rho}(x)\).

If \(0\le s\le \rho/\sqrt2\), then
\[
F_{\alpha,\rho}(x)=\frac12-s^2,
\]
so the minimum on this region is
\[
\begin{cases}
\frac12-\frac{\rho^2}{2}, & 0\le \rho\le 1,\\[1ex]
0, & \rho\ge 1.
\end{cases}
\]

Now suppose \(\rho/\sqrt2\le s\le 1/\sqrt2\). Then
\[
F_{\alpha,\rho}(x)
=
\frac12-s^2+\frac{(s-\rho/\sqrt2)^2}{1-\alpha}.
\]

If \(\alpha=0\), this becomes
\[
\frac12-\sqrt2\,\rho\,s+\frac{\rho^2}{2},
\]
which is decreasing in \(s\), hence minimized at \(s=1/\sqrt2\), with value
\[
\frac{(1-\rho)^2}{2}
\qquad (0\le \rho\le 1).
\]
Since
\[
\frac{(1-\rho)^2}{2}\le \frac12-\frac{\rho^2}{2}
\qquad (0\le \rho\le 1),
\]
it follows that
\[
\inf_{x\in[0,1/2]}F_{0,\rho}(x)
=
\begin{cases}
\frac{(1-\rho)^2}{2}, & 0\le \rho\le 1,\\[1ex]
0, & \rho\ge 1,
\end{cases}
=
\Psi_0(\rho).
\]

Assume now \(0<\alpha<1\). On \([\rho/\sqrt2,1/\sqrt2]\),
\[
G_{\alpha,\rho}(s):=
\frac12-s^2+\frac{(s-\rho/\sqrt2)^2}{1-\alpha}
\]
is a convex quadratic with
\[
G'_{\alpha,\rho}(s)=\frac{2\alpha}{1-\alpha}s-\frac{\sqrt2\,\rho}{1-\alpha},
\]
so its critical point is
\[
s^\star=\frac{\rho}{\sqrt2\,\alpha}.
\]

If \(0\le \rho\le \alpha\), then \(s^\star\in[\rho/\sqrt2,1/\sqrt2]\), and
\[
G_{\alpha,\rho}(s^\star)=\frac12-\frac{\rho^2}{2\alpha}.
\]
If \(\alpha<\rho\le 1\), then \(s^\star>1/\sqrt2\), so \(G_{\alpha,\rho}\) is decreasing on \([\rho/\sqrt2,1/\sqrt2]\), and its minimum is attained at \(s=1/\sqrt2\), giving
\[
G_{\alpha,\rho}(1/\sqrt2)=\frac{(1-\rho)^2}{2(1-\alpha)}.
\]
Finally, if \(\rho\ge 1\), choosing \(s=1/\sqrt2\) gives value \(0\).

Comparing with the first region,
\[
\frac12-\frac{\rho^2}{2\alpha}\le \frac12-\frac{\rho^2}{2}
\qquad (0\le \rho\le \alpha),
\]
and
\[
\frac{(1-\rho)^2}{2(1-\alpha)}\le \frac12-\frac{\rho^2}{2}
\qquad (\alpha<\rho\le 1).
\]
Hence
\[
\inf_{x\in[0,1/2]}F_{\alpha,\rho}(x)=\Psi_\alpha(\rho)
\]
for all \(\alpha\in[0,1)\).

It remains to maximize \(\rho\Psi_\alpha(\rho)\). Since \(\Psi_\alpha(\rho)=0\) for \(\rho\ge 1\), it suffices to work on \([0,1]\).

If \(\alpha=0\), then
\[
\rho\Psi_0(\rho)=\frac{\rho(1-\rho)^2}{2},\qquad 0\le \rho\le 1,
\]
whose derivative is
\[
\frac{(1-\rho)(1-3\rho)}{2}.
\]
Thus the maximum is attained at \(\rho=1/3\), with value
\[
\frac{(1/3)(2/3)^2}{2}=\frac{2}{27}=c_0.
\]

Now let \(0<\alpha<1\). We have
\[
\rho\Psi_\alpha(\rho)=
\begin{cases}
\displaystyle
\frac{\rho}{2}-\frac{\rho^3}{2\alpha}, & 0\le \rho\le \alpha,\\[2ex]
\displaystyle
\frac{\rho(1-\rho)^2}{2(1-\alpha)}, & \alpha<\rho\le 1.
\end{cases}
\]

On \([0,\alpha]\), the derivative is
\[
\frac12-\frac{3\rho^2}{2\alpha},
\]
so the critical point is \(\rho=\sqrt{\alpha/3}\). Hence the maximum on \([0,\alpha]\) is
\[
\frac{\sqrt{\alpha}}{3\sqrt3}
\quad\text{if }\alpha\ge \frac13,
\]
and otherwise is attained at \(\rho=\alpha\), with value
\[
\frac{\alpha(1-\alpha)}{2}.
\]

On \((\alpha,1]\), the derivative is
\[
\frac{(1-\rho)(1-3\rho)}{2(1-\alpha)},
\]
so the interior critical point is \(\rho=1/3\). Hence the maximum on \((\alpha,1]\) is
\[
\frac{2}{27(1-\alpha)}
\quad\text{if }\alpha\le \frac13,
\]
and otherwise is attained at the boundary \(\rho=\alpha\), with value
\[
\frac{\alpha(1-\alpha)}{2}.
\]

Since the two branches agree at \(\rho=\alpha\), it remains only to compare the boundary value with the interior one. For \(0\le \alpha\le 1/3\),
\[
\frac{\alpha(1-\alpha)}{2}\le \frac{2}{27(1-\alpha)},
\]
and for \(1/3\le \alpha<1\),
\[
\frac{\alpha(1-\alpha)}{2}\le \frac{\sqrt{\alpha}}{3\sqrt3}.
\]
Therefore
\[
\sup_{\rho\ge 0}\rho\,\Psi_\alpha(\rho)=
\begin{cases}
\displaystyle
\frac{2}{27(1-\alpha)}, & 0\le \alpha\le \frac13,\\[2ex]
\displaystyle
\frac{\sqrt{\alpha}}{3\sqrt3}, & \frac13\le \alpha<1,
\end{cases}
\]
which is exactly \(c_\alpha\).
\end{proof}

\begin{proof}[\textbf{Proof of Lemma~\ref{lem:two-point-hellinger-cvar}}]
Apply Corollary~\ref{cor:cvar-lower-bt} with \(f=H^2\) and reference law
\[
Q:=P_2.
\]
For each \(t\in\mathbb R\), define
\[
b_t:=\frac{1}{L_{\max}}\mathbb E_{\mu,Q}[(L-t)_+].
\]
Since \(\mu=\frac12(\delta_{M_1}+\delta_{M_2})\),
\[
b_t
=
\frac{1}{2L_{\max}}
\mathbb E_{P_2}\big[(L(M_1,X)-t)_+ + (L(M_2,X)-t)_+\big].
\]
Fix \(x\in\mathcal X\). By convexity of \(u\mapsto (u-t)_+\),
\[
\frac12\Big((L(M_1,x)-t)_+ + (L(M_2,x)-t)_+\Big)
\ge\]
\[
\left(\frac{L(M_1,x)+L(M_2,x)}{2}-t\right)_+.
\]
Using assumption \((\mathrm{H1})\),
\[
\frac12\Big((L(M_1,x)-t)_+ + (L(M_2,x)-t)_+\Big)
\ge
\left(\frac{C}{2}-t\right)_+.
\]
Taking \(\mathbb E_{P_2}\) and dividing by \(L_{\max}\) yields
\[
b_t\ge \left(\frac{C/2-t}{L_{\max}}\right)_+.
\tag{A.7}
\]

Next,
\[
B
=
\mathbb E_{M\sim\mu}D_{H^2}(P^{M,\textup{Alg}}\|Q)
=\]
\[
\frac12D_{H^2}(P_1\|P_2)+\frac12D_{H^2}(P_2\|P_2)
=
\frac12D_{H^2}(P_1\|P_2).
\]
Using \((\mathrm{H2})\),
\[
B\le \frac{\Gamma_H}{2}.
\tag{A.8}
\]

Corollary~\ref{cor:cvar-lower-bt} gives
\[
\textup{CVaR}_\alpha^{\mu\otimes P^{M,\textup{Alg}}}(L)
\ge
\min_{t\in\mathbb R}
\left\{
t+\frac{L_{\max}}{1-\alpha}\,a^-_{H^2}(B;b_t)
\right\}.
\tag{A.9}
\]
For Bernoulli laws, one has
\[
D_{H^2}\!\bigl(\textup{Bern}(a)\,\big\|\,\textup{Bern}(b)\bigr)
=
1-\sqrt{ab}-\sqrt{(1-a)(1-b)}
=\]
\[
\frac12\Bigl((\sqrt a-\sqrt b)^2+(\sqrt{1-a}-\sqrt{1-b})^2\Bigr).
\]
Hence
\[
D_{H^2}\!\bigl(\textup{Bern}(a)\,\big\|\,\textup{Bern}(b)\bigr)
\ge \frac12(\sqrt a-\sqrt b)^2.
\]
Therefore, if
\[
D_{H^2}\!\bigl(\textup{Bern}(a)\,\big\|\,\textup{Bern}(b)\bigr)\le B,
\]
then
\[
\sqrt a\ge (\sqrt b-\sqrt{2B})_+,
\]
and so
\[
a^-_{H^2}(B;b)\ge (\sqrt b-\sqrt{2B})_+^2.
\tag{A.10}
\]

Combining \((\mathrm{A.7})\), \((\mathrm{A.8})\), and \((\mathrm{A.10})\), we obtain
\[
a^-_{H^2}(B;b_t)
\ge
\left(
\sqrt{\left(\frac{C/2-t}{L_{\max}}\right)_+}
-\sqrt{\Gamma_H}
\right)_+^2.
\]
Substituting this into \((\mathrm{A.9})\) proves the claim.
\end{proof}

\begin{proof}[Proof of Corollary~\ref{cor:balanced-two-point-hellinger}]
Set \(C=L_{\max}\) in Lemma~\ref{lem:two-point-hellinger-cvar}. This gives
\[
\textup{CVaR}_\alpha^{\mu\otimes P^{M,\textup{Alg}}}(L)
\ge
\]
\[
\min_{t\in\mathbb R}
\left\{
t+\frac{L_{\max}}{1-\alpha}
\left(
\sqrt{\left(\frac12-\frac{t}{L_{\max}}\right)_+}
-\sqrt{\Gamma_H}
\right)_+^2
\right\}.
\]
\end{proof}

\begin{proof}[\textbf{Proof of Proposition~\ref{prop:cvar-dominates-bayes-risk}}]
By the dual representation of Expected Shortfall/\(\textup{CVaR}\) \cite{ang2018dual},
\[
\textup{CVaR}_\alpha(L)
=
\sup\Bigl\{
\mathbb{E}[LZ]:
0\le Z\le \tfrac{1}{1-\alpha},\ \mathbb{E}[Z]=1
\Bigr\}.
\]
Since the choice \(Z\equiv 1\) is feasible, it follows that
\[
\textup{CVaR}_\alpha(L)\ge \mathbb{E}[L].
\]
Applying this under the Bayesian law \(\mu\otimes P^{M,\textup{Alg}}\) gives
\[
\textup{CVaR}_\alpha^{\mu\otimes P^{M,\textup{Alg}}}(L)
\ge
\mathbb{E}_{\mu\otimes P^{M,\textup{Alg}}}[L].
\]
The remaining claims follow immediately.
\end{proof}

\begin{proof}[\textbf{Proof of Theorem~\ref{prop:gaussian-mean-hellinger-cvar}}]
We verify the assumptions of Corollary~\ref{cor:balanced-two-point-hellinger}.

\medskip
\noindent
\textit{Step 1: Boundedness and balanced-pair condition.}
By definition,
\[
0\le L(\theta,X)\le 2\Delta.
\]
Hence \(L_{\max}=2\Delta\). Fix any realized estimate \(z\in\mathbb R\). Then
\[
|z-\Delta|+|z+\Delta|\ge 2\Delta
\]
by the triangle inequality, so
\[
\min\{|z-\Delta|,2\Delta\}+\min\{|z+\Delta|,2\Delta\}\ge 2\Delta.
\]
Applying this with \(z=\widehat\theta\) gives
\[
L(+\Delta,X)+L(-\Delta,X)\ge 2\Delta=L_{\max}.
\]

\medskip
\noindent
\medskip
\noindent
\textit{Step 2: Hellinger budget.}
The observation vector \(Y=(Y_1,\dots,Y_n)\) has laws
\[
\mathcal N(\Delta,1)^{\otimes n}
\qquad\text{and}\qquad
\mathcal N(-\Delta,1)^{\otimes n}
\]
under \(M_+\) and \(M_-\), respectively. Since the transcript
\[
X=(Y,\widehat\theta),\qquad \widehat\theta=\textup{Alg}(Y),
\]
is a measurable function of \(Y\), the data-processing inequality gives
\[
D_{H^2}(P^{M_+,\textup{Alg}}\|P^{M_-,\textup{Alg}})\le
D_{\textup{KL}}\bigl(\mathcal N(\Delta,1)^{\otimes n}\,\big\|\,\mathcal N(-\Delta,1)^{\otimes n}\bigr).
\]
Now
\[
D_{\textup{KL}}\bigl(\mathcal N(\Delta,1)^{\otimes n}\,\big\|\,\mathcal N(-\Delta,1)^{\otimes n}\bigr)
=\]
\[
n\,D_{\textup{KL}}(\mathcal N(\Delta,1)\|\mathcal N(-\Delta,1))
=
n\cdot \frac{(2\Delta)^2}{2}
=
2n\Delta^2.
\]
Hence we may take
\[
\Gamma_H:=2n\Delta^2.
\]

\medskip
\noindent
\text{Step 3: Apply Corollary~\ref{cor:balanced-two-point-hellinger}.}
We obtain
\[
\textup{CVaR}_\alpha^{\mu\otimes P^{M,\textup{Alg}}}(L)
\ge
\]
\[
\min_{t\in\mathbb R}
\left\{
t+\frac{2\Delta}{1-\alpha}
\left(
\sqrt{\left(\frac{\Delta-t}{2\Delta}\right)_+}
-\sqrt{2n}\,\Delta
\right)_+^2
\right\}.
\tag{A.11}
\]

\medskip
\noindent
\text{Step 4: Closed form.}
Since \(0\le L\le 2\Delta\), the minimization in \((\mathrm{A.11})\) may be restricted to \(t\in[0,2\Delta]\). Set
\[
x:=\frac12-\frac{t}{2\Delta}\in[0,1/2].
\]
Then \(t=2\Delta(1/2-x)\), and \((\mathrm{A.11})\) becomes
\[
\textup{CVaR}_\alpha^{\mu\otimes P^{M,\textup{Alg}}}(L)
\ge\]
\[
2\Delta
\inf_{x\in[0,1/2]}
\left\{
\frac12-x+\frac{(\sqrt{x}-\sqrt{2n}\,\Delta)_+^2}{1-\alpha}
\right\}.
\]
Since
\(
\sqrt{2n}\,\Delta=\frac{2\sqrt n\,\Delta}{\sqrt2},
\)
Lemma~\ref{lem:scalar-min} with \(\rho=2\sqrt n\,\Delta\) yields
\[
\textup{CVaR}_\alpha^{\mu\otimes P^{M,\textup{Alg}}}(L)
\ge
2\Delta\,\Psi_\alpha(2\sqrt n\,\Delta).
\]

Finally, write \(\rho:=2\sqrt n\,\Delta\), so that \(2\Delta=\rho/\sqrt n\). Then
\[
2\Delta\,\Psi_\alpha(2\sqrt n\,\Delta)=\frac{\rho\,\Psi_\alpha(\rho)}{\sqrt n}.
\]
Maximizing over \(\rho\ge 0\) and using Lemma~\ref{lem:scalar-min} gives
\[
\sup_{\Delta>0}
\textup{CVaR}_\alpha^{\mu\otimes P^{M,\textup{Alg}}}(L)
\ge
\frac{c_\alpha}{\sqrt n}.
\]
\end{proof}

\begin{proof}[\textbf{Proof of Theorem~\ref{thm:two-armed-gaussian-hellinger-cvar}}]
We verify the assumptions of Corollary~\ref{cor:balanced-two-point-hellinger}.

\medskip
\noindent
\text{Step 1: Balanced-pair condition and boundedness.}
Under \(M_1\), arm \(1\) is optimal and arm \(2\) is suboptimal by gap \(g\), so
\[
R_T(M_1,H_T)=g\,N_2(H_T).
\]
Under \(M_2\), arm \(2\) is optimal and arm \(1\) is suboptimal by gap \(g\), so
\[
R_T(M_2,H_T)=g\,N_1(H_T).
\]
Hence, for every transcript \(h\),
\[
R_T(M_1,h)+R_T(M_2,h)
=
g\bigl(N_1(h)+N_2(h)\bigr)
=
gT.
\]
Also,
\[
0\le R_T(M_i,H_T)\le gT,\qquad i=1,2.
\]
Thus \(L(M,H_T):=R_T(M,H_T)\) satisfies the balanced-pair condition with \(L_{\max}=gT\).

\medskip
\noindent
\text{Step 2: Hellinger budget.}
For this two-armed Gaussian construction,
\[
D_{H^2}(P^{M_1,\textup{Alg}}\|P^{M_2,\textup{Alg}})
\le
D_{\textup{KL}}(P^{M_1,\textup{Alg}}\|P^{M_2,\textup{Alg}})
=
\frac{g^2T}{2}.
\]
Hence we may take
\(
\Gamma_H:=\frac{g^2T}{2}.
\)

\medskip
\noindent
\text{Step 3: Apply Corollary~\ref{cor:balanced-two-point-hellinger}.}
We obtain
\[
\textup{CVaR}_\alpha^{\mu\otimes P^{M,\textup{Alg}}}(R_T)
\ge\]
\[
\min_{t\in\mathbb R}
\left\{
t+\frac{gT}{1-\alpha}
\left(
\sqrt{\left(\frac{gT/2-t}{gT}\right)_+}
-
\frac{g\sqrt{T}}{\sqrt2}
\right)_+^2
\right\}.
\tag{A.12}
\]

\medskip
\noindent
\text{Step 4: Closed form.}
Since \(0\le R_T\le gT\), the minimization in \((\mathrm{A.12})\) may be restricted to \(t\in[0,gT]\). Set
\[
x:=\frac12-\frac{t}{gT}\in[0,1/2].
\]
Then \(t=gT(1/2-x)\), and \((\mathrm{A.12})\) becomes
\[
\textup{CVaR}_\alpha^{\mu\otimes P^{M,\textup{Alg}}}(R_T)
\ge\]
\[
gT
\inf_{x\in[0,1/2]}
\left\{
\frac12-x+\frac{(\sqrt{x}-g\sqrt{T}/\sqrt2)_+^2}{1-\alpha}
\right\}.
\]
Lemma~\ref{lem:scalar-min} with \(\rho=g\sqrt T\) yields
\[
\textup{CVaR}_\alpha^{\mu\otimes P^{M,\textup{Alg}}}(R_T)
\ge
gT\,\Psi_\alpha(g\sqrt T).
\]

Finally, write \(\rho:=g\sqrt T\), so that \(gT=\rho\sqrt T\). Then
\[
gT\,\Psi_\alpha(g\sqrt T)=\rho\,\Psi_\alpha(\rho)\,\sqrt T.
\]
Maximizing over \(\rho\ge 0\) and using Lemma~\ref{lem:scalar-min} gives
\[
\sup_{g>0}
\textup{CVaR}_\alpha^{\mu\otimes P^{M,\textup{Alg}}}(R_T)
\ge
c_\alpha\sqrt T.
\]
\end{proof}



\iffalse
\section{New Results}

This section states the main tools and examples used in the paper. Full proofs are deferred to the appendix.

We first record a reusable two-point instantiation of Corollaries~\ref{cor:hinge-lower} and~\ref{cor:cvar-lower-bt}. It packages the ingredients that recur in the examples below.

\begin{lemma}[Two-point Bayesian \textup{CVaR} lower bound]
\label{lem:two-point-cvar}
Let \(\mathcal{M}=\{M_1,M_2\}\), let
\[
\mu=\frac12\delta_{M_1}+\frac12\delta_{M_2},
\]
and fix an algorithm \(\textup{Alg}\). Let
\[
P_i:=P^{M_i,\textup{Alg}},\qquad i=1,2.
\]
Let \(L:\mathcal{M}\times \mathcal{X}\to [0,L_{\max}]\) be a bounded loss, where \(\mathcal{X}\) denotes the transcript space. Assume that, for some \(C\in[0,2L_{\max}]\),
\[
L(M_1,x)+L(M_2,x)\ge C
\qquad
\text{for all }x\in\mathcal{X},
\tag{A1}
\]
and that
\[
D_{\textup{KL}}(P_1\|P_2)\le \Gamma.
\tag{A2}
\]
Then, for every \(\alpha\in[0,1)\),
\[
\textup{CVaR}_\alpha^{\mu\otimes P^{M,\textup{Alg}}}(L)
\ge\]
\[
\min_{t\in\mathbb{R}}
\left\{
t+\frac{L_{\max}}{1-\alpha}
\left[
\left(\frac{C/2-t}{L_{\max}}\right)_+
-\frac{\sqrt{\Gamma}}{2}
\right]_+
\right\}.
\]
\end{lemma}

\textit{Proof sketch:}
Apply Corollary~\ref{cor:hinge-lower} with \(f=\textup{KL}\) and reference law \(Q=P_2\). The pairwise separation condition \((\mathrm{A1})\) yields a lower bound on the reference hinge term \(b_t\), while \((\mathrm{A2})\) controls the divergence budget \(B\). A Bernoulli Pinsker inversion then lower-bounds \(a_{\textup{KL}}^{-}(B;b_t)\), and Corollary~\ref{cor:cvar-lower-bt} converts the resulting hinge bound into a \(\textup{CVaR}\) lower bound. \qed

\begin{corollary}[Balanced-pair ]
\label{cor:balanced-two-point-cvar}
Under the assumptions of Lemma~\ref{lem:two-point-cvar}, if in addition
\[
L(M_1,x)+L(M_2,x)\ge L_{\max}
\qquad
\text{for all }x\in\mathcal{X},
\tag{A3}
\]
then
\[
\textup{CVaR}_\alpha^{\mu\otimes P^{M,\textup{Alg}}}(L)
\ge
L_{\max}\left[\frac12-\frac{\sqrt{\Gamma}}{2}\right]_+.
\]
\end{corollary}

\textit{Proof sketch:}
Substitute \(C=L_{\max}\) in Lemma~\ref{lem:two-point-cvar} and explicitly minimize the resulting one-dimensional lower bound over \(t\). The minimizer is obtained by a simple three-region analysis. \qed

\begin{proposition}[Bayesian \textup{CVaR} dominates Bayesian expected loss]
\label{prop:cvar-dominates-bayes-risk}
Let \(L\) be any integrable loss, and let \(\alpha\in[0,1)\). Then
\[
\textup{CVaR}_\alpha(L)\ge \mathbb{E}[L].
\]
Consequently, for any prior \(\mu\),
\[
\textup{CVaR}_\alpha^{\mu\otimes P^{M,\textup{Alg}}}(L)
\ge
\mathbb{E}_{\mu\otimes P^{M,\textup{Alg}}}[L].
\]
In particular, if for some prior \(\mu\) one has
\[
\mathbb{E}_{\mu\otimes P^{M,\textup{Alg}}}[L]\ge c\,\psi,
\]
then
\[
\textup{CVaR}_\alpha^{\mu\otimes P^{M,\textup{Alg}}}(L)\ge c\,\psi.
\]
Moreover,
\[
\sup_{\mu}\textup{CVaR}_\alpha^{\mu\otimes P^{M,\textup{Alg}}}(L)
\ge
\sup_{\mu}\mathbb{E}_{\mu\otimes P^{M,\textup{Alg}}}[L].
\]
\end{proposition}

\textit{Proof sketch:}
Use the dual representation of Expected Shortfall/\(\textup{CVaR}\):
\[
\textup{CVaR}_\alpha(L)
=
\sup\Bigl\{
\mathbb{E}[LZ]:
0\le Z\le \tfrac{1}{1-\alpha},\ \mathbb{E}[Z]=1
\Bigr\}.
\]
Choosing \(Z= 1\) gives \(\textup{CVaR}_\alpha(L)\ge \mathbb{E}[L]\), and the Bayesian statements follow immediately. \qed

\section{Examples}

\subsection{A passive estimation example}

We illustrate the reusable two-point Bayesian \textup{CVaR} lemma in a passive estimation problem.

\textit{Setup:}
Let \(\theta\in\{-\Delta,+\Delta\}\), where \(\Delta>0\), and suppose we observe
\[
Y_1,\dots,Y_n \;\;\text{i.i.d.}\;\; \mathcal{N}(\theta,1).
\]
An estimator \(\widehat{\theta}=\widehat{\theta}(Y_1,\dots,Y_n)\) is any measurable function of the sample. This fits the ISDM framework in the passive case:
\begin{itemize}
\item the model is \(M_\theta\), indexed by \(\theta\in\{-\Delta,+\Delta\}\);
\item there is no interaction, so the algorithm \(\textup{Alg}\) only outputs the final decision \(\widehat{\theta}\);
\item the transcript is
\[
X=(Y_1,\dots,Y_n).
\]
\end{itemize}

We equip this problem with the bounded estimation loss
\[
L(\theta,X):=\min\{|\widehat{\theta}(X)-\theta|,\;2\Delta\}.
\]
Then \(0\le L\le 2\Delta\). Let
\[
M_+ := M_{+\Delta},\qquad M_-:=M_{-\Delta},
\qquad
\mu:=\frac12\delta_{M_+}+\frac12\delta_{M_-}.
\]

\begin{proposition}[Bayesian \textup{CVaR} lower bound for Gaussian mean estimation]
\label{prop:gaussian-mean-cvar}
For every estimator \(\widehat{\theta}\) and every \(\alpha\in[0,1)\),
\[
\textup{CVaR}_\alpha^{\mu\otimes P^{M,\textup{Alg}}}(L)
\ge
2\Delta\left[\frac12-\Delta\sqrt{\frac{n}{2}}\right]_+.
\]
In particular, choosing
\[
\Delta=\frac{1}{2\sqrt{2n}}
\]
gives
\[
\textup{CVaR}_\alpha^{\mu\otimes P^{M,\textup{Alg}}}(L)
\ge
\frac{1}{4\sqrt{2n}}.
\]
\end{proposition}

\paragraph{Proof sketch.}
The loss is bounded by \(L_{\max}=2\Delta\). The triangle inequality gives the balanced-pair condition
\[
L(+\Delta,X)+L(-\Delta,X)\ge 2\Delta.
\]
The KL divergence between the two Gaussian product measures is
\[
D_{\textup{KL}}(P_+\|P_-)=2n\Delta^2.
\]
Corollary~\ref{cor:balanced-two-point-cvar} then yields the stated bound, and optimizing over \(\Delta\) gives the \(n^{-1/2}\) scaling. \qed

\subsection{Two-armed Gaussian bandits}

We now instantiate the same reusable two-point lemma for the two-armed Gaussian bandit.

Let \(T\ge 1\). Consider the two-armed Gaussian bandit with unit-variance rewards and horizon \(T\). For a parameter \(g>0\), define the two environments
\[
M_1:=\left(+\frac{g}{2},-\frac{g}{2}\right),
\qquad
M_2:=\left(-\frac{g}{2},+\frac{g}{2}\right).
\]
At each round \(t\), the algorithm \(\textup{Alg}\) chooses an arm \(A_t\in\{1,2\}\) based on the past history, then observes the reward
\[
Y_t=\mu_{A_t}+Z_t,
\qquad
Z_t\sim \mathcal{N}(0,1).
\]
Thus the transcript is
\[
H_T=(A_1,Y_1,\dots,A_T,Y_T),
\]
which fits the ISDM template with model \(M\in\{M_1,M_2\}\), decision \(A_t\), and observation \(Y_t\).

Let
\[
\mu=\frac12\delta_{M_1}+\frac12\delta_{M_2}.
\]
For \(a\in\{1,2\}\), let \(N_a(H_T)\) denote the number of pulls of arm \(a\) up to time \(T\).

\begin{theorem}[Bayesian \textup{CVaR} lower bound for two-armed Gaussian bandits]
\label{thm:two-armed-gaussian-cvar-reusable}
For every algorithm \(\textup{Alg}\) and every \(\alpha\in[0,1)\),
\[
\textup{CVaR}_\alpha^{\mu\otimes P^{M,\textup{Alg}}}(R_T)
\ge
gT\left[\frac12-\frac{g\sqrt{T}}{2\sqrt{2}}\right]_+.
\]
In particular, choosing
\[
g=\frac{1}{\sqrt{2T}}
\]
gives
\[
\textup{CVaR}_\alpha^{\mu\otimes P^{M,\textup{Alg}}}(R_T)
\ge
\frac{\sqrt{T}}{4\sqrt{2}}.
\]
\end{theorem}

\paragraph{Proof sketch.}
Under \(M_1\), \(R_T=gN_2\), and under \(M_2\), \(R_T=gN_1\). Hence
\[
R_T(M_1,h)+R_T(M_2,h)=gT
\]
for every transcript \(h\), so the balanced-pair condition holds with \(L_{\max}=gT\). The transcript KL divergence is
\[
D_{\textup{KL}}(P^{M_1,\textup{Alg}}\|P^{M_2,\textup{Alg}})=\frac{g^2T}{2}.
\]
Applying Corollary~\ref{cor:balanced-two-point-cvar} yields the first bound, and optimizing over \(g\) gives the \(\sqrt{T}\) scaling. \qed

\begin{remark}[Interpretation and relation to the examples]
\label{rem:comparison-with-bayes-risk}
Proposition~\ref{prop:cvar-dominates-bayes-risk} shows that any Bayesian expected-loss lower bound immediately yields a Bayesian \(\textup{CVaR}_\alpha\) lower bound of the same order. Thus, Bayesian expected risk provides a natural baseline benchmark for Bayesian \(\textup{CVaR}\). The value of the generalized-Fano \(\textup{CVaR}\) approach is that it targets tail risk directly under the prior-predictive law, provides a  instantiation template that applies in both passive and interactive settings, and can in principle yield more informative lower bounds than expectation alone, for example through explicit dependence on the risk level \(\alpha\).

For the passive Gaussian mean estimation example and the two-armed Gaussian bandit example, the explicit \(\textup{CVaR}_\alpha\) lower bounds obtained from the two-point template recover the same order as the corresponding Bayesian expected-loss benchmarks, namely \(n^{-1/2}\) in the estimation problem and \(\sqrt{T}\) in the two-armed bandit problem. In this sense, the examples verify that the corollary-based \(\textup{CVaR}\) approach is properly calibrated to the known Bayesian risk scaling in these canonical two-point settings, while also illustrating a reusable tail-sensitive lower-bound method that can be applied beyond expected-risk criteria.
\end{remark}
\section{Conclusion and Future Work}
This paper studies explicit Bayesian \(\textup{CVaR}\) lower bounds in interactive statistical decision making by instantiating the generalized-Fano framework of \cite{bongole2026generalizing}. Rather than extending the abstract theory itself, the focus is on making it concrete and usable through a reusable two-point Bayesian \(\textup{CVaR}\) lemma and two worked examples: passive Gaussian mean estimation and the two-armed Gaussian bandit. These examples show that the same generalized-Fano mechanism applies in both passive and interactive settings and recovers the same order as the corresponding Bayesian expected-loss benchmarks in these canonical two-point problems. Future work is to move beyond two-point constructions and obtain explicit Bayesian \(\textup{CVaR}\) lower bounds with sharper scaling in richer problems such as \(K\)-armed bandits and episodic Markov decision processes, to develop reusable multi-model instantiation tools in the spirit of interactive Le Cam-, Assouad-, and Fano-type arguments, and to compare the resulting lower bounds with risk-sensitive algorithmic upper bounds to understand when the induced Bayesian \(\textup{CVaR}\) scalings are tight.
\IEEEpeerreviewmaketitle

\bibliographystyle{IEEEtran}

\bibliography{references}

\newpage
\appendix
\section{Proofs of results}
\label{sec:proofs-results}

\begin{proof}[Proof of Lemma~\ref{lem:two-point-cvar}]
We apply Corollary~\ref{cor:hinge-lower} with \(f=\textup{KL}\) and reference law
\[
Q:=P_2.
\]
For each \(t\in\mathbb{R}\), define
\[
b_t:=\frac{1}{L_{\max}}\,
\mathbb{E}_{\mu,Q}\big[(L-t)_+\big].
\]
Since \(\mu=\frac12(\delta_{M_1}+\delta_{M_2})\) and \(Q=P_2\),
\[
b_t
=
\frac{1}{2L_{\max}}
\mathbb{E}_{P_2}\big[(L(M_1,X)-t)_+ + (L(M_2,X)-t)_+\big].
\]
Fix \(x\in\mathcal{X}\). By convexity of \(u\mapsto (u-t)_+\),
\[
\frac12\Big((L(M_1,x)-t)_+ + (L(M_2,x)-t)_+\Big)
\]
\[
\ge
\left(\frac{L(M_1,x)+L(M_2,x)}{2}-t\right)_+.
\]
Using assumption \((\mathrm{A1})\), we obtain
\[
\frac12\Big((L(M_1,x)-t)_+ + (L(M_2,x)-t)_+\Big)
\ge
\left(\frac{C}{2}-t\right)_+.
\]
Taking \(\mathbb{E}_{P_2}\) and dividing by \(L_{\max}\) yields
\[
b_t\ge \left(\frac{C/2-t}{L_{\max}}\right)_+.
\tag{1}
\]

Next, by the definition of \(B\) in Corollary~\ref{cor:hinge-lower},
\[
B
=
\mathbb{E}_{M\sim\mu}D_{\textup{KL}}(P^{M,\textup{Alg}}\|Q)
\]
\[=
\frac12D_{\textup{KL}}(P_1\|P_2)+\frac12D_{\textup{KL}}(P_2\|P_2)
=
\frac12D_{\textup{KL}}(P_1\|P_2).
\]
Using \((\mathrm{A2})\),
\[
B\le \frac{\Gamma}{2}.
\tag{2}
\]

Now Corollary~\ref{cor:hinge-lower} gives
\[
\mathbb{E}_{\mu,P^{M,\textup{Alg}}}[(L-t)_+]
\ge
L_{\max}\,a^{-}_{\textup{KL}}(B;b_t),
\tag{3}
\]
where
\[
a^{-}_{\textup{KL}}(B;b)
:=
\inf\Bigl\{a\in[0,1]:
D_{\textup{KL}}(\textup{Bern}(a)\|\textup{Bern}(b))\le B
\Bigr\}.
\]
By Pinsker's inequality,
\[
D_{\textup{KL}}(\textup{Bern}(a)\|\textup{Bern}(b))\le B
\quad\Longrightarrow\quad
|a-b|\le \sqrt{\frac{B}{2}}.
\]
Hence
\[
a^{-}_{\textup{KL}}(B;b)\ge \left[b-\sqrt{\frac{B}{2}}\right]_+.
\tag{4}
\]
Combining \((1)\), \((2)\), and \((4)\), we obtain
\[
a^{-}_{\textup{KL}}(B;b_t)
\ge
\left[
\left(\frac{C/2-t}{L_{\max}}\right)_+
-\frac{\sqrt{\Gamma}}{2}
\right]_+.
\tag{5}
\]
Substituting \((5)\) into \((3)\) gives
\[
\mathbb{E}_{\mu,P^{M,\textup{Alg}}}[(L-t)_+]
\ge
L_{\max}
\left[
\left(\frac{C/2-t}{L_{\max}}\right)_+
-\frac{\sqrt{\Gamma}}{2}
\right]_+.
\]
Finally, Corollary~\ref{cor:cvar-lower-bt} implies
\[
\textup{CVaR}_\alpha^{\mu\otimes P^{M,\textup{Alg}}}(L)\ge\]
\[
\min_{t\in\mathbb{R}}
\left\{
t+\frac{L_{\max}}{1-\alpha}
\left[
\left(\frac{C/2-t}{L_{\max}}\right)_+
-\frac{\sqrt{\Gamma}}{2}
\right]_+
\right\}.
\]
This proves the claim.
\end{proof}

\begin{proof}[Proof of Corollary~\ref{cor:balanced-two-point-cvar}]
By Lemma~\ref{lem:two-point-cvar}, it suffices to minimize
\[
F(t):=
t+\frac{L_{\max}}{1-\alpha}
\left[
\left(\frac12-\frac{t}{L_{\max}}\right)_+
-\frac{\sqrt{\Gamma}}{2}
\right]_+.
\]
Set
\[
c:=\frac{\sqrt{\Gamma}}{2}.
\]
We analyze \(F\) in three regions.

\smallskip
\noindent
\emph{Region 1: \(t\ge L_{\max}/2\).}
Then \(\left(\frac12-\frac{t}{L_{\max}}\right)_+=0\), so \(F(t)=t\), and hence
\[
\inf_{t\ge L_{\max}/2}F(t)=\frac{L_{\max}}{2}.
\]

\smallskip
\noindent
\emph{Region 2: \(L_{\max}(1/2-c)\le t\le L_{\max}/2\).}
Then
\[
\left(\frac12-\frac{t}{L_{\max}}\right)_+=\frac12-\frac{t}{L_{\max}},
\qquad
\frac12-\frac{t}{L_{\max}}-c\le 0,
\]
so again \(F(t)=t\), and therefore
\[
\inf_{L_{\max}(1/2-c)\le t\le L_{\max}/2}F(t)
=
L_{\max}\left(\frac12-c\right).
\]

\smallskip
\noindent
\emph{Region 3: \(t<L_{\max}(1/2-c)\).}
Then both positive parts are active, so
\[
F(t)
=
t+\frac{L_{\max}}{1-\alpha}\left(\frac12-\frac{t}{L_{\max}}-c\right)
=
t+\frac{L_{\max}(1/2-c)-t}{1-\alpha}.
\]
This is affine in \(t\), with derivative
\[
F'(t)=1-\frac{1}{1-\alpha}=-\frac{\alpha}{1-\alpha}<0.
\]
Hence \(F\) is decreasing on Region 3, and its infimum is attained at the right endpoint:
\[
\inf_{t<L_{\max}(1/2-c)}F(t)
=
L_{\max}\left(\frac12-c\right).
\]

Combining the three regions,
\[
\min_{t\in\mathbb{R}}F(t)
=
L_{\max}\left(\frac12-c\right)
=
L_{\max}\left(\frac12-\frac{\sqrt{\Gamma}}{2}\right),
\]
provided the quantity in parentheses is nonnegative. If it is negative, the lower bound is replaced by \(0\), since \(F(t)\ge 0\) for all \(t\). Therefore
\[
\min_{t\in\mathbb{R}}F(t)
=
L_{\max}\left[\frac12-\frac{\sqrt{\Gamma}}{2}\right]_+.
\]
This proves the corollary.
\end{proof}

\begin{proof}[Proof of Proposition~\ref{prop:cvar-dominates-bayes-risk}]
By the dual representation of Expected Shortfall/\(\textup{CVaR}\),
\[
\textup{CVaR}_\alpha(L)
=
\sup\Bigl\{
\mathbb{E}[LZ]:
0\le Z\le \tfrac{1}{1-\alpha},\ \mathbb{E}[Z]=1
\Bigr\}.
\]
Since the choice \(Z\equiv 1\) is feasible, it follows that
\[
\textup{CVaR}_\alpha(L)\ge \mathbb{E}[L].
\]
Applying this under the Bayesian law \(\mu\otimes P^{M,\textup{Alg}}\) gives
\[
\textup{CVaR}_\alpha^{\mu\otimes P^{M,\textup{Alg}}}(L)
\ge
\mathbb{E}_{\mu\otimes P^{M,\textup{Alg}}}[L].
\]
The remaining claims follow immediately.
\end{proof}

\begin{proof}[Proof of Proposition~\ref{prop:gaussian-mean-cvar}]
We verify the assumptions of Corollary~\ref{cor:balanced-two-point-cvar}.

\medskip
\noindent
\text{Step 1: Boundedness.}
By definition,
\[
0\le L(\theta,X)\le 2\Delta.
\]
Hence we may take
\[
L_{\max}:=2\Delta.
\]

\medskip
\noindent
\text{Step 2: Balanced-pair condition.}
Fix any realized estimate \(z\in\mathbb{R}\). Define
\[
a:=|z-\Delta|,
\qquad
b:=|z+\Delta|.
\]
By the triangle inequality,
\[
a+b = |z-\Delta|+|z+\Delta| \ge |(\,z+\Delta\,)-(\,z-\Delta\,)|=2\Delta.
\]
Therefore
\[
\min\{a,2\Delta\}+\min\{b,2\Delta\}\ge 2\Delta.
\]
Applying this with \(z=\widehat{\theta}(X)\), we obtain, for every transcript \(X\),
\[
L(+\Delta,X)+L(-\Delta,X)\ge 2\Delta.
\]
Thus the balanced-pair condition \((\mathrm{A3})\) of Corollary~\ref{cor:balanced-two-point-cvar} holds.

\medskip
\noindent
\text{Step 3: KL divergence between the two models.}
Under \(M_+\), the sample law is
\[
P_+ = \mathcal{N}(\Delta,1)^{\otimes n},
\]
and under \(M_-\), the sample law is
\[
P_- = \mathcal{N}(-\Delta,1)^{\otimes n}.
\]
Since the observations are i.i.d., the KL divergence adds:
\[
D_{\textup{KL}}(P_+\|P_-)
=
n\,D_{\textup{KL}}\bigl(\mathcal{N}(\Delta,1)\,\|\,\mathcal{N}(-\Delta,1)\bigr).
\]
For unit-variance Gaussians,
\[
D_{\textup{KL}}\bigl(\mathcal{N}(\mu_1,1)\,\|\,\mathcal{N}(\mu_2,1)\bigr)
=
\frac{(\mu_1-\mu_2)^2}{2}.
\]
Hence
\[
D_{\textup{KL}}(P_+\|P_-)
=
n\cdot \frac{(2\Delta)^2}{2}
=
2n\Delta^2.
\]
So we may take
\[
\Gamma:=2n\Delta^2.
\]

\medskip
\noindent
\text{Step 4: Apply the reusable two-point bound.}
Corollary~\ref{cor:balanced-two-point-cvar} now gives
\[
\textup{CVaR}_\alpha^{\mu\otimes P^{M,\textup{Alg}}}(L)
\ge
L_{\max}\left[\frac12-\frac{\sqrt{\Gamma}}{2}\right]_+.
\]
Substituting \(L_{\max}=2\Delta\) and \(\Gamma=2n\Delta^2\),
\[
\textup{CVaR}_\alpha^{\mu\otimes P^{M,\textup{Alg}}}(L)
\ge
2\Delta\left[\frac12-\frac{\sqrt{2n\Delta^2}}{2}\right]_+
=
2\Delta\left[\frac12-\Delta\sqrt{\frac{n}{2}}\right]_+.
\]
This proves the first claim.

\medskip
\noindent
\text{Step 5: Optimize the separation \(\Delta\).}
Consider
\[
\psi(\Delta):=
2\Delta\left(\frac12-\Delta\sqrt{\frac{n}{2}}\right)
=
\Delta-\sqrt{2n}\,\Delta^2.
\]
For \(\Delta\le 1/\sqrt{2n}\), the positive part is inactive, so maximizing the lower bound reduces to maximizing \(\psi(\Delta)\). Differentiating,
\[
\psi'(\Delta)=1-2\sqrt{2n}\,\Delta.
\]
Setting \(\psi'(\Delta)=0\) yields
\[
\Delta=\frac{1}{2\sqrt{2n}}.
\]
At this value,
\[
2\Delta=\frac{1}{\sqrt{2n}},
\qquad
\Delta\sqrt{\frac{n}{2}}=\frac14.
\]
Therefore
\[
\textup{CVaR}_\alpha^{\mu\otimes P^{M,\textup{Alg}}}(L)
\ge
\frac{1}{\sqrt{2n}}\left(\frac12-\frac14\right)
=
\frac{1}{4\sqrt{2n}}.
\]
This completes the proof.
\end{proof}

\begin{proof}[Proof of Theorem~\ref{thm:two-armed-gaussian-cvar-reusable}]
We verify the assumptions of Corollary~\ref{cor:balanced-two-point-cvar}.

Under \(M_1\), arm \(1\) is optimal and arm \(2\) is suboptimal by gap \(g\), so
\[
R_T(M_1,H_T)=g\,N_2(H_T).
\]
Under \(M_2\), arm \(2\) is optimal and arm \(1\) is suboptimal by gap \(g\), so
\[
R_T(M_2,H_T)=g\,N_1(H_T).
\]
Hence, for every transcript \(h\),
\[
R_T(M_1,h)+R_T(M_2,h)
=
g\bigl(N_1(h)+N_2(h)\bigr)
=
gT.
\]
Also,
\[
0\le R_T(M_i,H_T)\le gT,
\qquad i=1,2.
\]
Thus we may take
\[
L(M,H_T):=R_T(M,H_T),
\qquad
L_{\max}:=gT.
\]
The balanced-pair condition \((\mathrm{A3})\) holds with this \(L_{\max}\).

Finally, for this two-armed Gaussian construction, the transcript KL divergence satisfies
\[
D_{\textup{KL}}(P^{M_1,\textup{Alg}}\|P^{M_2,\textup{Alg}})
=
\frac{g^2T}{2}.
\]
Hence we may take
\[
\Gamma:=\frac{g^2T}{2}.
\]

Applying Corollary~\ref{cor:balanced-two-point-cvar}, we obtain
\[
\textup{CVaR}_\alpha^{\mu\otimes P^{M,\textup{Alg}}}(R_T)
\ge
gT\left[\frac12-\frac{1}{2}\sqrt{\frac{g^2T}{2}}\right]_+
=
gT\left[\frac12-\frac{g\sqrt{T}}{2\sqrt{2}}\right]_+.
\]
This proves the first claim.

To optimize over \(g\), define
\[
\psi(g):=
gT\left(\frac12-\frac{g\sqrt{T}}{2\sqrt{2}}\right).
\]
For \(g\le \sqrt{2/T}\), the positive part is inactive, so maximizing the lower bound reduces to maximizing \(\psi(g)\). Differentiating,
\[
\psi'(g)=\frac{T}{2}-\frac{gT^{3/2}}{\sqrt{2}}.
\]
Setting \(\psi'(g)=0\) yields
\[
g=\frac{1}{\sqrt{2T}}.
\]
At this value,
\[
gT=\frac{\sqrt{T}}{\sqrt{2}},
\qquad
\frac{g\sqrt{T}}{2\sqrt{2}}=\frac14.
\]
and therefore
\[
\textup{CVaR}_\alpha^{\mu\otimes P^{M,\textup{Alg}}}(R_T)
\ge
\frac{\sqrt{T}}{\sqrt{2}}\left(\frac12-\frac14\right)
=
\frac{\sqrt{T}}{4\sqrt{2}}.
\]
This completes the proof.
\end{proof}
\fi

\iffalse
We first record a reusable two-point instantiation of Corollaries~\ref{cor:hinge-lower} and~\ref{cor:cvar-lower-bt}. It packages the ingredients that recur in the examples below.

\begin{lemma}[Two-point Bayesian \textup{CVaR} lower bound]
\label{lem:two-point-cvar}
Let \(\mathcal{M}=\{M_1,M_2\}\), let
\[
\mu=\frac12\delta_{M_1}+\frac12\delta_{M_2},
\]
and fix an algorithm \(\textup{Alg}\). Let
\[
P_i:=P^{M_i,\textup{Alg}},\qquad i=1,2.
\]
Let \(L:\mathcal{M}\times \mathcal{X}\to [0,L_{\max}]\) be a bounded loss, where \(\mathcal{X}\) denotes the transcript space. Assume that, for some \(C\in[0,2L_{\max}]\),
\[
L(M_1,x)+L(M_2,x)\ge C
\qquad
\text{for all }x\in\mathcal{X},
\tag{A1}
\]
and that
\[
D_{\textup{KL}}(P_1\|P_2)\le \Gamma.
\tag{A2}
\]
Then, for every \(\alpha\in[0,1)\),
\[
\textup{CVaR}_\alpha^{\mu\otimes P^{M,\textup{Alg}}}(L)
\ge
\min_{t\in\mathbb{R}}
\left\{
t+\frac{L_{\max}}{1-\alpha}
\left[
\left(\frac{C/2-t}{L_{\max}}\right)_+
-\frac{\sqrt{\Gamma}}{2}
\right]_+
\right\}.
\]
\end{lemma}

\begin{proof}
We apply Corollary~\ref{cor:hinge-lower} with \(f=\textup{KL}\) and reference law
\[
Q:=P_2.
\]
For each \(t\in\mathbb{R}\), define
\[
b_t:=\frac{1}{L_{\max}}\,
\mathbb{E}_{\mu,Q}\big[(L-t)_+\big].
\]
Since \(\mu=\frac12(\delta_{M_1}+\delta_{M_2})\) and \(Q=P_2\),
\[
b_t
=
\frac{1}{2L_{\max}}
\mathbb{E}_{P_2}\big[(L(M_1,X)-t)_+ + (L(M_2,X)-t)_+\big].
\]
Fix \(x\in\mathcal{X}\). By convexity of \(u\mapsto (u-t)_+\),
\[
\frac12\Big((L(M_1,x)-t)_+ + (L(M_2,x)-t)_+\Big)
\ge
\left(\frac{L(M_1,x)+L(M_2,x)}{2}-t\right)_+.
\]
Using assumption \((\mathrm{A1})\), we obtain
\[
\frac12\Big((L(M_1,x)-t)_+ + (L(M_2,x)-t)_+\Big)
\ge
\left(\frac{C}{2}-t\right)_+.
\]
Taking \(\mathbb{E}_{P_2}\) and dividing by \(L_{\max}\) yields
\[
b_t\ge \left(\frac{C/2-t}{L_{\max}}\right)_+.
\tag{1}
\]

Next, by the definition of \(B\) in Corollary~\ref{cor:hinge-lower},
\[
B
=
\mathbb{E}_{M\sim\mu}D_{\textup{KL}}(P^{M,\textup{Alg}}\|Q)
=
\frac12D_{\textup{KL}}(P_1\|P_2)+\frac12D_{\textup{KL}}(P_2\|P_2)
=
\frac12D_{\textup{KL}}(P_1\|P_2).
\]
Using \((\mathrm{A2})\),
\[
B\le \frac{\Gamma}{2}.
\tag{2}
\]

Now Corollary~\ref{cor:hinge-lower} gives
\[
\mathbb{E}_{\mu,P^{M,\textup{Alg}}}[(L-t)_+]
\ge
L_{\max}\,a^{-}_{\textup{KL}}(B;b_t),
\tag{3}
\]
where
\[
a^{-}_{\textup{KL}}(B;b)
:=
\inf\Bigl\{a\in[0,1]:
D_{\textup{KL}}(\textup{Bern}(a)\|\textup{Bern}(b))\le B
\Bigr\}.
\]
By Pinsker's inequality,
\[
D_{\textup{KL}}(\textup{Bern}(a)\|\textup{Bern}(b))\le B
\quad\Longrightarrow\quad
|a-b|\le \sqrt{\frac{B}{2}}.
\]
Hence
\[
a^{-}_{\textup{KL}}(B;b)\ge \left[b-\sqrt{\frac{B}{2}}\right]_+.
\tag{4}
\]
Combining \((1)\), \((2)\), and \((4)\), we obtain
\[
a^{-}_{\textup{KL}}(B;b_t)
\ge
\left[
\left(\frac{C/2-t}{L_{\max}}\right)_+
-\frac{\sqrt{\Gamma}}{2}
\right]_+.
\tag{5}
\]
Substituting \((5)\) into \((3)\) gives
\[
\mathbb{E}_{\mu,P^{M,\textup{Alg}}}[(L-t)_+]
\ge
L_{\max}
\left[
\left(\frac{C/2-t}{L_{\max}}\right)_+
-\frac{\sqrt{\Gamma}}{2}
\right]_+.
\]
Finally, Corollary~\ref{cor:cvar-lower-bt} implies
\[
\textup{CVaR}_\alpha^{\mu\otimes P^{M,\textup{Alg}}}(L)
\ge
\min_{t\in\mathbb{R}}
\left\{
t+\frac{L_{\max}}{1-\alpha}
\left[
\left(\frac{C/2-t}{L_{\max}}\right)_+
-\frac{\sqrt{\Gamma}}{2}
\right]_+
\right\}.
\]
This proves the claim.
\end{proof}

\begin{corollary}[Balanced-pair]
\label{cor:balanced-two-point-cvar}
Under the assumptions of Lemma~\ref{lem:two-point-cvar}, if in addition
\[
L(M_1,x)+L(M_2,x)\ge L_{\max}
\qquad
\text{for all }x\in\mathcal{X},
\tag{A3}
\]
then
\[
\textup{CVaR}_\alpha^{\mu\otimes P^{M,\textup{Alg}}}(L)
\ge
L_{\max}\left[\frac12-\frac{\sqrt{\Gamma}}{2}\right]_+.
\]
\end{corollary}

\begin{proof}
By Lemma~\ref{lem:two-point-cvar}, it suffices to minimize
\[
F(t):=
t+\frac{L_{\max}}{1-\alpha}
\left[
\left(\frac12-\frac{t}{L_{\max}}\right)_+
-\frac{\sqrt{\Gamma}}{2}
\right]_+.
\]
Set
\[
c:=\frac{\sqrt{\Gamma}}{2}.
\]
We analyze \(F\) in three regions.

\smallskip
\noindent
\emph{Region 1: \(t\ge L_{\max}/2\).}
Then \(\left(\frac12-\frac{t}{L_{\max}}\right)_+=0\), so \(F(t)=t\), and hence
\[
\inf_{t\ge L_{\max}/2}F(t)=\frac{L_{\max}}{2}.
\]

\smallskip
\noindent
\emph{Region 2: \(L_{\max}(1/2-c)\le t\le L_{\max}/2\).}
Then
\[
\left(\frac12-\frac{t}{L_{\max}}\right)_+=\frac12-\frac{t}{L_{\max}},
\qquad
\frac12-\frac{t}{L_{\max}}-c\le 0,
\]
so again \(F(t)=t\), and therefore
\[
\inf_{L_{\max}(1/2-c)\le t\le L_{\max}/2}F(t)
=
L_{\max}\left(\frac12-c\right).
\]

\smallskip
\noindent
\emph{Region 3: \(t<L_{\max}(1/2-c)\).}
Then both positive parts are active, so
\[
F(t)
=
t+\frac{L_{\max}}{1-\alpha}\left(\frac12-\frac{t}{L_{\max}}-c\right)
=
t+\frac{L_{\max}(1/2-c)-t}{1-\alpha}.
\]
This is affine in \(t\), with derivative
\[
F'(t)=1-\frac{1}{1-\alpha}=-\frac{\alpha}{1-\alpha}<0.
\]
Hence \(F\) is decreasing on Region 3, and its infimum is attained at the right endpoint:
\[
\inf_{t<L_{\max}(1/2-c)}F(t)
=
L_{\max}\left(\frac12-c\right).
\]

Combining the three regions,
\[
\min_{t\in\mathbb{R}}F(t)
=
L_{\max}\left(\frac12-c\right)
=
L_{\max}\left(\frac12-\frac{\sqrt{\Gamma}}{2}\right),
\]
provided the quantity in parentheses is nonnegative. If it is negative, the lower bound is replaced by \(0\), since \(F(t)\ge 0\) for all \(t\). Therefore
\[
\min_{t\in\mathbb{R}}F(t)
=
L_{\max}\left[\frac12-\frac{\sqrt{\Gamma}}{2}\right]_+.
\]
This proves the corollary.
\end{proof}

\begin{proposition}[Bayesian \textup{CVaR} dominates Bayesian expected loss]
\label{prop:cvar-dominates-bayes-risk}
Let \(L\) be any integrable loss, and let \(\alpha\in[0,1)\). Then
\[
\textup{CVaR}_\alpha(L)\ge \mathbb{E}[L].
\]
Consequently, for any prior \(\mu\),
\[
\textup{CVaR}_\alpha^{\mu\otimes P^{M,\textup{Alg}}}(L)
\ge
\mathbb{E}_{\mu\otimes P^{M,\textup{Alg}}}[L].
\]
In particular, if for some prior \(\mu\) one has
\[
\mathbb{E}_{\mu\otimes P^{M,\textup{Alg}}}[L]\ge c\,\psi,
\]
then
\[
\textup{CVaR}_\alpha^{\mu\otimes P^{M,\textup{Alg}}}(L)\ge c\,\psi.
\]
Moreover,
\[
\sup_{\mu}\textup{CVaR}_\alpha^{\mu\otimes P^{M,\textup{Alg}}}(L)
\ge
\sup_{\mu}\mathbb{E}_{\mu\otimes P^{M,\textup{Alg}}}[L].
\]
\end{proposition}

\begin{proof}
By the dual representation of Expected Shortfall/\(\textup{CVaR}\),
\[
\textup{CVaR}_\alpha(L)
=
\sup\Bigl\{
\mathbb{E}[LZ]:
0\le Z\le \tfrac{1}{1-\alpha},\ \mathbb{E}[Z]=1
\Bigr\}.
\]
Since the choice \(Z\equiv 1\) is feasible, it follows that
\[
\textup{CVaR}_\alpha(L)\ge \mathbb{E}[L].
\]
Applying this under the Bayesian law \(\mu\otimes P^{M,\textup{Alg}}\) gives
\[
\textup{CVaR}_\alpha^{\mu\otimes P^{M,\textup{Alg}}}(L)
\ge
\mathbb{E}_{\mu\otimes P^{M,\textup{Alg}}}[L].
\]
The remaining claims follow immediately.
\end{proof}


\section{Examples}

\subsection{A passive estimation example}

We now illustrate the reusable two-point Bayesian \textup{CVaR} lemma in a passive estimation problem.

\paragraph{Setup.}
Let \(\theta\in\{-\Delta,+\Delta\}\), where \(\Delta>0\), and suppose we observe
\[
Y_1,\dots,Y_n \;\;\text{i.i.d.}\;\; \mathcal{N}(\theta,1).
\]
An estimator \(\widehat{\theta}=\widehat{\theta}(Y_1,\dots,Y_n)\) is any measurable function of the sample.
This fits the ISDM framework in the passive case:
\begin{itemize}
\item the model is \(M_\theta\), indexed by \(\theta\in\{-\Delta,+\Delta\}\);
\item there is no interaction, so the algorithm \(\textup{Alg}\) only outputs the final decision \(\widehat{\theta}\);
\item the transcript is simply
\[
X=(Y_1,\dots,Y_n).
\]
\end{itemize}

We equip this problem with the bounded estimation loss
\[
L(\theta,X):=\min\{|\widehat{\theta}(X)-\theta|,\;2\Delta\}.
\]
Then \(0\le L\le 2\Delta\).

Let
\[
M_+ := M_{+\Delta},\qquad M_-:=M_{-\Delta},
\qquad
\mu:=\frac12\delta_{M_+}+\frac12\delta_{M_-}.
\]

\begin{proposition}[Bayesian \textup{CVaR} lower bound for Gaussian mean estimation]
\label{prop:gaussian-mean-cvar}
For every estimator \(\widehat{\theta}\) and every \(\alpha\in[0,1)\),
\[
\textup{CVaR}_\alpha^{\mu\otimes P^{M,\textup{Alg}}}(L)
\ge
2\Delta\left[\frac12-\Delta\sqrt{\frac{n}{2}}\right]_+.
\]
In particular, choosing
\[
\Delta=\frac{1}{2\sqrt{2n}}
\]
gives
\[
\textup{CVaR}_\alpha^{\mu\otimes P^{M,\textup{Alg}}}(L)
\ge
\frac{1}{4\sqrt{2n}}.
\]
\end{proposition}

\begin{proof}
We verify the assumptions of Corollary~\ref{cor:balanced-two-point-cvar}.

\medskip
\noindent
\textbf{Step 1: Boundedness.}
By definition,
\[
0\le L(\theta,X)\le 2\Delta.
\]
Hence we may take
\[
L_{\max}:=2\Delta.
\]

\medskip
\noindent
\textbf{Step 2: Balanced-pair condition.}
Fix any realized estimate \(z\in\mathbb{R}\). Define
\[
a:=|z-\Delta|,
\qquad
b:=|z+\Delta|.
\]
By the triangle inequality,
\[
a+b = |z-\Delta|+|z+\Delta| \ge |(\,z+\Delta\,)-(\,z-\Delta\,)|=2\Delta.
\]
Therefore
\[
\min\{a,2\Delta\}+\min\{b,2\Delta\}\ge 2\Delta.
\]
Applying this with \(z=\widehat{\theta}(X)\), we obtain, for every transcript \(X\),
\[
L(+\Delta,X)+L(-\Delta,X)\ge 2\Delta.
\]
Thus the balanced-pair condition \((\mathrm{A3})\) of Corollary~\ref{cor:balanced-two-point-cvar} holds.

\medskip
\noindent
\textbf{Step 3: KL divergence between the two models.}
Under \(M_+\), the sample law is
\[
P_+ = \mathcal{N}(\Delta,1)^{\otimes n},
\]
and under \(M_-\), the sample law is
\[
P_- = \mathcal{N}(-\Delta,1)^{\otimes n}.
\]
Since the observations are i.i.d., the KL divergence adds:
\[
D_{\textup{KL}}(P_+\|P_-)
=
n\,D_{\textup{KL}}\bigl(\mathcal{N}(\Delta,1)\,\|\,\mathcal{N}(-\Delta,1)\bigr).
\]
For unit-variance Gaussians,
\[
D_{\textup{KL}}\bigl(\mathcal{N}(\mu_1,1)\,\|\,\mathcal{N}(\mu_2,1)\bigr)
=
\frac{(\mu_1-\mu_2)^2}{2}.
\]
Hence
\[
D_{\textup{KL}}(P_+\|P_-)
=
n\cdot \frac{(2\Delta)^2}{2}
=
2n\Delta^2.
\]
So we may take
\[
\Gamma:=2n\Delta^2.
\]

\medskip
\noindent
\textbf{Step 4: Apply the reusable two-point bound.}
Corollary~\ref{cor:balanced-two-point-cvar} now gives
\[
\textup{CVaR}_\alpha^{\mu\otimes P^{M,\textup{Alg}}}(L)
\ge
L_{\max}\left[\frac12-\frac{\sqrt{\Gamma}}{2}\right]_+.
\]
Substituting \(L_{\max}=2\Delta\) and \(\Gamma=2n\Delta^2\),
\[
\textup{CVaR}_\alpha^{\mu\otimes P^{M,\textup{Alg}}}(L)
\ge
2\Delta\left[\frac12-\frac{\sqrt{2n\Delta^2}}{2}\right]_+
=
2\Delta\left[\frac12-\Delta\sqrt{\frac{n}{2}}\right]_+.
\]
This proves the first claim.

\medskip
\noindent
\textbf{Step 5: Optimize the separation \(\Delta\).}
Consider
\[
\psi(\Delta):=
2\Delta\left(\frac12-\Delta\sqrt{\frac{n}{2}}\right)
=
\Delta-\sqrt{2n}\,\Delta^2.
\]
For \(\Delta\le 1/\sqrt{2n}\), the positive part is inactive, so maximizing the lower bound reduces to maximizing \(\psi(\Delta)\). Differentiating,
\[
\psi'(\Delta)=1-2\sqrt{2n}\,\Delta.
\]
Setting \(\psi'(\Delta)=0\) yields
\[
\Delta=\frac{1}{2\sqrt{2n}}.
\]
At this value,
\[
2\Delta=\frac{1}{\sqrt{2n}},
\qquad
\Delta\sqrt{\frac{n}{2}}=\frac14.
\]
Therefore
\[
\textup{CVaR}_\alpha^{\mu\otimes P^{M,\textup{Alg}}}(L)
\ge
\frac{1}{\sqrt{2n}}\left(\frac12-\frac14\right)
=
\frac{1}{4\sqrt{2n}}.
\]
This completes the proof.
\end{proof}

\subsection{Two-armed Gaussian bandits}

We now instantiate the same reusable two-point lemma for the two-armed Gaussian bandit.

Let \(T\ge 1\). Consider the two-armed Gaussian bandit with unit-variance rewards and horizon \(T\). For a parameter \(g>0\), define the two environments
\[
M_1:=\left(+\frac{g}{2},-\frac{g}{2}\right),
\qquad
M_2:=\left(-\frac{g}{2},+\frac{g}{2}\right).
\]
At each round \(t\), the algorithm \(\textup{Alg}\) chooses an arm \(A_t\in\{1,2\}\) based on the past history, then observes the reward
\[
Y_t=\mu_{A_t}+Z_t,
\qquad
Z_t\sim \mathcal{N}(0,1).
\]
Thus the transcript is
\[
H_T=(A_1,Y_1,\dots,A_T,Y_T),
\]
which fits the ISDM template with model \(M\in\{M_1,M_2\}\), decision \(A_t\), and observation \(Y_t\).

Let
\[
\mu=\frac12\delta_{M_1}+\frac12\delta_{M_2}.
\]
For \(a\in\{1,2\}\), let \(N_a(H_T)\) denote the number of pulls of arm \(a\) up to time \(T\).

\begin{theorem}[Bayesian \textup{CVaR} lower bound for two-armed Gaussian bandits]
\label{thm:two-armed-gaussian-cvar-reusable}
For every algorithm \(\textup{Alg}\) and every \(\alpha\in[0,1)\),
\[
\textup{CVaR}_\alpha^{\mu\otimes P^{M,\textup{Alg}}}(R_T)
\ge
gT\left[\frac12-\frac{g\sqrt{T}}{2\sqrt{2}}\right]_+.
\]
In particular, choosing
\[
g=\frac{1}{\sqrt{2T}}
\]
gives
\[
\textup{CVaR}_\alpha^{\mu\otimes P^{M,\textup{Alg}}}(R_T)
\ge
\frac{\sqrt{T}}{4\sqrt{2}}.
\]
\end{theorem}

\begin{proof}
We verify the assumptions of Corollary~\ref{cor:balanced-two-point-cvar}.

Under \(M_1\), arm \(1\) is optimal and arm \(2\) is suboptimal by gap \(g\), so
\[
R_T(M_1,H_T)=g\,N_2(H_T).
\]
Under \(M_2\), arm \(2\) is optimal and arm \(1\) is suboptimal by gap \(g\), so
\[
R_T(M_2,H_T)=g\,N_1(H_T).
\]
Hence, for every transcript \(h\),
\[
R_T(M_1,h)+R_T(M_2,h)
=
g\bigl(N_1(h)+N_2(h)\bigr)
=
gT.
\]
Also,
\[
0\le R_T(M_i,H_T)\le gT,
\qquad i=1,2.
\]
Thus we may take
\[
L(M,H_T):=R_T(M,H_T),
\qquad
L_{\max}:=gT.
\]
The balanced-pair condition \((\mathrm{A3})\) holds with this \(L_{\max}\).

Finally, for this two-armed Gaussian construction, the transcript KL divergence satisfies
\[
D_{\textup{KL}}(P^{M_1,\textup{Alg}}\|P^{M_2,\textup{Alg}})
=
\frac{g^2T}{2}.
\]
Hence we may take
\[
\Gamma:=\frac{g^2T}{2}.
\]

Applying Corollary~\ref{cor:balanced-two-point-cvar}, we obtain
\[
\textup{CVaR}_\alpha^{\mu\otimes P^{M,\textup{Alg}}}(R_T)
\ge
gT\left[\frac12-\frac{1}{2}\sqrt{\frac{g^2T}{2}}\right]_+
=
gT\left[\frac12-\frac{g\sqrt{T}}{2\sqrt{2}}\right]_+.
\]
This proves the first claim.

To optimize over \(g\), define
\[
\psi(g):=
gT\left(\frac12-\frac{g\sqrt{T}}{2\sqrt{2}}\right).
\]
For \(g\le \sqrt{2/T}\), the positive part is inactive, so maximizing the lower bound reduces to maximizing \(\psi(g)\). Differentiating,
\[
\psi'(g)=\frac{T}{2}-\frac{gT^{3/2}}{\sqrt{2}}.
\]
Setting \(\psi'(g)=0\) yields
\[
g=\frac{1}{\sqrt{2T}}.
\]
At this value,
\[
gT=\frac{\sqrt{T}}{\sqrt{2}},
\qquad
\frac{g\sqrt{T}}{2\sqrt{2}}=\frac14,
\]
and therefore
\[
\textup{CVaR}_\alpha^{\mu\otimes P^{M,\textup{Alg}}}(R_T)
\ge
\frac{\sqrt{T}}{\sqrt{2}}\left(\frac12-\frac14\right)
=
\frac{\sqrt{T}}{4\sqrt{2}}.
\]
This completes the proof.
\end{proof}
\begin{remark}[Interpretation and relation to the examples]
\label{rem:comparison-with-bayes-risk}
Proposition~\ref{prop:cvar-dominates-bayes-risk} shows that any Bayesian expected-loss lower bound immediately implies a Bayesian \(\textup{CVaR}_\alpha\) lower bound of the same order. Thus, the role of our explicit instantiations is not merely to recover what already follows from Bayes risk, but to show that the generalized-Fano \(\textup{CVaR}\) corollaries can be worked out directly in concrete problems.

For the passive Gaussian mean estimation example and the two-armed Gaussian bandit example, the explicit \(\textup{CVaR}_\alpha\) lower bounds obtained from the two-point template recover the same order as the corresponding Bayesian expected-loss benchmarks, namely \(n^{-1/2}\) in the estimation problem and \(\sqrt{T}\) in the two-armed bandit problem. In this sense, the examples verify that the corollary-based \(\textup{CVaR}\) approach is properly calibrated to the known Bayesian risk scaling in these canonical two-point settings.
\end{remark}
\fi


\iffalse
We now give a fully explicit instantiation of Corollaries~2--3 for the two-armed Gaussian bandit.

\begin{theorem}[Bayesian \textup{CVaR} lower bound for two-armed Gaussian bandits]
\label{thm:two-armed-gaussian-cvar}
Let \(T\ge 1\). Consider the two-armed Gaussian bandit with unit-variance rewards and horizon \(T\). For a parameter \(g>0\), define the two environments
\[
M_1:=\left(+\frac{g}{2},-\frac{g}{2}\right),
\qquad
M_2:=\left(-\frac{g}{2},+\frac{g}{2}\right),
\]
and let the prior be
\[
\mu=\frac12\delta_{M_1}+\frac12\delta_{M_2}.
\]
Then for every algorithm \(\textup{Alg}\) and every \(\alpha\in[0,1)\),
\[
\textup{CVaR}_\alpha^{\mu\otimes P^{M,\textup{Alg}}}(R_T)
\ge
gT\left[\frac12-\frac{g\sqrt{T}}{2\sqrt{2}}\right]_+.
\]
In particular, choosing
\[
g=\frac{1}{\sqrt{2T}}
\]
gives
\[
\textup{CVaR}_\alpha^{\mu\otimes P^{M,\textup{Alg}}}(R_T)
\ge
\frac{\sqrt{T}}{4\sqrt{2}}.
\]
\end{theorem}

\begin{proof}
We proceed in six steps.

\medskip
\noindent
\textbf{Step 1: Regret representation.}
Let \(N_a(H_T)\) denote the number of pulls of arm \(a\in\{1,2\}\) up to time \(T\).
Under \(M_1\), arm \(1\) is optimal and arm \(2\) is suboptimal by gap \(g\), so
\[
R_T(M_1,H_T)=g\,N_2(H_T).
\]
Under \(M_2\), arm \(2\) is optimal and arm \(1\) is suboptimal by gap \(g\), so
\[
R_T(M_2,H_T)=g\,N_1(H_T).
\]
Therefore, for every transcript \(h\),
\[
R_T(M_1,h)+R_T(M_2,h)
=
g\bigl(N_1(h)+N_2(h)\bigr)
=
gT.
\tag{1}
\]
Also,
\[
0\le R_T(M_i,H_T)\le gT,
\qquad i=1,2.
\]
Thus Corollaries~2--3 apply with
\[
L(M,H_T)=R_T(M,H_T),
\qquad
L_{\max}=gT.
\]

\medskip
\noindent
\textbf{Step 2: Apply Corollaries~2--3 with \(f=\textup{KL}\) and reference \(Q=P_2\).}
Let
\[
Q:=P^{M_2,\textup{Alg}}=:P_2.
\]
For each \(t\in\mathbb{R}\), Corollary~2 yields
\[
\mathbb{E}_{\mu,P^{M,\textup{Alg}}}\big[(R_T-t)_+\big]
\ge
gT\,a^{-}_{\textup{KL}}(B;b_t),
\tag{2}
\]
where
\[
B
=
\mathbb{E}_{M\sim\mu}
D_{\textup{KL}}(P^{M,\textup{Alg}}\|Q),
\qquad
b_t
=
\frac{1}{gT}\,
\mathbb{E}_{\mu,Q}\big[(R_T-t)_+\big].
\]
Corollary~3 then gives
\[
\textup{CVaR}_\alpha^{\mu\otimes P^{M,\textup{Alg}}}(R_T)
\ge
\min_{t\in\mathbb{R}}
\left\{
t+\frac{gT}{1-\alpha}\,a^{-}_{\textup{KL}}(B;b_t)
\right\}.
\tag{3}
\]

\medskip
\noindent
\textbf{Step 3: Lower bound the reference hinge term \(b_t\).}
By definition,
\[
b_t
=
\frac{1}{gT}\cdot \frac12
\mathbb{E}_{P_2}\big[(R_T(M_1,H_T)-t)_+\big]
\]
\[
+
\frac{1}{gT}\cdot \frac12
\mathbb{E}_{P_2}\big[(R_T(M_2,H_T)-t)_+\big].
\]
Using the identity \eqref{1}, for every transcript \(h\),
\[
R_T(M_1,h)+R_T(M_2,h)=gT.
\]
Since \(x\mapsto (x-t)_+\) is convex,
\[
\frac12\Bigl((R_T(M_1,h)-t)_+ + (R_T(M_2,h)-t)_+\Bigr)\]
\[
\ge
\left(\frac{R_T(M_1,h)+R_T(M_2,h)}{2}-t\right)_+
=
\left(\frac{gT}{2}-t\right)_+.
\]
Taking \(\mathbb{E}_{P_2}\) on both sides yields
\[
\frac12\mathbb{E}_{P_2}\big[(R_T(M_1,H_T)-t)_+\big]
+
\frac12\mathbb{E}_{P_2}\big[(R_T(M_2,H_T)-t)_+\big]\]
\[
\ge
\left(\frac{gT}{2}-t\right)_+.
\]
Hence
\[
b_t
\ge
\left[\frac12-\frac{t}{gT}\right]_+.
\tag{4}
\]

\medskip
\noindent
\textbf{Step 4: Compute the divergence budget \(B\).}
By definition,
\[
B
=
\frac12 D_{\textup{KL}}(P_1\|P_2)
+
\frac12 D_{\textup{KL}}(P_2\|P_2)
=
\frac12 D_{\textup{KL}}(P_1\|P_2).
\]
For the two-armed Gaussian construction, the transcript KL divergence is
\[
D_{\textup{KL}}(P_1\|P_2)=\frac{g^2T}{2},
\]
as computed in the quantile lower-bound paper for the same pair of environments.
Therefore
\[
B=\frac{g^2T}{4}.
\tag{5}
\]

\medskip
\noindent
\textbf{Step 5: Invert the Bernoulli KL ball.}
Pinsker's inequality for Bernoulli laws implies that
\[
D_{\textup{KL}}(\textup{Bern}(a)\|\textup{Bern}(b))\le B
\quad\Longrightarrow\quad
|a-b|\le \sqrt{\frac{B}{2}}.
\]
Hence
\[
a^{-}_{\textup{KL}}(B;b)\ge \left[b-\sqrt{\frac{B}{2}}\right]_+.
\tag{6}
\]
Combining \eqref{5} and \eqref{6},
\[
a^{-}_{\textup{KL}}(B;b_t)
\ge
\left[b_t-\frac{g\sqrt{T}}{2\sqrt{2}}\right]_+.
\tag{7}
\]
Using the lower bound \eqref{4} on \(b_t\), we obtain
\[
a^{-}_{\textup{KL}}(B;b_t)
\ge
\left[
\left(\frac12-\frac{t}{gT}\right)_+
-
\frac{g\sqrt{T}}{2\sqrt{2}}
\right]_+.
\tag{8}
\]

\medskip
\noindent
\textbf{Step 6: Minimize the \(\textup{CVaR}\) lower bound.}
Substituting \eqref{8} into \eqref{3} gives
\[
\textup{CVaR}_\alpha^{\mu\otimes P^{M,\textup{Alg}}}(R_T)
\ge
\min_{t\in\mathbb{R}}
\left\{
t+\frac{gT}{1-\alpha}
\left[
\left(\frac12-\frac{t}{gT}\right)_+
-
\frac{g\sqrt{T}}{2\sqrt{2}}
\right]_+
\right\}.
\tag{9}
\]
Set
\[
c:=\frac{g\sqrt{T}}{2\sqrt{2}}.
\]
We analyze the minimum of the right-hand side of \eqref{9}.

\smallskip
\noindent
\emph{Region 1: \(t\ge gT/2\).}
Then \(\left(\frac12-\frac{t}{gT}\right)_+=0\), so the bracket in \eqref{9} vanishes and the objective equals \(t\). Its minimum on this region is \(gT/2\).

\smallskip
\noindent
\emph{Region 2: \(gT(1/2-c)\le t\le gT/2\).}
Then \(\left(\frac12-\frac{t}{gT}\right)_+\ge 0\), but
\[
\frac12-\frac{t}{gT}-c\le 0,
\]
so the bracket in \eqref{9}  vanishes and the objective equals \(t\). Its minimum on this region is
\[
gT\left(\frac12-c\right).
\]

\smallskip
\noindent
\emph{Region 3: \(t<gT(1/2-c)\).}
Then the bracket in \eqref{9} is positive, so the objective becomes
\[
t+\frac{gT}{1-\alpha}\left(\frac12-\frac{t}{gT}-c\right)
=
t+\frac{gT(1/2-c)-t}{1-\alpha}.
\]
Its derivative with respect to \(t\) is
\[
1-\frac{1}{1-\alpha}
=
-\frac{\alpha}{1-\alpha}<0.
\]
Hence the objective is strictly decreasing on this region, and its minimum is attained at the largest admissible value of \(t\), namely
\[
t=gT\left(\frac12-c\right).
\]

\smallskip
Thus the global minimum in \eqref{9} is attained at
\[
t^\star = gT\left(\frac12-c\right),
\]
and equals
\[
gT\left(\frac12-c\right)
=
gT\left(\frac12-\frac{g\sqrt{T}}{2\sqrt{2}}\right).
\]
Adding the positive-part notation yields
\[
\textup{CVaR}_\alpha^{\mu\otimes P^{M,\textup{Alg}}}(R_T)
\ge
gT\left[\frac12-\frac{g\sqrt{T}}{2\sqrt{2}}\right]_+.
\]

Finally, optimize over \(g\). Consider
\[
\psi(g)
=
gT\left(\frac12-\frac{g\sqrt{T}}{2\sqrt{2}}\right)
=
\frac{gT}{2}-\frac{g^2T^{3/2}}{2\sqrt{2}}.
\]
Differentiating,
\[
\psi'(g)=\frac{T}{2}-\frac{gT^{3/2}}{\sqrt{2}}.
\]
Setting \(\psi'(g)=0\) yields
\[
g=\frac{1}{\sqrt{2T}}.
\]
Substituting this into the bound gives
\[
gT=\frac{\sqrt{T}}{\sqrt{2}},
\qquad
\frac{g\sqrt{T}}{2\sqrt{2}}=\frac14,
\]
and therefore
\[
\textup{CVaR}_\alpha^{\mu\otimes P^{M,\textup{Alg}}}(R_T)
\ge
\frac{\sqrt{T}}{\sqrt{2}}\left(\frac12-\frac14\right)
=
\frac{\sqrt{T}}{4\sqrt{2}}.
\]
This completes the proof.
\end{proof}

\fi
\iffalse
\subsection{Multi-armed Gaussian Bandits}
\section{A Bayesian CVaR lower bound for $K$-armed Gaussian bandits}

We give a simple worked example of the reference-model/Hellinger route for Bayesian CVaR lower bounds.
The result is a fixed-prior lower bound for stochastic $K$-armed Gaussian bandits.

\subsection{Setup}

Consider a stochastic $K$-armed Gaussian bandit with horizon $T$ and unit-variance rewards.
For each arm $a \in [K] := \{1,\dots,K\}$, pulling arm $a$ at time $t$ yields a reward
\[
Y_t = \mu_a + Z_t, \qquad Z_t \sim \mathcal N(0,1),
\]
with the noises independent across time.

For a gap parameter $\Delta>0$, define the model family
\[
M_i : \qquad \mu_i = \Delta,\qquad \mu_j = 0\ \ \text{for all }j\neq i,\qquad i\in[K].
\]
Thus, under $M_i$, arm $i$ is uniquely optimal and every suboptimal pull incurs instantaneous regret $\Delta$.

Let the prior be uniform on the hard family:
\[
I \sim \mu := \textup{Unif}([K]).
\]
Let $\textup{Alg}$ be any (possibly randomized, adaptive) bandit algorithm.
For a bandit history
\[
H_T := (A_1,Y_1,\dots,A_T,Y_T),
\]
write $N_i(H_T) := \sum_{t=1}^T \mathbf 1\{A_t=i\}$ for the number of pulls of arm $i$.

The cumulative regret under model $M_i$ is
\[
R_T(M_i,H_T) := \sum_{t=1}^T \bigl(\mu^\star(M_i)-\mu_{A_t}(M_i)\bigr)
= \Delta\,(T-N_i(H_T)).
\]
In particular,
\[
0\le R_T(M_i,H_T)\le \Delta T.
\]

We choose the reference model
\[
\bar M:\qquad \mu_1=\cdots=\mu_K=0,
\]
and denote by
\[
Q := P^{\bar M,\textup{Alg}},
\qquad
P_i := P^{M_i,\textup{Alg}}
\]
the transcript laws under $\bar M$ and $M_i$ respectively.

\subsection{A general Hellinger--CVaR reduction}

The following lemma is the direct Hellinger specialization of Corollaries~2--3 in the generalized-Fano CVaR paper.

\begin{lemma}[Hellinger--CVaR reduction]
\label{lem:hellinger-cvar-reduction}
Let $L$ be any loss satisfying $0\le L\le L_{\max}$ almost surely.
Fix a prior $\mu$ on models and a reference law $Q$ on transcripts.
Define, for each $t\in\mathbb R$,
\[
b_t := \frac{1}{L_{\max}} \mathbb{E}_{\mu,Q}[(L-t)_+],
\qquad
B := \mathbb{E}_{M\sim\mu} D_{H^2}(P^{M,\textup{Alg}}\|Q).
\]
Then
\[
\mathbb{E}_{\mu,P^{M,\textup{Alg}}}[(L-t)_+]
\ge
L_{\max}\, a^{-}_{H^2}(B;b_t),
\]
where
\[
a^{-}_{H^2}(B;b)
:=
\inf\Bigl\{a\in[0,1]:
D_{H^2}(\textup{Bern}(a)\|\textup{Bern}(b))\le B
\Bigr\}.
\]
Consequently,
\[
\textup{CVaR}_\alpha^{\mu\otimes P^{M,\textup{Alg}}}(L)
\ge
\min_{t\in\mathbb R}
\left\{
t+\frac{L_{\max}}{1-\alpha}\,a^{-}_{H^2}(B;b_t)
\right\}.
\]
\end{lemma}

\begin{proof}
Apply Corollary~2 of the generalized-Fano CVaR paper with $f=H^2$ and the bounded hinge
\[
\phi_t(\ell)=\frac{(\ell-t)_+}{L_{\max}},
\]
which yields
\[
\mathbb{E}_{\mu,P^{M,\textup{Alg}}}[(L-t)_+]
\ge
L_{\max}\, a^{-}_{H^2}(B;b_t).
\]
Then apply Corollary~3, which uses the Rockafellar--Uryasev identity
\[
\textup{CVaR}_\alpha(L)=\min_{t\in\mathbb R}\left\{t+\frac{1}{1-\alpha}\mathbb{E}[(L-t)_+]\right\}.
\]
This gives the stated bound. \qedhere
\end{proof}

\begin{lemma}[Bernoulli Hellinger inversion]
\label{lem:bernoulli-hellinger-inversion}
For every $a,b\in[0,1]$,
\[
D_{H^2}(\textup{Bern}(a)\|\textup{Bern}(b))
=
1-\sqrt{ab}-\sqrt{(1-a)(1-b)}.
\]
Moreover, if
\[
D_{H^2}(\textup{Bern}(a)\|\textup{Bern}(b))\le B,
\]
then
\[
a\ge \bigl(\sqrt b-\sqrt{2B}\bigr)_+^2.
\]
Equivalently,
\[
a^{-}_{H^2}(B;b)\ge \bigl(\sqrt b-\sqrt{2B}\bigr)_+^2.
\]
\end{lemma}

\begin{proof}
The formula for Bernoulli Hellinger distance is standard and follows from direct substitution into
\[
D_{H^2}(P\|Q)=1-\int \sqrt{\frac{dP}{d\lambda}\frac{dQ}{d\lambda}}\,d\lambda.
\]
For the inversion, note that
\[
D_{H^2}(\textup{Bern}(a)\|\textup{Bern}(b))
=
\frac12\Bigl(
(\sqrt a-\sqrt b)^2 + (\sqrt{1-a}-\sqrt{1-b})^2
\Bigr).
\]
Hence
\[
(\sqrt a-\sqrt b)^2 \le 2B.
\]
Therefore
\[
\sqrt a \ge (\sqrt b-\sqrt{2B})_+,
\]
and squaring yields
\[
a\ge (\sqrt b-\sqrt{2B})_+^2.
\]
Taking the infimum over admissible $a$ proves the last claim.
\end{proof}

\begin{corollary}[Explicit Hellinger--CVaR reduction]
\label{cor:explicit-hellinger-cvar}
Under the assumptions of Lemma~\ref{lem:hellinger-cvar-reduction},
\[
\textup{CVaR}_\alpha^{\mu\otimes P^{M,\textup{Alg}}}(L)
\ge
\min_{t\in\mathbb R}
\left\{
t+\frac{L_{\max}}{1-\alpha}
\bigl(\sqrt{b_t}-\sqrt{2B}\bigr)_+^2
\right\}.
\]
\end{corollary}

\begin{proof}
Combine Lemma~\ref{lem:hellinger-cvar-reduction} with Lemma~\ref{lem:bernoulli-hellinger-inversion}.
\end{proof}

\subsection{Bandit-specific ingredients}

We now compute the two quantities needed in Corollary~\ref{cor:explicit-hellinger-cvar}:
the ghost hinge $b_t$ and the average Hellinger budget $B$.

\begin{lemma}[Reference-model representation of the regret]
\label{lem:bandit-regret-form}
Under model $M_i$,
\[
R_T(M_i,H_T)=\Delta\,(T-N_i(H_T)).
\]
In particular,
\[
0\le R_T(M_i,H_T)\le \Delta T.
\]
\end{lemma}

\begin{proof}
Under $M_i$, arm $i$ has mean $\Delta$ and every other arm has mean $0$, so the instantaneous regret at time $t$ is
\[
\Delta - \mu_{A_t}(M_i)
=
\Delta\,\mathbf 1\{A_t\neq i\}.
\]
Summing over $t=1,\dots,T$ gives
\[
R_T(M_i,H_T)=\Delta \sum_{t=1}^T \mathbf 1\{A_t\neq i\}
=\Delta\,(T-N_i(H_T)).
\]
The bounds are immediate.
\end{proof}

\begin{lemma}[Ghost hinge lower bound]
\label{lem:ghost-hinge-bandit}
Let
\[
L(M_i,H_T):=R_T(M_i,H_T),
\qquad
L_{\max}:=\Delta T.
\]
Then for every $t\in\mathbb R$,
\[
b_t
=
\frac{1}{\Delta T}\mathbb{E}_{\mu,Q}\bigl[(R_T-t)_+\bigr]
\ge
\left[1-\frac1K-\frac{t}{\Delta T}\right]_+.
\]
\end{lemma}

\begin{proof}
Under the reference law $Q=P^{\bar M,\textup{Alg}}$, the transcript $H_T$ is generated when all arms are statistically identical.
The random index $I\sim\textup{Unif}([K])$ is independent of $H_T$ under $\mu\otimes Q$.
Conditioning on $H_T$ and using Lemma~\ref{lem:bandit-regret-form},
\[
\mathbb{E}_{\mu,Q}[R_T\mid H_T]
=
\frac1K\sum_{i=1}^K \Delta\,(T-N_i(H_T))
=
\Delta\left(T-\frac1K\sum_{i=1}^K N_i(H_T)\right).
\]
Since $\sum_{i=1}^K N_i(H_T)=T$ almost surely,
\[
\mathbb{E}_{\mu,Q}[R_T\mid H_T]
=
\Delta T\left(1-\frac1K\right).
\]
Therefore
\[
\mathbb{E}_{\mu,Q}[R_T] = 
\mathbb{E}_{Q}\!\left[\mathbb{E}_{\mu,Q}[R_T\mid H_T]\right]
=
\Delta T\left(1-\frac1K\right).
\]
Now use the convexity of $x\mapsto (x-t)_+$:
\[
\mathbb{E}_{\mu,Q}[(R_T-t)_+]
\ge
\bigl(\mathbb{E}_{\mu,Q}[R_T]-t\bigr)_+
=
\left(\Delta T\left(1-\frac1K\right)-t\right)_+.
\]
Dividing by $\Delta T$ yields
\[
b_t\ge \left[1-\frac1K-\frac{t}{\Delta T}\right]_+.
\]
\end{proof}

\begin{lemma}[One-step Gaussian Hellinger distance]
\label{lem:gaussian-hellinger}
For unit-variance Gaussians,
\[
D_{H^2}\bigl(\mathcal N(\Delta,1)\,\|\,\mathcal N(0,1)\bigr)
=
1-e^{-\Delta^2/8}.
\]
In particular,
\[
D_{H^2}\bigl(\mathcal N(\Delta,1)\,\|\,\mathcal N(0,1)\bigr)
\le \frac{\Delta^2}{8}.
\]
\end{lemma}

\begin{proof}
Let
\[
p(x)=\frac{1}{\sqrt{2\pi}}e^{-(x-\Delta)^2/2},
\qquad
q(x)=\frac{1}{\sqrt{2\pi}}e^{-x^2/2}.
\]
Then
\[
\int_{\mathbb R}\sqrt{p(x)q(x)}\,dx
=
\frac{1}{\sqrt{2\pi}}\int_{\mathbb R}
\exp\!\left(
-\frac{(x-\Delta)^2+x^2}{4}
\right)\,dx.
\]
Completing the square,
\[
\frac{(x-\Delta)^2+x^2}{4}
=
\frac{(x-\Delta/2)^2}{2}+\frac{\Delta^2}{8}.
\]
Hence
\[
\int_{\mathbb R}\sqrt{p(x)q(x)}\,dx
=
e^{-\Delta^2/8}
\cdot
\frac{1}{\sqrt{2\pi}}
\int_{\mathbb R} e^{-(x-\Delta/2)^2/2}\,dx
=
e^{-\Delta^2/8}.
\]
Therefore
\[
D_{H^2}(P\|Q)=1-\int \sqrt{pq}
=
1-e^{-\Delta^2/8}.
\]
Finally, using $1-e^{-u}\le u$ for $u\ge 0$ gives
\[
1-e^{-\Delta^2/8}\le \frac{\Delta^2}{8}.
\]
\end{proof}

\begin{lemma}[Average Hellinger budget under the zero reference]
\label{lem:hellinger-budget-bandit}
Let
\[
B:=\mathbb{E}_{I\sim\mu}D_{H^2}(P_I\|Q).
\]
Then
\[
B\le \frac{\Delta^2T}{2K}.
\]
\end{lemma}

\begin{proof}
Fix $i\in[K]$.
We compare the transcript laws $P_i$ and $Q$.

At round $t$, conditional on the past history $H_{t-1}$, the algorithm chooses an arm $A_t$ according to the same kernel under both $P_i$ and $Q$ because the algorithm itself is unchanged.
The only difference is the conditional law of the reward $Y_t$:
\[
Y_t \mid (H_{t-1},A_t=a)
\sim
\begin{cases}
\mathcal N(\Delta,1), & a=i \text{ under }P_i,\\
\mathcal N(0,1), & a\neq i \text{ under }P_i,\\
\mathcal N(0,1), & \text{for all }a \text{ under }Q.
\end{cases}
\]
Thus, under $P_i$ versus $Q$, the only informative rounds are those on which $A_t=i$.

Using the chain rule for KL divergence for sequential experiments,
\[
D_{\textup{KL}}(P_i\|Q)
=
\sum_{t=1}^T \mathbb{E}_{P_i}\!\left[
D_{\textup{KL}}\!\left(
P_i(Y_t\mid H_{t-1},A_t)\,\|\,Q(Y_t\mid H_{t-1},A_t)
\right)
\right].
\]
The one-step conditional KL is $0$ if $A_t\neq i$, and equals
\[
D_{\textup{KL}}\bigl(\mathcal N(\Delta,1)\,\|\,\mathcal N(0,1)\bigr)=\frac{\Delta^2}{2}
\]
if $A_t=i$.
Hence
\[
D_{\textup{KL}}(P_i\|Q)
=
\frac{\Delta^2}{2}\,\mathbb{E}_{P_i}[N_i(H_T)].
\]
Now use $D_{H^2}\le D_{\textup{KL}}$:
\[
D_{H^2}(P_i\|Q)\le \frac{\Delta^2}{2}\,\mathbb{E}_{P_i}[N_i(H_T)].
\]
At this point one may keep $\mathbb{E}_{P_i}[N_i]$.
For a uniform-prior upper bound, we only need a crude estimate:
\[
\mathbb{E}_{P_i}[N_i(H_T)]\le T.
\]
Averaging over $i$ would only give $B\le \Delta^2T/2$, which is too weak.
So instead we use the same chain rule with the expectation taken under $Q$.
By nonnegativity of KL and the fact that the algorithmic action kernel is identical under both laws,
the same decomposition gives
\[
D_{\textup{KL}}(P_i\|Q)
=
\frac{\Delta^2}{2}\,\mathbb{E}_Q[N_i(H_T)].
\]
Indeed, under $Q$ the only per-round discrepancy arises when arm $i$ is pulled, and the conditional log-likelihood ratio has $Q$-expectation equal to the KL divergence between the two reward laws.
Therefore
\[
D_{H^2}(P_i\|Q)\le D_{\textup{KL}}(P_i\|Q)=\frac{\Delta^2}{2}\,\mathbb{E}_Q[N_i(H_T)].
\]
Averaging over $I\sim\textup{Unif}([K])$,
\[
B
\le
\frac{\Delta^2}{2K}\sum_{i=1}^K \mathbb{E}_Q[N_i(H_T)].
\]
Since $\sum_{i=1}^K N_i(H_T)=T$ almost surely under $Q$,
\[
\sum_{i=1}^K \mathbb{E}_Q[N_i(H_T)] = T.
\]
Hence
\[
B\le \frac{\Delta^2T}{2K}.
\]
\end{proof}

\subsection{The main bound}

\begin{theorem}[Bayesian CVaR lower bound for $K$-armed Gaussian bandits]
\label{thm:k-armed-gaussian-cvar}
Consider the $K$-armed Gaussian bandit family $\{M_i\}_{i=1}^K$ defined above, with prior
\[
I\sim\mu=\textup{Unif}([K]).
\]
Let $R_T$ denote the cumulative regret.
Then for every algorithm $\textup{Alg}$ and every $\alpha\in[0,1)$,
\[
\textup{CVaR}_\alpha^{\mu\otimes P^{M,\textup{Alg}}}(R_T)
\ge
\min_{t\in\mathbb R}
\left\{
t+\frac{\Delta T}{1-\alpha}
\left(
\sqrt{\left[1-\frac1K-\frac{t}{\Delta T}\right]_+}
-
\Delta\sqrt{\frac{T}{K}}
\right)_+^2
\right\}.
\]
\end{theorem}

\begin{proof}
Apply Corollary~\ref{cor:explicit-hellinger-cvar} with
\[
L=R_T,
\qquad
L_{\max}=\Delta T,
\qquad
Q=P^{\bar M,\textup{Alg}}.
\]
By Lemma~\ref{lem:ghost-hinge-bandit},
\[
b_t\ge \left[1-\frac1K-\frac{t}{\Delta T}\right]_+.
\]
By Lemma~\ref{lem:hellinger-budget-bandit},
\[
B\le \frac{\Delta^2T}{2K},
\qquad\text{so}\qquad
\sqrt{2B}\le \Delta\sqrt{\frac{T}{K}}.
\]
Substituting into Corollary~\ref{cor:explicit-hellinger-cvar} yields
\[
\textup{CVaR}_\alpha(R_T)
\ge
\min_{t\in\mathbb R}
\left\{
t+\frac{\Delta T}{1-\alpha}
\left(
\sqrt{\left[1-\frac1K-\frac{t}{\Delta T}\right]_+}
-
\Delta\sqrt{\frac{T}{K}}
\right)_+^2
\right\}.
\]
\end{proof}

\subsection{A simpler closed form}

The previous theorem already gives an explicit lower bound.
We now evaluate the minimization.

Write
\[
a_K := 1-\frac1K,
\qquad
\gamma := \Delta\sqrt{\frac{T}{K}}.
\]
Since $0\le R_T\le \Delta T$, the Rockafellar--Uryasev minimization may be restricted to $t\in[0,\Delta T]$.
Set
\[
x:=\frac{t}{\Delta T}\in[0,1].
\]
Then Theorem~\ref{thm:k-armed-gaussian-cvar} becomes
\[
\textup{CVaR}_\alpha(R_T)
\ge
\Delta T \cdot
\min_{x\in[0,1]}
\left\{
x+\frac{1}{1-\alpha}
\bigl(\sqrt{[a_K-x]_+}-\gamma\bigr)_+^2
\right\}.
\]
The minimizer can be computed explicitly.

\begin{lemma}[Evaluation of the one-dimensional minimization]
\label{lem:minimization}
Let
\[
F_{\alpha,a,\gamma}(x)
:=
x+\frac{1}{1-\alpha}\bigl(\sqrt{[a-x]_+}-\gamma\bigr)_+^2,
\qquad x\in[0,1],
\]
where $a\in(0,1]$, $\alpha\in(0,1)$, and $\gamma\ge 0$.
Then
\[
\min_{x\in[0,1]}F_{\alpha,a,\gamma}(x)
=
\begin{cases}
a-\dfrac{\gamma^2}{\alpha}, & \gamma\le \alpha\sqrt a,\\[1.5ex]
\dfrac{(\sqrt a-\gamma)_+^2}{1-\alpha}, & \gamma> \alpha\sqrt a.
\end{cases}
\]
\end{lemma}

\begin{proof}
There are two regions.

\paragraph{Region 1: $x\ge a-\gamma^2$.}
Then $\sqrt{[a-x]_+}\le \gamma$, so the positive part vanishes and
\[
F_{\alpha,a,\gamma}(x)=x.
\]
Hence on this region the minimum is attained at $x=a-\gamma^2$ and equals
\[
a-\gamma^2.
\]

\paragraph{Region 2: $0\le x\le a-\gamma^2$.}
Set
\[
y:=\sqrt{a-x},
\]
so that $y\in[\gamma,\sqrt a]$ and $x=a-y^2$.
Then
\[
F_{\alpha,a,\gamma}(x)
=
a-y^2+\frac{(y-\gamma)^2}{1-\alpha}
=:G(y).
\]
Expanding,
\[
G(y)
=
a+\frac{\alpha}{1-\alpha}y^2-\frac{2\gamma}{1-\alpha}y+\frac{\gamma^2}{1-\alpha}.
\]
This is a convex quadratic in $y$.
Its derivative is
\[
G'(y)=\frac{2\alpha}{1-\alpha}y-\frac{2\gamma}{1-\alpha},
\]
so the unique critical point is
\[
y^\star=\frac{\gamma}{\alpha}.
\]

If $\gamma\le \alpha\sqrt a$, then $y^\star\in[\gamma,\sqrt a]$, and thus $y^\star$ is the minimizer in Region~2.
Substituting $y^\star=\gamma/\alpha$,
\[
G(y^\star)=a-\frac{\gamma^2}{\alpha}.
\]

If $\gamma>\alpha\sqrt a$, then $y^\star>\sqrt a$, so on the interval $[\gamma,\sqrt a]$ the function $G$ is decreasing, and the minimum in Region~2 is attained at $y=\sqrt a$, i.e. at $x=0$.
This yields
\[
G(\sqrt a)=\frac{(\sqrt a-\gamma)_+^2}{1-\alpha}.
\]

Finally, compare with Region~1.
When $\gamma\le \alpha\sqrt a$,
\[
a-\frac{\gamma^2}{\alpha}\le a-\gamma^2,
\]
so the Region~2 value is smaller.
When $\gamma>\alpha\sqrt a$, the minimizer is at $x=0$ and the value is
\[
\frac{(\sqrt a-\gamma)_+^2}{1-\alpha}.
\]
This proves the claim.
\end{proof}

Combining the theorem and the minimization lemma gives the following explicit formula.

\begin{corollary}[Closed-form lower bound]
\label{cor:closed-form}
Let
\[
a_K:=1-\frac1K,
\qquad
\gamma:=\Delta\sqrt{\frac{T}{K}}.
\]
Then
\[
\textup{CVaR}_\alpha^{\mu\otimes P^{M,\textup{Alg}}}(R_T)
\ge
\Delta T \times
\begin{cases}
a_K-\dfrac{\gamma^2}{\alpha}, & \gamma\le \alpha\sqrt{a_K},\\[1.5ex]
\dfrac{(\sqrt{a_K}-\gamma)_+^2}{1-\alpha}, & \gamma>\alpha\sqrt{a_K}.
\end{cases}
\]
\end{corollary}

\begin{proof}
Apply Lemma~\ref{lem:minimization} with $a=a_K$ and $\gamma=\Delta\sqrt{T/K}$ to the bound of Theorem~\ref{thm:k-armed-gaussian-cvar}.
\end{proof}

\subsection{Rate consequence}

\begin{corollary}[$\sqrt{KT}$ scaling]
\label{cor:sqrtKT}
Fix $\alpha\in(0,1)$ and $K\ge 2$.
Choose
\[
\Delta = c_\alpha \sqrt{\frac{K}{T}},
\]
where $c_\alpha>0$ is small enough that
\[
c_\alpha \le \alpha\sqrt{1-\frac1K}.
\]
Then there exists a constant $c'_\alpha>0$ depending only on $\alpha$ such that
\[
\textup{CVaR}_\alpha^{\mu\otimes P^{M,\textup{Alg}}}(R_T)
\ge
c'_\alpha \sqrt{KT}.
\]
\end{corollary}

\begin{proof}
With the stated choice,
\[
\gamma = \Delta\sqrt{\frac{T}{K}} = c_\alpha,
\]
so the first branch of Corollary~\ref{cor:closed-form} applies:
\[
\textup{CVaR}_\alpha(R_T)
\ge
\Delta T\left(a_K-\frac{c_\alpha^2}{\alpha}\right).
\]
Since $a_K=1-1/K\ge 1/2$ for $K\ge 2$, choosing $c_\alpha$ sufficiently small ensures
\[
a_K-\frac{c_\alpha^2}{\alpha}\ge c''_\alpha>0.
\]
Therefore
\[
\textup{CVaR}_\alpha(R_T)
\ge
c''_\alpha \Delta T
=
c''_\alpha c_\alpha \sqrt{KT}.
\]
Setting $c'_\alpha:=c''_\alpha c_\alpha$ proves the claim.
\end{proof}


\fi





%------------------------------------------------%
\iffalse
\begin{corollary}[Clipped mean]
\label{cor:upper-clipped}
Assume $0\le L(M,X)\le L_{\max}$ and set $\phi_{\mathrm{clip}}(l)=\min\{\,l/L_{\max},\,1\,\}$.
Let $\rho_{\mathrm{clip},Q}:=\mathbb{E}_{M\sim\mu,\,X\sim Q}[\phi_{\mathrm{clip}}(L(M,X))]$.
Then
\[
\mathbb{E}_{M\sim\mu,\,X\sim P^{M,\textup{ALG}}}\!\big[\phi_{\mathrm{clip}}(L(M,X))\big]
\;\le\;
\inf_{Q\in\Delta(\mathcal X)}\;
\inf\Big\{\theta\in[0,1]:
\ \mathbb{E}_{M\sim\mu} D_f\!\big(P^{M,\textup{ALG}}\big\|Q\big)
\ <\ d_{f,\theta}\!\big(\rho_{\mathrm{clip},Q}\big)\Big\}.
\]
\end{corollary}

\begin{corollary}[Exponential transform / Laplace functional]
\label{cor:upper-exp}
Assume $L(M,X)\ge 0$ and fix $\lambda>0$. Let $\phi_\lambda(l)=e^{-\lambda l}$ and
$\rho_{\lambda,Q}:=\mathbb{E}_{M\sim\mu,\,X\sim Q}\!\big[e^{-\lambda L(M,X)}\big]$.
Then
\[
\mathbb{E}_{M\sim\mu,\,X\sim P^{M,\textup{ALG}}}\!\big[e^{-\lambda L(M,X)}\big]
\;\le\;
\inf_{Q\in\Delta(\mathcal X)}\;
\inf\Big\{\theta\in[0,1]:
\ \mathbb{E}_{M\sim\mu} D_f\!\big(P^{M,\textup{ALG}}\big\|Q\big)
\ <\ d_{f,\theta}\!\big(\rho_{\lambda,Q}\big)\Big\}.
\]
\end{corollary}
\fi

\iffalse
\begin{corollary}[CVaR lower bound via bounded transforms]
\label{cor:cvar-bt}
Assume $0\le L(M,X)\le L_{\max}$. For $t\in\mathbb R$ define the bounded hinge
\[
\phi_t(l):=\frac{(l-t)_+}{L_{\max}}\in[0,1],\qquad
a_t:=\mathbb{E}_{\mu,\;P^{M,\textup{ALG}}}[\phi_t(L)],\quad
b_t:=\mathbb{E}_{\mu,\;Q}[\phi_t(L)].
\]
Let
\[
B:=\mathbb{E}_{M\sim\mu}\, D_f\!\big(P^{M,\textup{ALG}}\,\big\|\,Q\big),
\qquad
\underline a_f(B;b)\;:=\;\sup\Big\{a\in[0,1]:\ D_f\!\big(\mathrm{Bern}(a)\,\big\|\,\mathrm{Bern}(b)\big)\le B\Big\}.
\]
Then, for every $t\in\mathbb R$,
\[
\mathbb{E}_{\mu,\;P^{M,\textup{ALG}}}[(L-t)_+]
= L_{\max}\,a_t
\ \ge\
L_{\max}\,\underline a_f\!\big(B;\,b_t\big),
\qquad
b_t=\frac{1}{L_{\max}}\mathbb{E}_{\mu,\;Q}[(L-t)_+].
\]
Consequently, by the Rockafellar--Uryasev identity
$\mathrm{CVaR}_\alpha(L)=\min_{t\in\mathbb R}\{\,t+\tfrac{1}{1-\alpha}\,\mathbb{E}[(L-t)_+]\,\}$,
\begin{equation}
\label{eq:cvar-bt}
\mathrm{CVaR}_\alpha(L)\ \ge\
\min_{t\in\mathbb R}\Big\{\,t+\frac{L_{\max}}{1-\alpha}\,
\underline a_f\!\big(B;\,b_t\big)\Big\}
\quad\text{with}\quad
b_t=\frac{1}{L_{\max}}\mathbb{E}_{\mu,\;Q}[(L-t)_+].
\end{equation}
\end{corollary}

\begin{proof}
Apply Theorem~\ref{thm:bt-interactive-Fano-calibrated} with $\phi=\phi_t$ to get
$D_f(\mathrm{Bern}(a_t)\,\|\,\mathrm{Bern}(b_t))\le B$, hence $a_t\ge \underline a_f(B;b_t)$ by definition of the inverse calibration. Multiply by $L_{\max}$ and minimize over $t$ using the CVaR representation.
\end{proof}

\begin{corollary}[KL + mixture explicit form]
\label{cor:cvar-kl-mixture}
Let $f=\mathrm{KL}$ and choose the mixture reference $Q^\star=\mathbb{E}_{M\sim\mu}P^{M,\textup{ALG}}$. Then $B=I_{\mu,\textup{ALG}}(M;X)$. By Pinsker’s inequality,
\(
\underline a_{\mathrm{KL}}(B;b)\ \ge\ \max\{0,\,b-\sqrt{B/2}\}.
\)
Therefore,
\begin{equation}
\label{eq:cvar-kl-mixture}
\mathrm{CVaR}_\alpha(L)\ \ge\
\min_{t\in\mathbb R}\Big\{\,t+\frac{L_{\max}}{1-\alpha}\,
\big[b_t-\sqrt{\,I_{\mu,\textup{ALG}}(M;X)/2\,}\big]_+\Big\},
\quad
b_t=\frac{1}{L_{\max}}\mathbb{E}_{\mu,\;Q^\star}[(L-t)_+]\,,
\end{equation}
where $[z]_+=\max\{z,0\}$.
\end{corollary}

\begin{proof}
With $f=\mathrm{KL}$ and $Q=Q^\star$, Theorem~\ref{thm:bt-interactive-Fano-calibrated} yields
$\mathrm{KL}(\mathrm{Bern}(a_t)\,\|\,\mathrm{Bern}(b_t))\le B=I_{\mu,\textup{ALG}}(M;X)$.
Total variation between Bernoulli$(a_t)$ and Bernoulli$(b_t)$ equals $|a_t-b_t|$, and Pinsker implies
$|a_t-b_t|\le \sqrt{B/2}$, hence $a_t\ge [\,b_t-\sqrt{B/2}\,]_+$. Insert into
\eqref{eq:cvar-bt}.
\end{proof}





\end{proof}
\fi
\iffalse
\cite{ma2024high} points out that the benefit of Theorem \ref{thm:lower-minimax-quantile-relation} is that if for a fixed $r$,
\[
    \inf_{\textup{ALG}} \sup_MP^{M,\textup{ALG}}(L(M,X) >r) > \delta
\]
then,
\[
    \mathfrak{M}(\delta) \geq \mathfrak{M}_{-}(\delta) \geq r
\]
\fi

\iffalse
\begin{proposition}[DV stepping stone for bounded transforms]
Fix $f=\mathrm{KL}$. For any measurable $\phi:[0,\infty)\to[0,1]$, any $\lambda>0$, and any $Q\in\Delta(\mathcal X)$,
\[
\mathbb{E}_{M\sim\mu,\,X\sim P^{M,\textup{ALG}}}\!\big[\phi(L(M,X))\big]
\;\ge\;
-\frac{1}{\lambda}\,\mathbb{E}_{M\sim\mu}\mathrm{KL}\!\big(P^{M,\textup{ALG}}\big\|Q\big)
-\frac{1}{\lambda}\,\log\mathbb{E}_{M\sim\mu,\,X\sim Q}\!\big[e^{-\lambda\phi(L(M,X))}\big].
\]
\end{proposition}

\begin{proof}[Proof (one line via DV)]
Apply the DV variational formula to $P_0(m,x)=\mu(m)P^{M,\textup{ALG}}(x)$ vs. $P_1(m,x)=\mu(m)Q(x)$ with $g(m,x)=-\lambda\,\phi(L(m,x))$:
\[
\mathbb{E}_{P_0}[g]\;\le\;\mathrm{KL}(P_0\|P_1)+\log\mathbb{E}_{P_1}[e^{g}],
\]
rearrange and divide by $-\lambda$.
\end{proof}

\paragraph{Convenient envelope (Bernoulli mgf).}
Let $b:=\mathbb{E}_{\mu,Q}[\phi(L)]\in[0,1]$. Since $z\mapsto e^{-\lambda z}$ is convex on $[0,1]$,
\[
\log\mathbb{E}_{\mu,Q}[e^{-\lambda\phi(L)}]\ \le\ \log\!\big((1-b)+b\,e^{-\lambda}\big).
\]
Hence, for any $\lambda>0$,
\[
\mathbb{E}_{\mu,P^{M,\textup{ALG}}}[\phi(L)]
\ \ge\
b\ -\ \frac{1}{\lambda}\,\mathbb{E}_{M}\mathrm{KL}(P^{M,\textup{ALG}}\|Q)\ -\ \frac{1}{\lambda}\log\!\big((1-b)+b\,e^{-\lambda}\big).
\]
Optimizing in $\lambda$ gives a tight expectation lower bound; choosing $Q^\star=\mathbb{E}_{M}P^{M,\textup{ALG}}$ turns the KL budget into $I_{\mu,\textup{ALG}}(M;X)$.
\fi

% --- Attribution snippets you can paste near the proof sketch ---
% Concise:
% \emph{Proof sketch.} Our argument follows the interactive Fano framework of \cite[Thm.~2, App.~C.1]{chenInteraction} verbatim. 
% The only change is the use of the non-strict success event $\{L\le \Delta\}$, which yields a bound for the strict tail $\mathbb P(L>\Delta)$. 
% We also state a uniform-in-\textup{ALG} form that we later use for minimax quantile lower bounds.

% Fuller:
% \emph{Proof sketch (after \cite{chenInteraction}).} We adapt the interactive Fano method of \cite[Thm.~2]{chenInteraction} to the non-strict success indicator $\{L\le \Delta\}$. 
% The data-processing and Bernoulli monotonicity steps are unchanged; the modification only affects the tail event on the left-hand side, which becomes $\mathbb P(L>\Delta)$. 
% We further package the bound in a uniform-in-\textup{ALG} form (Cor.~\ref{cor:minimax-quantile}) to obtain minimax strict-quantile lower bounds used later.




\iffalse
\begin{theorem}[High probability interactive Fano]

Fix an $f$-divergence $D_f$. Let $\mathrm{ALG}$ be a given algorithm, $\mu\in\Delta(\mathcal{M})$ a prior over models, and $\Delta>0$ a risk level. For any reference distribution $Q\in\Delta(\mathcal{X})$, define
\[
\bar\rho_{\Delta,Q}\;:=\;\mathbb{P}_{M\sim\mu,\;X\sim Q}\!\bigl(L(M,X)\le \Delta\bigr).
\]
Then the following quantile lower bound holds:
\begin{align*}
&\mathbb{P}_{M\sim\mu,\;X\sim \mathbb{P}_{M,\mathrm{ALG}}}
   \!\bigl(L(M,X)>\Delta\bigr)
   \;\ge\; \\[0.5ex]
   &\sup_{Q\in\Delta(\mathcal{X}),\,\epsilon\in[0,1]}
     \Bigl\{\,\epsilon:\;
       \mathbb{E}_{M\sim\mu}\!\bigl[
         D_f(\mathbb{P}_{M,\mathrm{ALG}}\|Q)
       \bigr]<
       d_{f,\epsilon}(\bar\rho_{\Delta,Q}) 
     \Bigr\} = \epsilon^*.
\end{align*}



where 
\[
d_{f,\delta}(p)\;:=\; 
\begin{cases}
D_f\!\bigl(\mathrm{Bern}(1-\delta),\mathrm{Bern}(p)\bigr), & \text{if } p\le 1-\delta,\\[3pt]
0, & \text{otherwise.}
\end{cases}
\]
Assume further that supremum is at $\epsilon^*$.
This further implies: %(only if $\sup = \max$)%

\begin{align*}
\mathfrak{M}(\delta) \geq\mathfrak{M}\_(\delta) \geq \Delta,
 \forall \delta \in [0,\epsilon^*) .
\end{align*}
\end{theorem}

\begin{proof}
The proof is identical to that of Chen et al., Theorem~2, with the change that we replace the indicator $\mathbf{1}\{L(M,X)<\Delta\}$ by $\mathbf{1}\{L(M,X)\le \Delta\}$.
To derive the bound on the minimax quantile we can note:
\textcolor{red}{Write the below better and shorter}
\begin{align*}
    &\sup_{M} \mathbb{P}_{M\sim\mu,\;X\sim \mathbb{P}_{M,\mathrm{ALG}}}
   \!\bigl(L(M,X)>\Delta\bigr)
   \geq
   \\ &\mathbb{P}_{M\sim\mu,\;X\sim \mathbb{P}_{M,\mathrm{ALG}}}
   \!\bigl(L(M,X)>\Delta\bigr)
\end{align*}
Since 
\begin{align*}
&\sup_{M} \mathbb{P}_{M\sim\mu,\;X\sim \mathbb{P}_{M,\mathrm{ALG}}}
   \!\bigl(L(M,X)>\Delta\bigr) =  \\&\sup_{M}\mathbb{P}_{M,\mathrm{ALG}}
   \!\bigl(L(M,X)>\Delta\bigr),
\end{align*}
we get:
\begin{align*}
    &\inf_{\textup{ALG}}\sup_{M}  \mathbb{P}_{M,\mathrm{ALG}}
   \!\bigl(L(M,X)>\Delta\bigr)
   \geq
   \\ &\mathbb{P}_{M\sim\mu,\;X\sim \mathbb{P}_{M,\mathrm{ALG}}}
   \!\bigl(L(M,X)>\Delta\bigr) > \delta. 
\end{align*}
And we can then use theorem 1 to find the lower bound on the minimax quantile.

\end{proof}

\iffalse
\begin{definition}[$\Delta$- Quantile risk]
    The $\Delta$ quantile risk for a prior $\mu$ and any distribution $\mathbb{Q}$ for any $a \in (0,1)$ is given by $\rho_{\Delta,Q,a} = \ \mathbb{P}_{M\sim\mu, X\sim\mathbb{Q}}({L(M,X) \leq a\Delta})$
    \textcolor{blue}{Why $\mu$ is not inside $\rho$ definition?$\rho_{\mu,\delta,...}$ and what is the point of $a$? why not $a$ inside $\Delta$? Can't we scale $\Delta$?}
\end{definition}

\begin{definition}[Threshold divergence]
    The Threshold divergence $d_{f,\beta}$ is defined by  $d_{f,\beta}(p) = D_f(\textup{Bern}(1-\beta),\textup{Bern}(p))$ if $p \leq 1- \beta$ and $d_{f,\beta}(p) = 0$ otherwise. 
\end{definition}
\begin{definition}[Best case quantile parameter]
   Suppose a prior $\mu$ over the models is specified. The best case quantile parameter $\beta(\textup{ALG}, \Delta, \textcolor{blue}{a}?)$ is the maximum parameter such that the average f-divergence is less than the threshold divergence i.e.,
\begin{align*}
\beta(\textup{ALG}, \Delta, a) 
&= \sup_{\substack{\mathbb{Q} \in \Delta(\mathcal{X}) \\ \beta' \in [0,1]}} 
\bigl\{ \beta': \mathbb{E}_{M \sim \mu} \bigl[ D_f(P^{M,\textup{ALG}}, \mathbb{Q}) \bigr] \\
&\quad < d_{f,\beta'}(\rho_{\Delta, \mathbb{Q}, a}) \bigr\}
\end{align*}
\end{definition}
\fi

\begin{theorem}[High probability Interactive Fano version 2]
Let f-divergence $D_f$, an algorithm $\textup{ALG}$ and a prior $\mu \in \mathcal{P}(\mathcal{M})$, a risk level $\Delta \geq 0$ be fixed. Let us define:
$\rho_{\Delta,\mathbb{Q}} = \mathbb{P}_{M\sim\mu, X\sim\mathbb{Q}}(L(M,X) < \Delta)$. 
Suppose there exists $\epsilon \in (0,1]$ such that 
\[
  \inf_{\textup{ALG}}\beta(\textup{ALG}, \Delta, a) \geq \epsilon 
\] for any $a \in (0,1)$, then the lower minimax quantile is bounded below as follows:
\[
\mathfrak{M}\_(\delta) \geq \Delta
\] for all $\delta \in (0,\epsilon).$
\end{theorem}
\begin{proof}
%\textcolor{red}{Fix dependence on algorithm}
The proof is similar to the proof from \cite[Lemma 7]{ma2024high} while we modify the kernel and the Fano's inequality to use the interactive versions from \cite{chen2024assouad}. 
Define $h(x) = \mathbf{1}_{\{x>a\Delta\}} $
\begin{align*}
\inf_{\textup{ALG} \in \mathcal{D}} \sup_{M \in \mathcal{M}} 
P^{M,\textup{ALG}}(L(M,X) > a\Delta) 
&= \inf_{\textup{ALG} \in \mathcal{D}} \sup_{M \in \mathcal{M}} \\
&\quad \mathbb{E}^{M,\textup{ALG}}[\tilde{h}(L(M,X))]
\end{align*}
\begin{align*}
\geq \Delta' \biggl\{ 
\sup_{\substack{\mathbb{Q} \in \Delta(\mathcal{X}) \\ \beta \in [0,1]}} 
\bigl\{ \beta : \mathbb{E}_{M \sim \mu} \bigl[ D_f(P^{M,\textup{ALG}}, \mathbb{Q}) \\
< d_{f,\beta}(\rho'_{1,\mathbb{Q}}) \bigr] \bigr\} 
\biggr\} 
\geq \epsilon > \delta
\end{align*}
\begin{align*}
= \biggl\{ 
\sup_{\substack{\mathbb{Q} \in \Delta(\mathcal{X}) \\ \beta \in [0,1]}} 
\bigl\{ \beta : \mathbb{E}_{M \sim \mu} \bigl[ D_f(P^{M,\textup{ALG}}, \mathbb{Q}) \\
< d_{f,\beta}(\rho'_{1,\mathbb{Q}}) \bigr] \bigr\} 
\biggr\} 
\geq \epsilon > \delta
\end{align*}
where
\[
\rho'_{1,\mathbb{Q}} = \mathbb{P}_{M\sim\mu, X\sim\mathbb{Q}}(\mathbf{1}_{L(M,X) > a} < 1)
\]
which means 
\[
\rho'_{1,\mathbb{Q}} = \mathbb{P}_{M\sim\mu, X\sim\mathbb{Q}}({L(M,X) \leq a\Delta})
\]
And
\[
\rho_{\Delta,\mathbb{Q}} = \mathbb{P}_{M\sim\mu, X\sim\mathbb{Q}}(L(M,X) < \Delta)
\]
We know $\rho'_{1,\mathbb{Q}}\leq\rho_{\Delta,\mathbb{Q}}$ since
\[
\{M,X: L(M,X)\leq a.\Delta\} \subset\{M,X: L(M,X)< \Delta\}
\]

\end{proof}
\fi


\iffalse
\begin{theorem}[Fano inequality for VAR and other risk measures]
\label{cor:info-criterion}

\end{theorem}

\begin{proof}[Proof sketch]

\end{proof}





\begin{theorem}[Fano inequality using strong DPI]
\end{theorem}

\begin{proof}[Proof]


\end{proof}
\fi













\iffalse
\section{How does this work differ from \cite{chen2024assouad}?}



\cite{chen2024assouad} give a bound for the minimax quantile probability as follows

Let any $T \geq 1$ and $\delta \in [0,1)$ be given, and define
\[
\varepsilon_{\delta}(T) := \frac{1}{14}\sqrt{\frac{\delta}{T}}.
\]
Then, for any algorithm $\mathrm{ALG}$, there exists $M^\star \in \mathcal{M}$ such that
\[
\mathbb{P}_{M^\star,\mathrm{ALG}}\left(\mathrm{Risk}_{\mathrm{DM}}(T)\geq \text{p-dec}_{\varepsilon_{\delta}(T),\delta}^{q}(\mathcal{M})\right)\geq \frac{\delta}{2}.
\]

Note however that this is a different tail quantity:
\begin{align*}
\mathbb{P}_{M^\star,\mathrm{ALG}} \Bigl( \mathrm{Risk}_{\mathrm{DM}}(T) 
&\geq \text{p-dec}_{\varepsilon_{\delta}(T),\delta}^{q}(\mathcal{M}) \Bigr) \\
&\geq 
\mathbb{P}_{M^\star,\mathrm{ALG}} \Bigl( \mathrm{Risk}_{\mathrm{DM}}(T) \\
&\quad > \text{p-dec}_{\varepsilon_{\delta}(T),\delta}^{q}(\mathcal{M}) \Bigr)
\end{align*}
This implies that we cannot find a lower bound for our minimax quantile based on Chen et al.'s quantile. But if we define a new minimax quantile with non strict inequality as:

\begin{definition}[Weak Quantile]
For a given level $\delta \in (0,1]$, the (1-$\delta$)th quantile is defined as:
\begin{align*}
\textup{WQuantile}(1-\delta, P^{M,\textup{ALG}}, L) 
= \inf \bigl\{ r \in [0, \infty] :\  
P^{M,\textup{ALG}}( \\
L(M,X) \geq r ) \leq \delta \bigr\}
\end{align*}\end{definition}

Then we get a bound for the minimax quantile based on the Weak Quantile as follows:
\[
\mathfrak{M}_W(\delta) \geq \text{p-dec}_{\varepsilon_{\delta}(T),\delta}^{q}(\mathcal{M})
\]

\textcolor{red}{Simplify this part to show only the gist of the discussion}

Further \cite{chen2024assouad} give a bound on the quantile-DEC as follows:

\begin{proposition}[Proposition D.2]
Suppose that Assumption 4 holds for $\mathcal{M}$.  
Then, for any $\widehat{M}\in\mathcal{M}^{+}$, it holds that
\[
\operatorname{r\text{-}dec}^{\mathrm{c}}_{\varepsilon}\!\bigl(\mathcal{M}\cup\{\widehat{M}\},\widehat{M}\bigr)
   \;\le\;
   2\cdot\operatorname{r\text{-}dec}^{q}_{\varepsilon,\delta}(\mathcal{M},\widehat{M})
   \;+\;
   c_{\delta}\,C_{r}\,\varepsilon,
\]
where
\[
c_{\delta}\;=\;\max\!\Bigl\{\tfrac{\delta}{1-\delta},\,1\Bigr\}.
\]

In particular,
\[
\operatorname{r\text{-}dec}^{q}_{\varepsilon,\,1/2}(\mathcal{M})
   \;\ge\;
   \frac{1}{2}\,
   \Bigl(
      \max_{\widehat{M}\in\mathcal{M}^{+}}
      \operatorname{r\text{-}dec}^{\mathrm{c}}_{\varepsilon}\!\bigl(\mathcal{M}\cup\{\widehat{M}\},\widehat{M}\bigr)
      \;-\;
      C_{r}\,\varepsilon
   \Bigr).
\]
\end{proposition}

This implies a bound for the minimax quantile as follows based on the constrained DEC.
\[
\mathfrak{M}_W(\delta) \geq f(\max_{\widehat{M}\in\mathcal{M}^{+}}
      \operatorname{r\text{-}dec}^{\mathrm{c}}_{\varepsilon}\!\bigl(\mathcal{M}\cup\{\widehat{M}\},\widehat{M}\bigr)
      \;-\;
      C_{r}\,\varepsilon)
\]
We know that 
\[
\mathbb{P}(L\geq \Delta + \epsilon) \leq \mathbb{P}(L>\Delta )
\]
Since,
\[
(L\geq \Delta + \epsilon) \implies (L>\Delta )
\]
This implies if we have a lower bound on every $\Delta$ level then we have a lower bound on $\Delta+ \epsilon$ level.
\[
\delta'\leq\mathbb{P}(L\geq \Delta + \epsilon) \leq \mathbb{P}(L>\Delta )
\]
But we cannot take limit and hence we cannot find a bound for 
$
 \mathbb{P}(L>\Delta )
$ level.
There are  things that can be improved in this line of chen et al's work.
There are two things that can be improved in this line of chen et al's work.
\begin{itemize}
    \item A lower bound on the weak quantile does not give a bound on the quantile and hence there the minimax quantile cannot be bounded in terms of the  minimax weak quantile
    \item the framework does not give level specific bounds. i.e. we must find a lower bound on the quantile-DEC that depends on $\delta$ and this has not been done yet.
\end{itemize}


%\begin{theorem}
%    Relation between minimax quantile and minimax probability:
%    \[
%     \mathbb{P}(L> r) > \delta \implies   \inf_{r'}\{\mathbb{P}(L\geq r') \leq \delta\} \geq r
%    \]
%     \[
%     \inf_\textup{ALG} \sup_M \mathbb{P}(L\geq r) \geq %\delta \implies   \inf_\textup{ALG} \sup_M  \inf_{r'}\%{\mathbb{P}(L\geq r') \leq \delta\} \geq r
 %   \]
%    This implies:
%    \[
 %    \mathfrak{M}(\delta) \geq r
%    \]
%\end{theorem}

The above theorem says that any upper bound on the joint probability can be an upper bound on the minimax quantile probability when the minimax theorem is satisfied.

\fi








\iffalse
The following will demonstrate the simplicity of our framework to derive high probability/quantile lower bounds for the minimax risk. 
In this section, we will define a simple 2-arm bandit problem as an instance of ISDM and apply Theorem 5 to derive a lower bound on the lower minimax quantile. 


A two-armed Gaussian bandit can be viewed as an instance of the ISDM framework.
Let the model class be
\[
\mathcal{M}
=\big\{\,\big(\mathcal{N}(\mu_1,1),\,\mathcal{N}(\mu_2,1)\big)\ :\ (\mu_1,\mu_2)\in\mathbb{R}^2\,\big\}.
\]

Let the history space be $\mathcal{H}^t := (\{1,2\}\times\mathbb{R})^t$ and let
$\mathrm{ALG}$ be a sequence of stochastic decision kernels
$\{\Pi_t\}_{t=1}^T$ with $\Pi_t(\cdot \mid h^{t-1}) \in \Delta(\{1,2\})$ for each
$h^{t-1}\in\mathcal{H}^{t-1}$. Then the interaction evolves as
\[
\pi_t \sim \Pi_t(\cdot \mid H^{t-1}),\qquad
o_t \mid (\pi_t,H^{t-1}) \sim M_\mu(\pi_t)=\mathcal{N}(\mu_{\pi_t},1),
\]
with $(o_t)$ conditionally independent of $H^{t-1}$ given $\pi_t$.
The (random) trajectory is
\[
H^T := (\pi_1,o_1,\ldots,\pi_T,o_T)\in(\{1,2\}\times\mathbb{R})^T.
\]
The data is then given by \(X=H^T\),
with distribution \(\mathbb{P}^{M,\mathrm{ALG}}\). Further, let the loss be the regret given by \(L(M,X)  = \sum^T_{t=1} \mu^* - \mu_{\pi_t}\) where $\mu^* = \max(\mu_1,\mu_2)$.
\begin{theorem}
    For the two-armed Gaussian bandit problem above with $T\geq1$, for $\delta \in (0,1/2)$, we have $\mathfrak{M}_{-}(\delta) \geq \sqrt{\frac{T\log(1/4\delta(1-\delta))}{2}}.$
\end{theorem}


\begin{proof}
Fix any (possibly randomized, adaptive) algorithm. For a model \(M=(\mu_1,\mu_2)\) and trajectory \(x\) with actions \(A_{1:T}\), let \(\mathfrak{R}_T(M,x)=\sum_{t=1}^T(\mu^\star(M)-\mu_{A_t})\), where \(\mu^\star(M)=\max\{\mu_1,\mu_2\}\). Pick \(g>0\) and consider \(M_1:(\mu_1,\mu_2)=(+g/2,-g/2)\) and \(M_2:(\mu_1,\mu_2)=(-g/2,+g/2)\). If \(N_a(x)\) is the number of pulls of arm \(a\) in \(x\), then \(\mathfrak{R}_T(M_1,x)=g\,N_2(x)\) and \(\mathfrak{R}_T(M_2,x)=g\,N_1(x)\); thus \(\mathfrak{R}_T(M_1,x)+\mathfrak{R}_T(M_2,x)=g(N_1(x)+N_2(x))=gT\), i.e., the separation condition holds with \(\Delta=gT/2\). Let \(\mathbb{P}_i\) be the trajectory law under \(M_i\) and set \(\Lambda_\delta=\log(1/(4\delta(1-\delta)))\). By the divergence decomposition lemma \cite[Lemma~15.1]{lattimore2020bandit}, \(\mathrm{KL}(\mathbb{P}_1\|\mathbb{P}_2)=\sum_{a=1}^2 \mathbb{E}_{M_1}[N_a(T)]\,\mathrm{KL}(\mathcal N(\mu_a^{(1)},1)\Vert \mathcal N(\mu_a^{(2)},1))=(g^2/2)\,(\mathbb{E}_{M_1}[N_1(T)]+\mathbb{E}_{M_1}[N_2(T)])=(g^2/2)\,T\). By Theorem~\ref{theorem: High probability version for Le Cam with interaction using Divergence}, if \(\mathrm{KL}(\mathbb{P}_1\|\mathbb{P}_2)<\Lambda_\delta\) and the separation holds, then \(\mathfrak{M}_{-}(\delta)\ge \Delta\). Hence whenever \(g^2T/2<\Lambda_\delta\), we have \(\mathfrak{M}_{-}(\delta)\ge gT/2\). Choose any \(\eta\in(0,1)\) and take \(g=(1-\eta)\sqrt{2\Lambda_\delta/T}\), which yields \(\mathfrak{M}_{-}(\delta)\ge (1-\eta)\sqrt{(T/2)\Lambda_\delta}=(1-\eta)\sqrt{(T/2)\log(1/(4\delta(1-\delta)))}\); letting \(\eta\downarrow 0\) proves the claim.
\end{proof}



Chapter~17 of \cite{lattimore2020bandit} gives high-probability lower bounds for the same loss using the Bretagnolle–Huber inequality with a KL–decomposition, yielding a weak-tail bound of the form $\mathbb{P}(L\ge r)\ge \delta$; in the two-arm Gaussian case the scaling matches our $\sqrt{T\log(1/\delta)}$ up to constants. By contrast, our result is a $\delta$-explicit minimax quantile (strong-tail) lower bound derived via an interactive Le Cam condition at the KL threshold $\log\!\tfrac{1}{4\delta(1-\delta)}$. Instantiating this condition for the two-armed Gaussian bandit yields the bound as an immediate corollary.
 While the theorem above is bandit-specific, Theorem~\ref{theorem: High probability version for Le Cam with interaction using Divergence} is within the ISDM framework; thus this bandit result serves as a template for analogous high-probability lower bounds in other interactive decision problems.
\fi




\iffalse
\begin{proof}
Fix any (possibly randomized, adaptive) algorithm. For a model $M=(\mu_1,\mu_2)$ and trajectory $x$ with actions $A_{1:T}$, define
\[
L(M,x)\;=\;\mathfrak{R}_T(M,x)\;\coloneqq\;\sum_{t=1}^T\big(\mu^\star(M)-\mu_{A_t}\big),.
\]
where $\mu^\star(M)\coloneqq\max\{\mu_1,\mu_2\}$.
Pick $g>0$ and set
\[
M_1:\ (\mu_1,\mu_2)=\Bigl(+\tfrac{g}{2},-\tfrac{g}{2}\Bigr),\qquad
M_2:\ (\mu_1,\mu_2)=\Bigl(-\tfrac{g}{2},+\tfrac{g}{2}\Bigr).
\]
Let $N_a(x)$ be the number of pulls of arm $a$ in $x$. Then
\[
\mathfrak{R}_T(M_1,x)=g\,N_2(x),\qquad
\mathfrak{R}_T(M_2,x)=g\,N_1(x),
\]
and hence, for all $x$,
\[
\mathfrak{R}_T(M_1,x)+\mathfrak{R}_T(M_2,x)=g\big(N_1(x)+N_2(x)\big)=gT,
\]
so the separation condition holds with $\Delta=\tfrac{gT}{2}$.

 Let $\mathbb{P}_i$ be the trajectory law under $M_i$ and write
$\Lambda_\delta\coloneqq\log\!\tfrac{1}{4\delta(1-\delta)}$.
By the divergence decomposition (\cite[Lemma~15.1]{lattimore2020bandit}),
\begin{align*}
\mathrm{KL}(\mathbb{P}_1\|\mathbb{P}_2)
&=\sum_{a=1}^2 \mathbb{E}_{M_1}[N_a(T)]\,
\mathrm{KL}\!\Big(\mathcal N(\mu_a^{(1)},1)\,\Big\|\,\mathcal N(\mu_a^{(2)},1)\Big)\\
&=\frac{g^2}{2}\big(\mathbb{E}_{M_1}[N_1(T)]+\mathbb{E}_{M_1}[N_2(T)]\big)
=\frac{g^2}{2}\,T,
\end{align*}
since $N_1(T)+N_2(T)=T$.

By Theorem~\ref{theorem: High probability version for Le Cam with interaction using Divergence}, if
$\mathrm{KL}(\mathbb{P}_1\|\mathbb{P}_2)<\Lambda_\delta$ and the separation holds, then
$\mathfrak{M}_{-}(\delta)\ge\Delta$. The KL condition is met whenever
$\tfrac{g^2}{2}T<\Lambda_\delta$, hence $\mathfrak{M}_{-}(\delta)\ge \tfrac{gT}{2}$
for any such $g$. Fix $\eta\in(0,1)$ and choose
\[
g=(1-\eta)\sqrt{\frac{2\Lambda_\delta}{T}}.
\]
Then
\[
\mathfrak{M}_{-}(\delta)\;\ge\;(1-\eta)\sqrt{\frac{T}{2}\,\Lambda_\delta}
\;=\;(1-\eta)\sqrt{\frac{T}{2}\,\log\!\frac{1}{4\delta(1-\delta)}}.
\]
Letting $\eta\downarrow 0$ yields the claim.
\textcolor{red}{Give discussion of the above result}
\end{proof}
\fi






        
\iffalse
\textcolor{red}{Show how the ISDM setting includes the MAB}
Standard lower bound proofs involve:
i) build two nearly identical and indistuinguishable worlds. 
ii) Since the worlds are indistinuishable we need to pull even the wrong arm many times
iii) this will ensure high regret
\subsection{Stochastic Multi-arm Bandit problem}
\textcolor{red}{Check this example. Fix the proof}
\begin{theorem}[High–probability minimax quantile lower bound, two–armed Gaussian bandit]%
\label{thm:minimax-quantile}
Let $n \ge 1$ be the horizon, and fix a confidence level 
\(
  \delta \in (0,\tfrac14).
\)
For every bandit algorithm\/\footnote{%
An algorithm may be adaptive and randomised; at round $t$ it chooses an arm 
$A_t\in\{1,2\}$ as an arbitrary measurable function of past actions, past rewards and an internal random seed.}
$\pi$ there exists a Gaussian two–armed environment%
\footnote{%
Arm~$i$ produces i.i.d.\ rewards $X_t^{(i)}\sim\mathcal N(\mu_i,1)$.
}
such that the \((1-\delta)\)-quantile of the average pseudo-regret
\[
  L \;=\; \frac1n\sum_{t=1}^{n}\bigl(\mu^{\star}-\mu_{A_t}\bigr),
  \quad
  \mu^{\star}\;=\;\max\{\mu_1,\mu_2\},
\]
satisfies
\[
  Q_{1-\delta}\!\bigl(L\bigr)
  \;\;\ge\;
  \sqrt{\frac{\log\bigl(1/4\delta(1-\delta)\bigr)}{2n}}
  \;\;\ge\;
  \frac{1}{\sqrt{8n}}\sqrt{\log\!\bigl(1/\delta\bigr)} .
\]
Consequently, the minimax $(1-\delta)$-quantile
\(
  M(\delta)=\inf_{\pi}\sup_{(\mu_1,\mu_2)}Q_{1-\delta}\bigl(L\bigr)
\)
obeys
\(
  M(\delta)\;=\;\Omega\!\bigl(\sqrt{\log(1/\delta)/n}\bigr).
\)
\end{theorem}

\begin{proof}
The argument is an instantiation of Corollary~6 from
\cite{ma2024high}.
We proceed in five steps.

\paragraph{\textbf{~Hard\, pair of environments.}}
Fix a \emph{gap} $\Delta>0$ (to be chosen later) and define
\[
  \theta_1:\; (\mu_1,\mu_2)=(\Delta,0),
  \qquad
  \theta_2:\; (\mu_1,\mu_2)=(0,\Delta).
\]
Let $\mathbb{P}_1$ (resp.\ $\mathbb{P}_2$) denote the law of the full
trajectory $(A_{1:n},X_{1:n})$ generated by $\pi$ in~$\theta_1$ (resp.\ $\theta_2$).

\paragraph{\textbf{~KL divergence.}}
Because the decision kernel of~$\pi$ is identical under the two environments,
the chain rule gives
\[
  \operatorname{KL}(\mathbb{P}_1,\mathbb{P}_2)
    \;=\;
    \sum_{t=1}^{n} 
      \frac{\Delta^{2}}{2}
    \;=\;
    \frac{\Delta^{2}n}{2}.
\tag{$\star$}%
\]

\paragraph{\textbf{~Choose $\Delta$ with $\varepsilon$--slack.}}
Fix an arbitrary $\varepsilon\in(0,1)$ and set
\[
  \Delta
  \;:=\;
  (1-\varepsilon)\,
  \sqrt{
    \frac{2\log\!\bigl(1/4\delta(1-\delta)\bigr)}{n}
  } .
\]
With this choice, $(\star)$ yields
\(
  \operatorname{KL}(\mathbb{P}_1,\mathbb{P}_2)
  < \log\!\bigl(1/4\delta(1-\delta)\bigr),
\)
so the KL pre-condition of Corollary~6 is \emph{strictly} satisfied.

\paragraph{\textbf{~Metric and loss mapping.}}
Define a pseudo-metric
\(
  d(\theta_1,\theta_2)=|\mu_1-\mu_2|=\Delta ,
\)
and take the non-decreasing map $g(x)=x$.
Then
\(
  g\!\bigl(d/2\bigr)=\Delta/2 .
\)

\paragraph{\textbf{~Invoke Theorem 5}}
Theorem 5 now gives, for \emph{every} algorithm~$\pi$,
\[
  \max_{\theta\in\{\theta_1,\theta_2\}}
     \mathbb{P}_\theta\!\bigl\{L>\tfrac{\Delta}{2}\bigr\}
  \;\ge\;
  \delta .
\]
That inequality is equivalent to
\(
  Q_{1-\delta}\bigl(L\bigr)\ge\Delta/2 .
\)
Letting $\varepsilon\downarrow0$ delivers
\[
  Q_{1-\delta}\bigl(L\bigr)
    \;\ge\;
    \sqrt{\frac{\log\bigl(1/4\delta(1-\delta)\bigr)}{2n}}
    \;\ge\;
    \tfrac{1}{\sqrt{8n}}\sqrt{\log(1/\delta)} .
\]

\paragraph{\textbf{~From $Q_{1-\delta}$ to $M(\delta)$}}
Because the bound holds
\emph{for every $\pi$}, taking $\inf_{\pi}$ over algorithms yields the same lower bound on the minimax quantile $M^{-}(\delta)$.
The ``sandwich lemma'' (Thm.~4 of Ma~et~al.) asserts
\(
  M^{-}(\delta)
  \le
  M(\delta)
  \le
  M^{-}(\delta-\xi)
\)
for any $\xi\in(0,\delta)$; letting $\xi\downarrow0$ completes the proof.
\end{proof}
\textcolor{red}{can we use the high probability interactive Fano derived?}
\fi






\iffalse
Another interesting avenue for further research is to find conditions where the minimax theorem holds and then to find upper bounds for the minimax quantile interms of the upper bound on the joint probability.
\fi







% For peer review papers, you can put extra information on the cover
% page as needed:
% \ifCLASSOPTIONpeerreview
% \begin{center} \bfseries EDICS Category: 3-BBND \end{center}
% \fi
%
% For peerreview papers, this IEEEtran command inserts a page break and
% creates the second title. It will be ignored for other modes.


\end{document}
